% paper16/banded_toeplitz_conditioning.tex
%
\documentclass[11pt]{amsart}

\usepackage[T1]{fontenc}
\usepackage[utf8]{inputenc}
\usepackage{lmodern}
\usepackage{microtype}
\usepackage{amsmath,amssymb,amsthm,mathtools}
\usepackage{mathrsfs}
\usepackage{enumitem}
\usepackage[colorlinks=true,linkcolor=blue,citecolor=blue,urlcolor=blue]{hyperref}

\numberwithin{equation}{section}

\newtheorem{theorem}{Theorem}[section]
\newtheorem{proposition}[theorem]{Proposition}
\newtheorem{corollary}[theorem]{Corollary}
\newtheorem{lemma}[theorem]{Lemma}
\newtheorem{question}[theorem]{Question}
\newtheorem{problem}[theorem]{Problem}
\newtheorem{conjecture}[theorem]{Conjecture}
\newtheorem{hypothesis}[theorem]{Hypothesis}
\theoremstyle{definition}
\newtheorem{definition}[theorem]{Definition}
\theoremstyle{remark}
\newtheorem{remark}[theorem]{Remark}

\newcommand{\C}{\mathbb C}
\newcommand{\R}{\mathbb R}
\newcommand{\Z}{\mathbb Z}
\newcommand{\F}{\mathbb F}
\newcommand{\T}{\mathbb T}
\newcommand{\D}{\mathbb D}
\newcommand{\Prob}{\mathbb P}
\newcommand{\Exp}{\mathbb E}
\newcommand{\Var}{\operatorname{Var}}
\newcommand{\Cov}{\operatorname{Cov}}
\newcommand{\Span}{\operatorname{span}}
\newcommand{\dist}{\operatorname{dist}}
\newcommand{\Id}{I}
\newcommand{\wind}{\operatorname{wind}}
\newcommand{\smin}{\sigma_{\min}}
\newcommand{\kdn}{\kappa}
\newcommand{\norm}[1]{\left\lVert #1\right\rVert}
\newcommand{\abs}[1]{\left\lvert #1\right\rvert}
\newcommand{\eps}{\varepsilon}

\title[A condition-number bound for balanced banded Toeplitz matrices]{A least singular value bound for random balanced banded Toeplitz matrices}
\author{Yitzchak Shmalo}
\address{Einstein Institute of Mathematics, The Hebrew University of Jerusalem, Jerusalem, Israel}
\email{yitzchak.shmalo@gmail.com}
\date{June 2026}
\subjclass[2020]{Primary 15B05, 60B20; Secondary 15A18, 15A12, 42A05, 30C15}
\keywords{banded Toeplitz matrices, least singular value, condition number, random polynomials, small-ball probability, negative second moment, Konyagin--Schlag}

\begin{document}

\begin{abstract}
Let $P_{r,r}(z)=\sum_{j=-r}^{r}\xi_jz^{j}$ be a random symbol with
balanced band $r=s$, bandwidth $m=2r$, and coefficients $\xi_j$ that are
independent, mean zero, variance one and have a bounded density; the Gaussian
case is included. Write $G_m(z)=z^rP_{r,r}(z)=\sum_{j=0}^m\xi_jz^j$ for the
associated Kac polynomial and $T_n=T_n(P_{r,r})$ for the $n\times n$ finite
sections. The companion paper \cite{Parity} proves that the stability
probability of $(T_n)$ is at most $\tfrac12+o(1)$ for \emph{every} band,
through a parity obstruction, and proves exponential instability only for
strongly unbalanced bands $\abs{r-s}>2m^{9/10}$; in the balanced regime it
gives no quantitative bound on the smallest singular value on the stable
subevent. This note isolates the missing conditioning mechanism. Importing the
least singular value machinery of the phase retrieval paper \cite{PhaseRetrieval}
--- the negative second moment small-ball identity and a good-family second
moment philosophy --- we formulate the boundary small-ball estimates for $G_m$
on the unit circle that close the Gaussian balanced case up to the zero-winding
event. The expected single-arc rate on arcs of length $c/m$ is
$\Prob\{\min_I\abs{G_m}\le t\}\le C(1+c)\,t/\sqrt m$, governed by
one-dimensional curve crossing rather than two-dimensional point
anti-concentration; we prove this sharp rate away from $z=\pm1$ by a Kac--Rice
crossing count for \(|G_m(e^{i\theta})|^2\), and prove the endpoint contribution
by a separate endpoint Kac--Rice estimate. We also compute the exact
two-endpoint Gaussian formula and the sharp second-order endpoint expansion
through the \(t^2/N\) term. Together these give the sharp uniform
Gaussian boundary small-ball upper law
\(\Prob\{\min_{\T}\abs{G_m}\le t\}\le C\sqrt m\,t\). In the fixed bulk we
compute the leading Kac--Rice crossing constant and the resulting first-moment
probability upper bound at \(t=\tau/\sqrt N\). At the continuum scale
\(t=\tau/\sqrt N\) the two parity endpoints contribute only
\(2\sqrt{2/\pi}\,\tau/N+O_\tau(N^{-2})\), and the full natural
\(O(1/m)\) endpoint neighborhoods have probability \(O_{C_0}(\tau/N)\), so the
\(O(1)\) candidate law is a non-real-arc phenomenon. We also isolate the exact
remaining lower-tail input: the fixed-bulk upper bound and endpoint-neighborhood
separation give the sharp whole-circle continuum upper constant
\(\sqrt{\pi/3}\,\tau\), while a second-factorial-moment bound for the bulk
down-crossing count implies, by Paley--Zygmund, a fixed-bulk lower tail of order
\((2\pi-4\delta)\tau/\sqrt{12\pi}\) up to the expected quadratic correction.
Together with the sharp upper bounds, the same second-moment input gives the
full first-order fixed-bulk law as \(\tau\downarrow0\), and a version uniform as
\(\delta\downarrow0\) gives the corresponding first-order whole-circle lower
law with leading constant \(\sqrt{\pi/3}\).  The sharp global upper tail also
gives a useful endpoint-integrability consequence: the inverse scaled minimum
\((\sqrt N\min_{\T}|G_m|)^{-1}\) is uniformly weak \(L^1\), has sharp
finite-window weak constant at most \(\sqrt{\pi/3}\), and has
\(O(\alpha^{-1})\) inverse upper-tail quantiles with sharp limsup constant
at most \(\sqrt{\pi/3}\,\alpha^{-1}\).  The same weak endpoint and quantile
bounds hold simultaneously for finite batches, with the expected extra factor
equal to the batch size; the corresponding finite-batch package also controls
bad-clearance counts, clearance order-statistic quantiles, total bad-mass,
powered and jump/Stieltjes-weighted deficit summability, subcritical inverse-mass,
Stieltjes inverse-gauge budgets, confidence, menu, and aggregate-menu bounds,
normalized deficit-budget schedules, and sparse bad-count Borel--Cantelli schedules,
subcritical inverse-power moments, critical truncated
expectation, and logarithmic moments with the same factor.  It has only
logarithmic critical truncated expectation with the same leading coefficient,
and has uniform logarithmic inverse moments of every positive order, with
limsup constant at most \(\sqrt{\pi/3}\Gamma(p+1)\) for the \(p\)-th
logarithmic moment.
Therefore the scaled minimum has uniformly bounded negative moments
of every order below one, with the sharp small-window upper coefficient
inherited from \(\sqrt{\pi/3}\).
For bounded-density laws we
prove the matching fixed-point and deterministic-grid anti-concentration, its
independent-random grid conditioning consequence, plus local and uniform
boundary bounds conditional on a deterministic speed cutoff;
subexponential tails, and more generally uniformly bounded moments of all
orders via Fuk--Nagaev, verify the full polynomial-scale boundary-law
hypothesis. At the continuum threshold, subquadratic deterministic grids, and
subquadratic independent random grids, see only the pointwise small-ball scale,
while a natural polynomial derivative-speed tail makes \(O(N^2\sqrt{\log N})\)
deterministic grids complete for continuum excursions up to arbitrary
polynomial miss probability; independent uniform random grids recover the same
completeness with the extra coupon-collector logarithm
\(O(N^2(\log N)^{3/2})\), while stratified one-point-per-cell random grids and
random shifts of a uniform lattice keep the deterministic
\(O(N^2\sqrt{\log N})\) scale without the coupon overhead. For the actual conditioning theorem, however, no tail input beyond
the variance normalization is needed under the same circle-clearance input: a
Chebyshev--Bernstein polynomial derivative cutoff supplies the single
boundary-minimum event used by the finite-section proof, made explicit below as
tail-free polynomial boundary separation with exceptional probability
\(O_{A,K}(m^{-2(A-1)/3})\) at threshold \(m^{-A}\).  The same variance-only
input has a finite-confidence form:
\(\Prob\{\sup_{\T}|G_m'|>R\}\le Cm^4/R^2\), giving \(\eta\)-level grid
detection at mesh \(O(\tau\sqrt\eta/(m^2\sqrt N))\) and bounded-density
boundary loss \(\eta+C_{B,K}mt/\sqrt\eta\). The cutoff-free
continuum crossing step remains open only for the full boundary small-ball law,
not for the conditioning conclusion under clearance. The finite-section step
itself needs no extra probabilistic coupling: once a Konyagin--Schlag
circle-clearance event holds, Proposition~\ref{prop:fs-uniform} gives a
deterministic threshold uniform over all zero-winding symbols with that
clearance floor. The quantitative loss in the stable-event conclusion is just
the clearance failure probability, the explicit boundary-separation error, and
a tunable coefficient-size error; choosing the boundary floor and coefficient
envelope together gives a prescribed \(q_m+O_D(m^{-D})\) stable-event loss after
increasing only the deterministic finite-section threshold; equivalently, after
conditioning on the zero-winding event, the loss is divided by
\(\Prob\{\nu=r\}\). Consequently, the Gaussian balanced-conditioning theorem is
unconditional on the zero-winding event by the cited Konyagin--Schlag clearance
theorem, while every
bounded-density law satisfying the same circle-clearance estimate has the same
conditioning conclusion. We are careful to separate what is proved --- the
finite-section reduction on the zero-winding event, the sharp Gaussian boundary
separation, fixed-grid and speed-conditional bounded-density anti-concentration,
and the clearance-conditional bounded-density deduction --- from the continuum
non-Gaussian universality estimates that remain open. Finally, the parity
observable \(\nu\equiv\mathbf1\{G(1)G(-1)<0\}\bmod2\) of \cite{Parity} is a
boundary-evaluation effect, and lattice coefficient laws require separate
treatment: for Rademacher coefficients, non-real roots of unity of order \(q\)
obey \(O(q/N)\) fixed-point anti-concentration, denominators \(q\le Q\) obey a
\(Q^3/N\) union bound, every fixed rational angle has matching \(N^{-1}\)
small-ball mass at a constant lattice radius, and every fixed primitive
\(q\)-th root has an exact atom local limit
\(\gamma_qN^{-\varphi(q)/2}\), with an explicit kernel-covolume constant
\(\gamma_q\), on the admissible parity residue classes, and has no atom on the
inadmissible classes. Hence any fixed finite set of cyclotomic orders has an
exact-zero asymptotic on its first nonempty totient layer. Consequently a
fixed bounded-denominator non-real grid has a peeled exact-zero layer formula
after removing all orders below any fixed cyclotomic degree threshold, and the
unpeeled fixed-\(Q\) union has the full expansion
\(c_Q(N)/N+e_Q(N)/N^2+o_Q(N^{-2})\).  Here \(c_Q(N)\) receives contributions
\(2\sqrt3/\pi\) from each of the primitive cubic and sixth pairs on \(3\mid N\),
and \(4/\pi\) from the primitive quartic pair on \(4\mid N\), while \(e_Q(N)\)
combines the second-order corrections to those degree-two atoms, their explicit
pairwise overlaps, and the first nonquadratic layer.  After removing the
cyclotomic degree-two roots, that nonquadratic layer is also explicit: orders
\(5,8,10,12\) contribute at scale \(N^{-2}\), with constants
\(20\sqrt5/\pi^2,64/\pi^2,20\sqrt5/\pi^2,48/\pi^2\) on their admissible
divisibility classes.  If the two endpoints are included, the full
bounded-denominator root-of-unity grid has, through order \(N^{-1}\), the exact
central-binomial endpoint-union term plus the same non-real coefficient
\(c_Q(N)/N\).
\end{abstract}

\maketitle

\section{Introduction}\label{sec:intro}

A banded Toeplitz matrix is constant along its diagonals and supported on finitely many of them. Given nonnegative integers $r,s$ and numbers $x_{-r},\dots,x_s$, the $n\times n$ matrix $T_n$ has entries $(T_n)_{jk}=x_{k-j}$ for $-r\le k-j\le s$ and $0$ otherwise, and the asymptotic theory of $(T_n(P))_{n\ge1}$ is governed by the symbol $P(z)=\sum_{j=-r}^{s}x_jz^j$; the reference is B\"ottcher--Grudsky \cite{BG}. The sequence $(T_n(P))$ is \emph{stable} --- eventually invertible with uniformly bounded inverses --- if and only if $P$ has no zeros on the unit circle $\T$ and zero winding number \cite[Ch.~3]{BG}. Writing
\begin{equation}\label{eq:symbol}
  G_m(z):=z^rP(z)=\sum_{j=0}^m\xi_{j-r}z^j,\qquad m:=r+s,
\end{equation}
and $\nu:=\#\{z\in\D:G_m(z)=0\}$ for the number of zeros in the open unit disk, stability is the requirement that $P\ne0$ on $\T$ and that exactly $r$ of the $m$ zeros of $G_m$ lie inside $\T$.

Randomizing the coefficients turns $G_m$ into a Kac polynomial, and the stability question becomes a delicate question about the exact disk-count $\nu$. The companion paper \cite{Parity} settles the qualitative side of this question and locates a gap in the literature. Its parity obstruction shows that for any independent, mean-zero, unit-variance coefficients with a finite third moment, $\nu\equiv\mathbf1\{G_m(1)G_m(-1)<0\}\pmod2$ on the event $G_m(1)G_m(-1)\ne0$, and the right-hand side is an asymptotically fair coin; hence
\begin{equation}\label{eq:half}
  \Prob\bigl((T_n(P_{r,s}))_n\text{ stable}\bigr)\le\tfrac12+o(1)
\end{equation}
for \emph{every} band. On the unstable side, \cite{Parity} proves that strongly unbalanced bands $\abs{r-s}>2m^{9/10}$ are exponentially ill-conditioned: $\norm{T_n^{-1}}\ge Ce^{\alpha_mn}$ and $\kdn(T_n)$ grows exponentially in $n$, with probability $1-O((\log m)^{-1/2})$. The mechanism is operator-theoretic: a nonzero winding number forces an exponentially blowing-up inverse \cite[Thm.~4.1]{BG}, and \eqref{eq:half} together with the zero-count window confines $\nu$ near $m/2$, so that a band far from balanced cannot have winding zero.

What \cite{Parity} does not provide, and what it explicitly leaves open, is any quantitative control of the smallest singular value $\smin(T_n)$ in the \emph{balanced} regime $r=s$ on the stable subevent. There the parity obstruction caps the stability probability at $\tfrac12$ but says nothing about the size of $\smin(T_n)$ once the zero-winding condition $\nu=r$ actually holds. The operator-theoretic exponential lower bound of the unbalanced case is simply unavailable on that zero-winding event. This paper isolates the zero-winding replacement: for every bounded-density variance-one coefficient law satisfying the same circle-clearance estimate, balanced finite sections are polynomially conditioned on the stable subevent. In the Gaussian case the needed clearance is supplied by Konyagin--Schlag, so the balanced-conditioning conclusion is unconditional on the zero-winding event.

\subsection{The mechanism: from the circle to the finite section}\label{sub:mechanism}

The link between $\smin(T_n)$ and the symbol is classical and local to the unit circle. If $P$ does not vanish on $\T$, then the finite sections inherit a lower bound governed by how close $P$ comes to zero on $\T$ together with the rate at which the nearest zeros of $G_m$ clear the circle. Concretely, for banded symbols there is a constant $c_0=c_0(m)$, polynomial in $m$, such that
\begin{equation}\label{eq:fs-reduction}
  \smin\bigl(T_n(P)\bigr)\;\ge\;c_0(m)^{-1}\,\min_{z\in\T}\abs{P(z)}\,\bigl(\,\mathrm{dist}(\text{zeros of }G_m,\T)\,\bigr)^{O(1)}
\end{equation}
once $n$ exceeds a threshold depending on $m$; we make a clean form of this
precise in Section~\ref{sec:main} (Proposition~\ref{prop:fs}). Two independent
random quantities therefore control $\smin(T_n)$:
\begin{enumerate}[label=(\roman*)]
\item the \emph{boundary minimum} $\min_{z\in\T}\abs{G_m(z)}$, which one needs to keep away from $0$; and
\item the \emph{circle clearance} $\min\{\bigl|\,\abs{\zeta}-1\,\bigr|:G_m(\zeta)=0\}$, the distance of the zero cloud to $\T$.
\end{enumerate}
The clearance (ii) is exactly the Konyagin--Schlag estimate already used in \cite{Parity,KMS}: with high probability no zero of $G_m$ lies within $m^{-2}(\log m)^{-3}$ of $\T$. The boundary minimum (i) is the genuinely new ingredient, and it is here that the phase-retrieval machinery enters.

\subsection{Importing the least singular value machinery}\label{sub:import}

The phase retrieval paper \cite{PhaseRetrieval} develops, for a tall Gaussian matrix $A\in\F^{N\times m}$, two tools that we transplant from the linear-algebra setting to the boundary-evaluation setting on $\T$.

The first is the \emph{negative second moment identity}. For an invertible $m\times m$ matrix $G$ with rows $R_1,\dots,R_m$ and $d_i=\dist(R_i,\Span\{R_j:j\ne i\})$,
\begin{equation}\label{eq:nsm}
  \sum_{k=1}^m\sigma_k(G)^{-2}=\sum_{i=1}^m d_i^{-2},
\end{equation}
where $\sigma_k$ denotes the $k$-th singular value \cite[Lem.~2.1]{PhaseRetrieval}. If $\smin(G)\le t$ the left side is $\ge t^{-2}$, so some distance $d_i\le\sqrt m\,t$; since each $d_i$ is, conditionally on the other rows, a one-dimensional Gaussian projection, a union bound gives the one-row small-ball bound
\begin{equation}\label{eq:onerow}
  \Prob\{\smin(G)\le t\}\le C_\F\,m^{1+d_\F/2}\,t^{d_\F},\qquad d_\F=\dim_\R\F.
\end{equation}
The point is that \eqref{eq:nsm} converts a small \emph{least singular value} (a global, geometric event) into a small \emph{distance of a single row to a hyperplane} (a one-dimensional small-ball event), which is elementary to bound; the small-ball estimate for such a one-dimensional projection is the linear-image small-ball of Rudelson--Vershynin \cite{RudelsonVershynin2015}. The same one-dimensional reduction reappears below as the curve-crossing mechanism on an arc, where the role of the projected coordinate is played by the planar curve $\{G_m(z):z\in I\}$.

The second tool addresses the obstruction that makes the problem hard. To bound a minimum over an exponentially large family one cannot afford a union bound: in \cite{PhaseRetrieval} there are $\binom{N-1}{m-1}=\exp(\Theta(m))$ hyperplanes, and in our setting the analogue is a family of $\Theta(m)$ arcs together with the continuum of points within each. The resolution in \cite{PhaseRetrieval} is a \emph{good-family second moment argument}: one fixes an apex (one Gaussian row) and counts, via $Z_t=\sum_S\mathbf1\{|\langle x,u_S\rangle|\le t\}$, how many hyperplanes the apex nearly meets; the first moment $\Exp Z_t$ is large, and the pairwise overlap combinatorics (Lemmas on the two-slab estimate and the overlap distribution there) make $\Var Z_t/(\Exp Z_t)^2\to0$, so that $Z_t>0$ with high probability by Chebyshev. The crucial structural fact that powers the second moment is \emph{decorrelation}: two row sets sharing few rows give nearly independent hyperplane normals. In the Toeplitz problem the same role is played by the Dirichlet kernel
\begin{equation}\label{eq:dirichlet}
  \Cov\bigl(G_m(z),G_m(w)\bigr)=\sum_{j=0}^m z^j\overline{w}^{\,j}=\frac{1-(z\overline w)^{m+1}}{1-z\overline w},
\end{equation}
which for $z,w\in\T$ at angular separation $\gg1/m$ is small: the evaluations of $G_m$ at points of $\T$ that are more than a Dirichlet width apart are weakly correlated, so $\T$ splits into $\Theta(m)$ arcs that decorrelate, exactly the analogue of the weakly-overlapping row sets. For the upper bound proved below, the continuum of points inside one arc is handled by a curve-crossing argument; the good-family second moment is what the matching \emph{lower} bound, the exponential law of \cite{YakirZeitouni,CookNguyen}, and the sharp whole-circle picture require.

\subsection{Results}\label{sub:results}

The boundary small-ball estimates are the analytic input to the balanced-regime conditioning bound. We first give the proved Gaussian single-arc estimate, then the proved sharp uniform whole-circle estimate. The single-arc estimate is governed not by the two-dimensional point anti-concentration of $G_m(z)$ but by the geometry of the planar curve $\{G_m(z):z\in I\}$: as $z$ traverses an arc of length $c/m$ this curve has length of order $\sqrt m$, so the event that it comes within $t$ of the origin has the \emph{one-dimensional} curve-crossing probability \(t/\sqrt m\), rather than the pointwise rate \(t^2/m\).

\begin{theorem}[Boundary small-ball, single arc]\label{thm:singlearc}
Let $G_m(z)=\sum_{j=0}^m\xi_jz^j$ with $\xi_j$ independent standard real Gaussians. For every fixed $\gamma>0$ there is a constant $C_\gamma>0$ such that for every $c\ge1$, every arc $I\subset\T$ with arclength $\abs I=c/m$ and angular distance at least $\gamma/m$ from the two real points $z=\pm1$, and every $t\in(0,1]$,
\begin{equation}\label{eq:singlearc}
  \Prob\Bigl\{\min_{z\in I}\abs{G_m(z)}\le t\Bigr\}\;\le\;C_\gamma\,(1+c)\,\frac{t}{\sqrt m}.
\end{equation}
The exponent of $t$ is $1$ and that of $m$ is $-\tfrac12$; the bound is the curve-crossing rate, not the pointwise $t^2/m$ rate.
\end{theorem}

The proof, given below, is a one-dimensional Kac--Rice crossing count for
\(|G_m(e^{i\theta})|^2\).  Proposition~\ref{prop:singlearc-log} records the
older occupation-measure proof with a logarithmic speed loss because that
argument is robust and is reused for the speed-conditional bounded-density
estimates.

\begin{theorem}[Uniform boundary small-ball]\label{thm:uniform}
With $G_m$ as above (independent standard real Gaussian coefficients), there is an absolute constant $C'>0$ such that
\begin{equation}\label{eq:uniform}
  \Prob\Bigl\{\min_{z\in\T}\abs{G_m(z)}\le t\Bigr\}\;\le\;C'\,\sqrt m\,t
  \qquad(0<t\le1),
\end{equation}
equivalently
\(\Prob\{\min_{\T}\abs{G_m}\le\tau/\sqrt m\}\le C'\tau\). The expected
constant is governed by the Kac--Rice first moment for down-crossings, of order
\(t\sqrt m\), in agreement with the minimum-modulus picture of
\cite{KMS,YakirZeitouni,CookNguyen}; away from the real points this first
moment is two-sided by Proposition~\ref{prop:away-firstmoment}, in the fixed
bulk its leading density is computed in Proposition~\ref{prop:bulk-kr-constant},
and Corollary~\ref{cor:bulk-continuum-intensity} gives the resulting
continuum-scale fixed-bulk first-moment limit at \(t=\tau/\sqrt N\).
Corollary~\ref{cor:bulk-prob-upper} converts that first moment into an explicit
fixed-bulk probability upper bound with the down-crossing constant, while
Proposition~\ref{prop:bulk-secondmoment-transfer} proves the complementary
conditional lower-tail transfer from a second-factorial-moment estimate for the
down-crossing count, and Corollary~\ref{cor:bulk-firstorder-sandwich} sandwiches
the fixed-bulk probability to first order in \(\tau\).
Corollary~\ref{cor:global-secondmoment-transfer} records the endpoint-cutoff
removal: if the same factorial second-moment estimate is uniform over fixed
bulks as \(\delta\downarrow0\), then the whole-circle lower tail has leading
constant \(\sqrt{\pi/3}\).  At the two real endpoints the local
order is already sharp, by Proposition~\ref{prop:endpoint-lower};
Corollaries~\ref{cor:endpoint-two-point},
\ref{cor:endpoint-continuum-negligible}, and
\ref{cor:endpoint-neighborhood-continuum} give the exact bivariate Gaussian
endpoint formula, the sharp second-order expansion, and the fact that at
\(t=\tau/\sqrt N\) both the two endpoints and their natural \(O(1/m)\)
neighborhoods have only \(O_\tau(N^{-1})\) probability.
Corollary~\ref{cor:global-prob-upper} combines this endpoint separation with
the fixed-bulk upper bound to give the sharp whole-circle continuum-scale upper
constant \(\sqrt{\pi/3}\), and
Corollary~\ref{cor:global-firstorder-sandwich} gives the conditional
whole-circle first-order law under the uniform factorial second-moment input.
The missing part of the matching global lower tail is therefore precisely the
anti-clustering and decorrelation estimate that bounds the second factorial
moment of the \(O(m)\) non-real down-crossing opportunities.
\end{theorem}

For the finite-section theorem one does not need the sharp log-free estimate at
all scales.  It is enough to know that a polynomially small boundary minimum is
unlikely.  The following corollary is the polynomial-scale form used in
Theorem~\ref{thm:gaussian-main}.

\begin{corollary}[Uniform Gaussian boundary small ball at polynomial scales]\label{prop:uniform-polylog}
Let \(G_m(z)=\sum_{j=0}^m\xi_jz^j\) with \(\xi_j\) independent standard real Gaussians. For every \(A,B>0\) there are constants \(C_{A,B}<\infty\) and \(m_0(A,B)\) such that, for every \(m\ge m_0(A,B)\) and every \(m^{-B}\le t\le1\),
\begin{equation}\label{eq:uniform-polylog}
  \Prob\Bigl\{\min_{z\in\T}\abs{G_m(z)}\le t\Bigr\}
  \le C_{A,B}\,\sqrt m\,t.
\end{equation}
In particular, Hypothesis~\ref{hyp:boundary-law} holds for standard real Gaussian coefficients.
\end{corollary}

\begin{hypothesis}[Polynomial-scale uniform boundary small ball for the coefficient law]\label{hyp:boundary-law}
For the coefficient law in Theorem~\ref{thm:main}, and the associated
\(G_m(z)=\sum_{j=0}^m\xi_jz^j\), for every \(A,B>0\) there are constants
\(C_{A,B,K}<\infty\) and \(m_0(A,B,K)\), where \(K\) denotes the density bound
of the law, such that, for every
\(m\ge m_0(A,B,K)\) and every \(m^{-B}\le t\le1\),
\[
  \Prob\Bigl\{\min_{z\in\T}\abs{G_m(z)}\le t\Bigr\}
  \le C_{A,B,K}\sqrt{m\log m}\,t+C_{A,B,K}m^{-A}.
\]
\end{hypothesis}

Theorem~\ref{thm:uniform} implies Hypothesis~\ref{hyp:boundary-law} for
Gaussian coefficients, and Corollary~\ref{prop:uniform-polylog} records the
polynomial-scale form used below.  For non-Gaussian
bounded-density laws, Hypothesis~\ref{hyp:boundary-law} is the natural
polynomial-scale boundary small-ball input.  The finite-dimensional
part of that input is unconditional: Proposition~\ref{prop:nongaussian-grid}
below proves the correct fixed-point and deterministic-grid small-ball estimates
for every bounded-density law, and Corollary~\ref{cor:bd-uniform-speed} gives the
corresponding continuum bound after intersecting with a deterministic speed
cutoff.  Corollary~\ref{cor:bd-speedtail-to-boundary} records the exact
conditional upgrade: a polynomially small speed-tail estimate at the natural
\(\sqrt{\log m}\) scale implies Hypothesis~\ref{hyp:boundary-law};
Corollaries~\ref{cor:bd-polymoment-boundary} and
\ref{cor:bd-subexp-boundary} check that speed-tail input for
all-polynomial-moment and subexponential bounded-density laws, respectively
(hence also for subgaussian ones). For the finite-section conditioning theorem
itself, Theorem~\ref{thm:main} needs only a much cruder
polynomial speed cutoff, available from variance one. What remains open in full
bounded-density generality, without a tail assumption, is the sharp continuum
boundary law itself: proving a natural-scale speed/crossing input, or replacing
it by a genuine local-limit argument.

The full boundary law is stronger than the finite-section conditioning theorem
needs.  The conditioning proof uses only three high-probability inputs: circle
clearance, one polynomial boundary-minimum event, and a crude coefficient-size
bound.  Corollary~\ref{cor:variance-boundary-separation} supplies that one
boundary-minimum event for every bounded-density variance-one law, at the cost
of a deliberately polynomial speed cutoff in its proof; the exceptional
probability is \(O_{A,K}(m^{-2(A-1)/3})\) for the threshold \(m^{-A}\).  Thus the only
coefficient-law input left in the conditioning theorem, beyond bounded density
and variance one, is circle clearance.
The point is elementary but important: the final probability is measured as
\[
  \Prob\{\nu=r\}-o(1),
\]
not as a conditional probability given \(\nu=r\).  Hence a high-probability clearance event may be intersected with the zero-winding event at additive cost \(o(1)\), without any independence or transport of the clearance law through the event \(\{\nu=r\}\).  The finite-section threshold is deterministic once a clearance floor is fixed; see Proposition~\ref{prop:fs-uniform}.

We therefore isolate the remaining genuine probabilistic input.

\begin{hypothesis}[Circle-clearance estimate for the coefficient law]\label{hyp:clearance}
Let the coefficients be as in Theorem~\ref{thm:main} below. With $\rho:=\min\{\,\abs{\,\abs\zeta-1\,}:G_m(\zeta)=0\}$ the distance of the zero cloud of $G_m$ to $\T$, there are a constant \(C_0>0\) and events \(\mathcal C_m\) with \(\Prob(\mathcal C_m)=1-o(1)\) such that, on \(\mathcal C_m\),
\[
  \rho\ge\frac{1}{C_0\,m^{2}(\log m)^{3}}
\]
and in particular \(G_m\) has no zeros on \(\T\).
\end{hypothesis}

For standard real Gaussian coefficients this hypothesis is exactly the circle-clearance theorem of Konyagin--Schlag \cite[Theorem~1.1]{KMS}, as used in \cite{Parity}.  For other coefficient laws, Theorem~\ref{thm:main} assumes this same clearance input explicitly.

\begin{theorem}[Balanced conditioning on the stable event under clearance]\label{thm:main}
Let $r=s$, $m=2r$, and let the coefficients $\xi_{-r},\dots,\xi_r$ be independent, mean zero, variance one, with a density bounded by $K$ (the Gaussian case is included). Set
\[
  \nu:=\#\{z\in\D:G_m(z)=0\}.
\]
Assume Hypothesis~\ref{hyp:clearance}. Then, for every \(b>0\), there exist a
deterministic function \(n_{0,b}(m)<\infty\) and events
\(\Omega_m\) with \(\Prob(\Omega_m)=1-o(1)\) such that, on
\(\Omega_m\cap\{\nu=r\}\), the symbol \(P_{r,r}\) has no zeros on \(\T\), has
winding number zero, and for all \(n\ge n_{0,b}(m)\)
\begin{equation}\label{eq:main}
  \smin\bigl(T_n(P_{r,r})\bigr)\;\ge\;n^{-b}
  \qquad\text{and}\qquad
  \kdn\bigl(T_n(P_{r,r})\bigr)\;\le\;n^{\,b}.
\end{equation}
Equivalently, for any prescribed polynomial exponent, the conditioning conclusion
holds with probability at least \(\Prob\{\nu=r\}-o(1)\), after increasing the
deterministic finite-section threshold. The theorem does not assert that a
balanced random symbol is stable with high probability; that would require a
separate local theorem for \(\nu\). This is the zero-winding counterpart to the
exponential ill-conditioning of the strongly unbalanced regime
\cite[Thm.~1.2]{Parity}.
\end{theorem}

\begin{corollary}[Quantitative stable-event conditioning]\label{cor:quant-main}
In the setting of Theorem~\ref{thm:main}, assume more quantitatively that for
some \(C_0>0\), events \(\mathcal C_m\), and numbers \(q_m\in[0,1]\),
\[
  \Prob(\mathcal C_m^c)\le q_m,\qquad
  \rho\ge\frac{1}{C_0m^2(\log m)^3}\quad\text{on }\mathcal C_m .
\]
Fix \(A_0>2\), \(b>0\), and a deterministic envelope \(R_m>0\).  Then there
are deterministic thresholds \(n_{0,b,A_0,R}(m)<\infty\) and events
\(\Omega_{m,A_0,R}\) such that
\[
  \Prob(\Omega_{m,A_0,R}^c)
  \le q_m+C_{A_0,K}m^{-2(A_0-1)/3}+\frac{m+1}{R_m^2},
\]
and, on \(\Omega_{m,A_0,R}\cap\{\nu=r\}\), the symbol \(P_{r,r}\) has no zeros on
\(\T\), has winding number zero, and for all \(n\ge n_{0,b,A_0,R}(m)\),
\[
  \smin\bigl(T_n(P_{r,r})\bigr)\ge n^{-b},
  \qquad
  \kdn\bigl(T_n(P_{r,r})\bigr)\le n^b .
\]
Consequently the same conclusion holds with probability at least
\[
  \Prob\{\nu=r\}-q_m-C_{A_0,K}m^{-2(A_0-1)/3}-\frac{m+1}{R_m^2}.
\]
In particular, for every \(D>0\), choosing \(R_m=m^{(D+1)/2}\) replaces the last
loss by \(2m^{-D}\) for \(m\ge1\), after changing only the deterministic
threshold.
\end{corollary}

\begin{corollary}[Prescribed polynomial stable-event loss]
\label{cor:quant-main-polynomial-loss}
In the setting of Corollary~\ref{cor:quant-main}, fix \(b>0\) and \(D>0\).
Then there are deterministic thresholds \(n_{0,b,D}(m)<\infty\) and events
\(\Omega_{m,D}\) such that
\[
  \Prob(\Omega_{m,D}^c)\le q_m+C_{D,K}m^{-D},
\]
and, on \(\Omega_{m,D}\cap\{\nu=r\}\), the symbol \(P_{r,r}\) has no zeros on
\(\T\), has winding number zero, and for all \(n\ge n_{0,b,D}(m)\),
\[
  \smin\bigl(T_n(P_{r,r})\bigr)\ge n^{-b},
  \qquad
  \kdn\bigl(T_n(P_{r,r})\bigr)\le n^b .
\]
Consequently the same conclusion holds with probability at least
\[
  \Prob\{\nu=r\}-q_m-C_{D,K}m^{-D}.
\]
\end{corollary}

\begin{proof}
Apply Corollary~\ref{cor:quant-main} with
\[
  A_0=1+\frac32(D+1),\qquad R_m=m^{(D+2)/2}.
\]
Then \(A_0>2\) and \(2(A_0-1)/3=D+1\), so the boundary-separation loss is
\(C_{D,K}m^{-(D+1)}\).  Also
\[
  \frac{m+1}{R_m^2}\le 2m^{-(D+1)}\le 2m^{-D}\qquad(m\ge1).
\]
Absorbing constants and renaming the deterministic threshold gives the stated
event and probability bound.
\end{proof}

\begin{corollary}[Conditional stable-event probability]
\label{cor:quant-main-conditional}
In the setting of Corollary~\ref{cor:quant-main-polynomial-loss}, suppose
\(p_m:=\Prob\{\nu=r\}>0\).  Let \(\mathcal G_{m,b,D}\) be the event that
\(P_{r,r}\) has no zeros on \(\T\), has winding number zero, and for all
\(n\ge n_{0,b,D}(m)\),
\[
  \smin\bigl(T_n(P_{r,r})\bigr)\ge n^{-b},
  \qquad
  \kdn\bigl(T_n(P_{r,r})\bigr)\le n^b .
\]
Then
\[
  \Prob\left(\mathcal G_{m,b,D}\,\middle|\,\nu=r\right)
  \ge 1-\frac{q_m+C_{D,K}m^{-D}}{p_m}.
\]
In particular, if \(p_m\ge p_*>0\) along a subsequence, then the conditional
failure probability on that subsequence is at most
\(p_*^{-1}(q_m+C_{D,K}m^{-D})\).
\end{corollary}

\begin{proof}
Corollary~\ref{cor:quant-main-polynomial-loss} gives an event
\(\Omega_{m,D}\) with complement probability at most \(q_m+C_{D,K}m^{-D}\) and
\(\Omega_{m,D}\cap\{\nu=r\}\subset\mathcal G_{m,b,D}\cap\{\nu=r\}\).  Hence
\[
  \Prob\left(\mathcal G_{m,b,D}^c\cap\{\nu=r\}\right)
  \le q_m+C_{D,K}m^{-D}.
\]
Dividing by \(p_m=\Prob\{\nu=r\}\) proves the conditional estimate.
\end{proof}

\begin{theorem}[Gaussian balanced conditioning]\label{thm:gaussian-main}
In the setting of Theorem~\ref{thm:main}, assume the coefficients are
independent standard real Gaussians. Then the conclusion of
Theorem~\ref{thm:main} holds unconditionally: Hypothesis~\ref{hyp:clearance} is
supplied by the Konyagin--Schlag clearance theorem \cite{KMS}.
\end{theorem}

\begin{corollary}[Subexponential boundary law and conditioning]\label{cor:subexp-main}
In the setting of Theorem~\ref{thm:main}, assume in addition that
\(\norm{\xi_j}_{\psi_1}\le\kappa\) uniformly in \(j\) and \(m\). Then
Hypothesis~\ref{hyp:boundary-law} holds for this coefficient law. If the
coefficient law also satisfies the circle-clearance estimate
Hypothesis~\ref{hyp:clearance}, then the conclusion of Theorem~\ref{thm:main}
holds. Thus, for subexponential bounded-density laws -- in particular for
subgaussian bounded-density laws -- the full boundary-law input is available;
this is stronger than what the conditioning proof under clearance ultimately
needs.
\end{corollary}

\begin{corollary}[All-moment bounded-density boundary law and conditioning]\label{cor:allmoment-main}
In the setting of Theorem~\ref{thm:main}, assume instead that for every
\(p<\infty\) the coefficients satisfy
\[
  \sup_{m\ge1}\max_{0\le j\le m}\Exp\abs{\xi_j}^{p}<\infty .
\]
Then Hypothesis~\ref{hyp:boundary-law} holds for this coefficient law. If the
coefficient law also satisfies the circle-clearance estimate
Hypothesis~\ref{hyp:clearance}, then the conclusion of Theorem~\ref{thm:main}
holds. Thus the exponential-tail assumption is not the boundary input required
for the full polynomial-scale boundary law; uniformly bounded moments of all
orders already supply the needed natural-scale continuum speed tail.
\end{corollary}

Indeed, Corollaries~\ref{cor:bd-subexp-boundary} and
\ref{cor:bd-polymoment-boundary} verify Hypothesis~\ref{hyp:boundary-law} from,
respectively, the subexponential derivative tail and the Fuk--Nagaev
polynomial-moment derivative tail. Theorem~\ref{thm:main} shows that the weaker
conditioning conclusion only needs the one-threshold substitute packaged in
Corollary~\ref{cor:variance-boundary-separation}, while
Hypothesis~\ref{hyp:clearance} is assumed here.

The boundary estimate over the whole circle is what upgrades ``one arc is good'' to ``all arcs are good simultaneously'': a single favorable arc cannot give the finite-section bound \eqref{eq:main}, because Proposition~\ref{prop:fs} needs the full boundary minimum $\min_{z\in\T}\abs{G_m(z)}$.  For the full polynomial-scale boundary law this remains the role of Hypothesis~\ref{hyp:boundary-law}; for the conditioning conclusion under clearance, the tail-free polynomial separation in Corollary~\ref{cor:variance-boundary-separation} is already sufficient.

\subsection{The bridge to the parity obstruction}\label{sub:bridge}

The boundary evaluations that drive Theorem~\ref{thm:singlearc} are the same two evaluations that drive the parity obstruction of \cite{Parity}. The parity of $\nu$ is the parity of the number of real zeros of $G_m$ in $(-1,1)$, which is controlled by the signs of the two boundary values $G_m(1)=\sum_j\xi_j$ and $G_m(-1)=\sum_j(-1)^j\xi_j$, the evaluations of $G_m$ at the two real points of $\T$. The boundary small-ball estimate is the quantitative refinement of the same observable: instead of asking only for the \emph{sign} of $G_m(\pm1)$, it asks how close $\abs{G_m}$ comes to zero along $\T$.

This shared dependence on boundary evaluations is precisely where the two papers meet, and the meeting point is the Bernoulli case. In the balanced setup of this paper \(m=2r\), so \(m+1\) is odd and the two sums \(G_m(1)\) and \(G_m(-1)\) cannot vanish for Rademacher coefficients. Nevertheless their lattice spacing is a boundary obstruction to density-type small-ball estimates at very fine scales. In adjacent parities, where an even number of signs is present, exact boundary zeros occur with polynomial probability at \(z=\pm1\). This is the Toeplitz analogue of the obstruction recorded in \cite[Rem.~6.2]{PhaseRetrieval}: for $\pm1$ measurement matrices the least singular value over square submatrices can vanish by a coordinate hyperplane collision, so continuity of the entry law is necessary for the clean lower-bound theorem. We develop this correspondence in Section~\ref{sec:bridge}.

\subsection{What is proved and what is conjectured}\label{sub:honest}

Because the bounded-density theorem still separates the coefficient-law inputs, we state plainly where the line falls.
\begin{itemize}
\item \textbf{Proved.} The finite-section symbol reduction \eqref{eq:fs-reduction} on the zero-winding event (Propositions~\ref{prop:fs} and \ref{prop:fs-uniform}), the sharp Gaussian single-arc bound away from \(z=\pm1\) (Theorem~\ref{thm:singlearc}), the away-from-endpoint two-sided first moment for crossings, the explicit fixed-bulk Kac--Rice down-crossing constant \(1/\sqrt{12\pi}\), the fixed-bulk continuum-scale crossing limit at \(t=\tau/\sqrt N\), the corresponding fixed-bulk first-moment probability upper bound, the Paley--Zygmund transfer from a bulk down-crossing factorial second moment to a fixed-bulk lower tail, the resulting fixed-bulk first-order probability sandwich, the sharp global continuum-scale probability upper bound, the weak-\(L^1\) inverse-scaled-minimum endpoint, single and finite-family inverse-minimum quantile bounds, finite-family clearance order-statistic quantile and total/sparse bad-count schedule bounds, finite-family Stieltjes-weighted/powered clearance-deficit summability, normalized deficit-budget schedules, Stieltjes inverse-gauge mass schedules, inverse-gauge confidence budgets, gauge-menu budgets, and aggregate gauge-menu summability, finite-family subcritical inverse-power and truncated/logarithmic endpoint moments, the logarithmic critical truncated inverse-moment bound, logarithmic inverse-moment bounds of all positive orders, and subcritical negative-moment bounds for the scaled global boundary minimum, the conditional global first-order probability sandwich under a uniform factorial moment bound, and the uniform-in-\(\delta\) transfer from fixed bulks to the whole-circle lower-tail constant \(\sqrt{\pi/3}\) (Propositions~\ref{prop:away-firstmoment}, \ref{prop:bulk-kr-constant}, and \ref{prop:bulk-secondmoment-transfer}, Corollaries~\ref{cor:bulk-continuum-intensity}, \ref{cor:bulk-prob-upper}, \ref{cor:bulk-firstorder-sandwich}, \ref{cor:global-prob-upper}, \ref{cor:global-negative-moments}, \ref{cor:global-inverse-weak}, \ref{cor:global-inverse-quantiles}, \ref{cor:global-inverse-family-quantiles}, \ref{cor:global-inverse-family-counts}, \ref{cor:global-inverse-family-order-quantiles}, \ref{cor:global-inverse-family-borel}, \ref{cor:global-inverse-family-count-sum}, \ref{cor:global-inverse-family-deficit-sum}, \ref{cor:global-inverse-family-deficit-budget}, \ref{cor:global-inverse-family-deficit-powers}, \ref{cor:global-inverse-family-deficit-weighted}, \ref{cor:global-inverse-family-subcritical-sum}, \ref{cor:global-inverse-family-gauge-sum}, \ref{cor:global-inverse-family-gauge-confidence}, \ref{cor:global-inverse-family-gauge-menu}, \ref{cor:global-inverse-family-gauge-menu-sum}, \ref{cor:global-inverse-family-count-borel}, \ref{cor:global-inverse-family-moments}, \ref{cor:global-inverse-family-subcritical}, \ref{cor:global-inverse-truncated}, \ref{cor:global-inverse-logmoments}, \ref{cor:global-firstorder-sandwich}, and \ref{cor:global-secondmoment-transfer}), the endpoint lower bound showing the real-point arc scale is sharp, the exact two-endpoint Gaussian formula and sharp second-order expansion, the \(O_\tau(N^{-1})\) endpoint and natural-endpoint-neighborhood contributions at the continuum scale, and the resulting global endpoint floor (Proposition~\ref{prop:endpoint-lower} and Corollaries~\ref{cor:endpoint-two-point}, \ref{cor:endpoint-continuum-negligible}, \ref{cor:endpoint-neighborhood-continuum}, and \ref{cor:endpoint-global-floor}), the sharp uniform Gaussian boundary small-ball bound (Theorem~\ref{thm:uniform}), the exact Gaussian non-real and full root-grid law, including endpoint survival factors, and its pointwise Rayleigh-limit/moment-expansion and exact continuum-scale grid-coefficient corollaries (Proposition~\ref{prop:gaussian-root-grid} and Corollaries~\ref{cor:gaussian-root-grid-rayleigh} and \ref{cor:gaussian-root-grid-misses}), the fixed-point/deterministic-grid bounded-density small-ball estimate, deterministic and independent-random subquadratic grid-miss consequences, matching quadratic-resolution detection under a speed cutoff, independent uniform random-grid detection with one extra covering logarithm, and Gaussian quadratic-log high-probability grid detection (Proposition~\ref{prop:nongaussian-grid}, Corollaries~\ref{cor:subquadratic-grid-misses}, \ref{cor:random-subquadratic-grid-misses}, \ref{cor:uniform-random-grid-detects}, and \ref{cor:gaussian-grid-detects}, and Lemma~\ref{lem:grid-speed-detection}), the bounded-density local and uniform estimates under a deterministic speed cutoff (Proposition~\ref{prop:bd-speed} and Corollary~\ref{cor:bd-uniform-speed}), the variance-only polynomial derivative cutoff (Lemma~\ref{lem:variance-speed}), the resulting bounded-density conditioning theorem under circle clearance and its quantitative stable-event loss (Theorem~\ref{thm:main} and Corollary~\ref{cor:quant-main}), the reduction from a polynomial speed tail to the full polynomial-scale boundary law and to high-probability grid detection (Corollaries~\ref{cor:bd-speedtail-to-boundary} and \ref{cor:speedtail-grid-detects}), the verification of that speed-tail hypothesis for all-polynomial-moment and subexponential bounded-density laws (Corollaries~\ref{cor:bd-polymoment-boundary} and \ref{cor:bd-subexp-boundary}), the endpoint, order-\(q\) root-of-unity, fixed-radius lower, rational-grid, bounded-denominator, fixed-root, primitive-quadratic, primitive-quadratic second-order/overlap, degree-four, finite-cyclotomic-layer, and bounded-denominator exact-zero Rademacher lattice anti-concentration/atom scales, including the exact parity-residue criterion for zero atoms, the peeled layer formula, the first two exact-zero coefficients, the full unpeeled expansion through \(N^{-2}\), and the endpoint-inclusive bounded-denominator expansion through order \(N^{-1}\) (Propositions~\ref{prop:rademacher-endpoint}, \ref{prop:rademacher-root-unity}, \ref{prop:rademacher-root-fixed-radius-lower}, and \ref{prop:rademacher-small-denom-lower}, and Corollaries~\ref{cor:rademacher-root-grid}, \ref{cor:rademacher-denom-grid}, \ref{cor:rademacher-finite-layer-zero}, \ref{cor:rademacher-quadratic-atoms}, \ref{cor:rademacher-quadratic-overlaps}, \ref{cor:rademacher-quadratic-second-order}, \ref{cor:rademacher-degree-four-atoms}, \ref{cor:rademacher-denom-zero}, \ref{cor:rademacher-denom-zero-layer}, \ref{cor:rademacher-denom-zero-leading}, \ref{cor:rademacher-denom-zero-degree-four-leading}, \ref{cor:rademacher-denom-zero-second-order}, and \ref{cor:rademacher-denom-zero-with-endpoints}), and the unconditional Gaussian deduction in Theorem~\ref{thm:gaussian-main}.
The finite-family endpoint package also records the weak tail in
Corollary~\ref{cor:global-inverse-family-weak}, bad-clearance counts in
Corollary~\ref{cor:global-inverse-family-counts}, clearance order statistics in
Corollary~\ref{cor:global-inverse-family-order-quantiles}, and the summable
growing-batch first-bad, total bad-mass, deficit-mass, normalized
deficit-budget, powered and jump/Stieltjes-weighted deficit, subcritical inverse-mass,
Stieltjes inverse-gauge, inverse-gauge confidence, gauge-menu, aggregate gauge-menu,
and sparse bad-count schedules of
Corollaries~\ref{cor:global-inverse-family-borel},
\ref{cor:global-inverse-family-count-sum},
\ref{cor:global-inverse-family-deficit-sum},
\ref{cor:global-inverse-family-deficit-budget},
\ref{cor:global-inverse-family-deficit-powers},
\ref{cor:global-inverse-family-deficit-weighted},
\ref{cor:global-inverse-family-subcritical-sum},
\ref{cor:global-inverse-family-gauge-sum},
\ref{cor:global-inverse-family-gauge-confidence},
\ref{cor:global-inverse-family-gauge-menu},
\ref{cor:global-inverse-family-gauge-menu-sum}, and
\ref{cor:global-inverse-family-count-borel}.
The same variance-only cutoff also gives polynomial-grid continuum-excursion
detection with failure probability \(O_\delta(m^{1-2\delta})\)
(Corollaries~\ref{cor:variance-speed-confidence} and
\ref{cor:variance-grid-detects}); its finite-confidence form gives
\(\eta\)-level grid detection at mesh
\(O(\tau\sqrt\eta/(m^2\sqrt N))\) and bounded-density boundary loss
\(\eta+C_{B,K}mt/\sqrt\eta\)
(Corollaries~\ref{cor:variance-grid-confidence} and
\ref{cor:variance-boundary-confidence}); in fixed-power form, every
uniform \(N^\beta\)-grid with \(\beta>5/2\) has miss probability
\(O_\beta(N^{5-2\beta})\)
(Corollary~\ref{cor:variance-grid-fixed-power}), hence
  \(O_A(m^{-A})\) failure after increasing the grid size to
\(O_{A,\tau}(N^{(A+5)/2})\)
(Corollary~\ref{cor:variance-grid-polynomial-confidence}), and the combined
subquadratic-obstruction/fine-grid-sufficiency sampling bracket
(Corollary~\ref{cor:grid-sampling-bracket}).
Under the natural polynomial derivative-speed tail, the fine-grid side improves
to \(O_{A,\tau}(N^2\sqrt{\log N})\) points with the same \(O_A(m^{-A})\)
miss probability, while the subquadratic obstruction is unchanged
(Corollary~\ref{cor:natural-speed-sampling-bracket}).
The same natural-speed scale holds for stratified one-point-per-cell random
grids and random shifts of a uniform lattice, without the iid coupon-collector
overhead (Corollary~\ref{cor:stratified-grid-detects}).
Under bounded density it also gives the tail-free polynomial separation
\(\min_{\T}\abs{G_m}>m^{-A}\) with exceptional probability
\(O_{A,K}(m^{-2(A-1)/3})\) for every \(A>1\)
(Corollary~\ref{cor:variance-boundary-separation}).
\item \textbf{Conjectured/conditional.} The full continuum boundary-law version of universality beyond Gaussian coefficients remains open at bounded-density generality, without tail or all-moment input. The matching Gaussian lower-tail law is now reduced, on every fixed bulk \(B_\delta\), to the factorial second-moment/decorrelation estimate in \eqref{eq:bulk-factorial-secondmoment}; Proposition~\ref{prop:bulk-secondmoment-transfer} proves the transfer from that estimate to the lower-tail probability, Corollary~\ref{cor:bulk-firstorder-sandwich} combines it with the sharp upper bound to give the fixed-bulk first-order law as \(\tau\downarrow0\), Corollary~\ref{cor:global-prob-upper} gives the sharp whole-circle upper constant, and Corollary~\ref{cor:global-firstorder-sandwich} shows that a version uniform as \(\delta\downarrow0\) gives the whole-circle first-order law as \(\tau\downarrow0\). The conditioning theorem no longer needs that full law once circle clearance is assumed; for coefficient laws outside the available Konyagin--Schlag hypotheses the remaining conditioning input is the circle-clearance estimate in Hypothesis~\ref{hyp:clearance}. There is no additional clearance-coupling hypothesis: the intersection with \(\{\nu=r\}\) costs only the unconditional exceptional probability.
\end{itemize}
We do not use an asymptotic formula for $\Var(\nu)$ here; the disk-count fluctuation problem belongs to the companion Toeplitz-parity analysis.

\subsection*{Relation to the companions and the literature}
This paper is a companion to \cite{Parity} and \cite{PhaseRetrieval}. From \cite{Parity} we take the operator-theoretic framework \eqref{eq:half}, the unbalanced exponential instability, and the Konyagin--Schlag clearance \cite{KMS}; we complement it by isolating the zero-winding conditioning mechanism in the balanced regime. From \cite{PhaseRetrieval} we take the negative-second-moment small-ball engine and the good-family second moment philosophy, transplanted from row submatrices to boundary arcs. The conditioning of \emph{dense} random Toeplitz and circulant matrices is treated in the references collected in \cite{Parity}; the operator theory of banded symbols is \cite{BG}.

\section{Setup and the negative second moment engine}\label{sec:setup}

Throughout, $m=2r$ is even, $\xi_0,\dots,\xi_m$ are independent with $\Exp\xi_j=0$, $\Exp\xi_j^2=1$, and (unless stated otherwise) with a density bounded by $K$; $G_m(z)=\sum_{j=0}^m\xi_jz^j$. We write $\T=\{z:\abs z=1\}$ and $\D=\{z:\abs z<1\}$, and parametrize $\T$ by $z=e^{i\theta}$. The symbol is $P_{r,r}(z)=z^{-r}G_m(z)$, so $\abs{P_{r,r}(z)}=\abs{G_m(z)}$ on $\T$. Constants $c,C,\dots$ are absolute or depend only on $K$, and may change line to line.

\subsection{Boundary values as a Gaussian field}
Fix points $z=e^{i\theta},w=e^{i\phi}\in\T$. Then $G_m(z)=\sum_j\xi_je^{ij\theta}$, and in the real-coefficient case the real and imaginary parts
\[
  \Re G_m(e^{i\theta})=\sum_{j=0}^m\xi_j\cos(j\theta),\qquad
  \Im G_m(e^{i\theta})=\sum_{j=0}^m\xi_j\sin(j\theta)
\]
are linear forms in the $\xi_j$ with covariance structure governed by \eqref{eq:dirichlet}. The diagonal is $\Exp\abs{G_m(z)}^2=m+1$, and the covariance between two points is the Dirichlet kernel
\begin{equation}\label{eq:cov-again}
  \Exp\bigl[G_m(z)\overline{G_m(w)}\bigr]=\sum_{j=0}^m e^{ij(\theta-\phi)}=D_m(\theta-\phi),
  \qquad
  \abs{D_m(\psi)}\le\min\Bigl\{m+1,\frac{1}{\abs{\sin(\psi/2)}}\Bigr\}.
\end{equation}
Thus $\abs{D_m(\theta-\phi)}\le 1/\abs{\sin((\theta-\phi)/2)}$ is $o(m+1)$ once $\abs{\theta-\phi}\gg1/m$: points more than a Dirichlet width apart are weakly correlated. This is the decorrelation that organizes $\T$ into $\Theta(m)$ effectively independent arcs and is the analogue, in the present problem, of the weak overlap of well-separated row sets in \cite{PhaseRetrieval}.

\subsection{The negative second moment identity on an arc}
We will reduce the minimum of $\abs{G_m}$ over an arc to a one-dimensional small-ball event by the same device as \eqref{eq:nsm}, applied not to a single matrix but to a well-chosen finite sample of the field on the arc. Fix an arc $I$ of length $\abs I\ge c/m$ and sample $G_m$ at $L\asymp m\abs I\ge c$ equally spaced points $z_1,\dots,z_L\in I$ at spacing $\asymp1/m$. Consider the random vector $\mathbf g=(G_m(z_1),\dots,G_m(z_L))^\top=V\xi$, where $V\in\C^{L\times(m+1)}$ has entries $V_{\ell j}=z_\ell^{\,j}$ (a Vandermonde block) and $\xi=(\xi_0,\dots,\xi_m)^\top$.

\begin{lemma}[Sample covariance is well-conditioned on a short arc]\label{lem:vander}
Fix \(c'<\infty\) and spacing constants \(0<a\le A<\infty\).  There are constants \(0<\lambda_-\le\lambda_+<\infty\) and \(m_0<\infty\), depending only on \(a,A,c'\), such that the following holds for \(m\ge m_0\).  If \(z_\ell=e^{i\theta_\ell}\) are sample points on an arc of length at most \(c'/m\), with \(L\le A c'\) and pairwise angular separation at least \(a/m\), then the \(L\times L\) Gram matrix \(\Sigma:=\Exp[\mathbf g\,\mathbf g^*]=VV^*\) satisfies
\begin{equation}\label{eq:vander}
  \lambda_-(m+1)\,\Id_L\;\preceq\;\Sigma\;\preceq\;\lambda_+(m+1)\,\Id_L.
\end{equation}
\end{lemma}

\begin{proof}
The entries of $\Sigma$ are $\Sigma_{\ell\ell'}=D_m(\theta_\ell-\theta_{\ell'})$ by \eqref{eq:cov-again}.  Write \(\theta_\ell=\theta_0+x_\ell/m\).  The parameters \(x_\ell\) lie in a compact interval of length \(c'\), are separated by at least \(a\), and there are only boundedly many of them.  Uniformly over this compact parameter set,
\[
  \frac{1}{m+1}D_m\!\left(\frac{x_\ell-x_{\ell'}}{m}\right)
  \longrightarrow
  \int_0^1 e^{iu(x_\ell-x_{\ell'})}\,du .
\]
The limiting matrix is the Gram matrix of the distinct exponentials \(u\mapsto e^{iux_\ell}\) in \(L^2[0,1]\).  It is positive definite, and compactness of the separated configuration space gives a positive lower bound for its smallest eigenvalue.  The upper bound is immediate from bounded \(L\) and bounded entries.  The same bounds therefore hold for \(\Sigma/(m+1)\) for all \(m\ge m_0(a,A,c')\).
\end{proof}

For circular complex Gaussian coefficients, Lemma~\ref{lem:vander} says that on a short arc the sampled field $\mathbf g$ is a non-degenerate Gaussian vector in $\C^L$ at scale $\sqrt{m+1}$, with condition number $O(1)$.  In the real-coefficient problem of the present paper, the Hermitian covariance alone does not give full real \(2L\)-dimensional nondegeneracy; the pointwise real nondegeneracy away from \(z=\pm1\) is supplied by Lemma~\ref{lem:onepoint}, and the endpoint degeneration is handled separately by Lemma~\ref{lem:endpoint-occ}.  The Vandermonde lemma is included to record the decorrelation geometry behind the negative-second-moment analogy.

\section{Boundary small-ball estimates and targets}\label{sec:smallball}

We record the route from pointwise anti-concentration to arc estimates and then
to the remaining uniform problem.  The mechanism is not the discrete strategy
\eqref{eq:nsm}--\eqref{eq:onerow} applied pointwise, which would give the
pointwise \(t^2/m\) rate and lose the arc geometry; it is a one-dimensional
curve-crossing argument.  The sharp Gaussian single-arc theorem is proved below
by Kac--Rice.  We also keep the older occupation-measure argument, based on a
derivative cutoff and Lemma~\ref{lem:arcsampling}, because it is the robust
template for bounded-density laws.

\subsection{The three ingredients}

\begin{lemma}[One-point anti-concentration]\label{lem:onepoint}
Let $\xi_j$ be independent standard real Gaussians and fix $z=e^{i\theta}\in\T$. Then \(G_m(z)\) is a centered real Gaussian vector in \(\R^2\). For every fixed \(\gamma>0\) there are constants \(c_\gamma,C_\gamma>0\) and \(m_0(\gamma)\) such that, for \(m\ge m_0(\gamma)\), if \(z\) is separated from \(\pm1\) by an angular distance at least \(\gamma/m\), its covariance has both eigenvalues at least \(c_\gamma(m+1)\); at the two real points \(z=\pm1\), \(G_m(z)\) is real Gaussian of variance \(m+1\). Consequently, for \(t\in(0,2]\),
\[
  \Prob\{\abs{G_m(z)}\le t\}\;\le\;\frac{C_\gamma\,t^2}{m+1}\quad(\operatorname{dist}_\T(z,\{\pm1\})\ge\gamma/m),
  \qquad
  \Prob\{\abs{G_m(z)}\le t\}\;\le\;\frac{C\,t}{\sqrt{m+1}}\quad(z=\pm1).
\]
\end{lemma}

\begin{proof}
$\Re G_m(z)=\sum_j\xi_j\cos(j\theta)$ and $\Im G_m(z)=\sum_j\xi_j\sin(j\theta)$ are jointly Gaussian. Its covariance matrix is
\[
  \frac12
  \begin{pmatrix}
    m+1+\Re S_m(\theta) & \Im S_m(\theta)\\
    \Im S_m(\theta) & m+1-\Re S_m(\theta)
  \end{pmatrix},
  \qquad
  S_m(\theta):=\sum_{j=0}^m e^{2ij\theta}.
\]
Thus the eigenvalues are \(\frac12(m+1\pm |S_m(\theta)|)\). Write
\(u=\operatorname{dist}(2\theta,2\pi\Z)\). The separation from \(0\) and
\(\pi\) gives \(u\ge 2\gamma/m\), and
\[
  |S_m(\theta)|
  =\left|\frac{\sin((m+1)u/2)}{\sin(u/2)}\right|.
\]
For fixed \(\gamma\), the ratio \(|S_m(\theta)|/(m+1)\) is bounded by some
\(q_\gamma<1\) uniformly for \(m\ge m_0(\gamma)\): on \(u\le u_0\) this follows
from the uniform bound \(\sup_{x\ge\gamma}|\sin x|/x<1\), after putting
\(x=(m+1)u/2\) and choosing \(u_0\) so that \((u/2)/\sin(u/2)\) is close to
\(1\), while on \(u\ge u_0\) the ratio is \(O_{u_0}(1/m)\). Hence the
smallest covariance eigenvalue is at least \(c_\gamma(m+1)\), so the planar
Gaussian density is bounded by \(C_\gamma/(m+1)\). Integrating over a disk of
radius \(t\le2\) gives the first estimate. At \(z=\pm1\) the imaginary part
vanishes identically and \(G_m(\pm1)\) is a single real Gaussian of variance
\(m+1\), whose density is \(\le C/\sqrt{m+1}\), giving the second estimate.
\end{proof}

The two real points are exactly the parity points $z=\pm1$ of \cite{Parity}: at those points the field is one-real-dimensional, the small-ball exponent drops from $t^2$ to $t^1$, and the anti-concentration is weaker by a factor $\sqrt{m+1}$. This is the first appearance, inside the small-ball calculation, of the boundary mechanism developed in Section~\ref{sec:bridge}, and it is exactly the two endpoints that Theorem~\ref{thm:singlearc} separates off through the $\gamma/m$ condition.

\begin{proposition}[Bounded-density fixed points and deterministic grids]\label{prop:nongaussian-grid}
Let \(\xi_0,\dots,\xi_m\) be independent real random variables whose densities
are bounded by \(K\).  Fix \(\gamma>0\).  There are constants
\(C_{\gamma,K},C_K<\infty\) and \(m_0(\gamma)\) such that, for \(m\ge m_0(\gamma)\), all
\(a\in\C\), and all \(t>0\),
\[
  \Prob\{\abs{G_m(z)-a}\le t\}
  \le \frac{C_{\gamma,K}t^2}{m+1}
  \qquad
  \bigl(\dist_\T(z,\{\pm1\})\ge\gamma/m\bigr),
\]
while at the two real points
\[
  \Prob\{\abs{G_m(\pm1)-a}\le t\}
  \le \frac{C_Kt}{\sqrt{m+1}}.
\]
Consequently, for every deterministic finite set \(Q\subset\T\),
\[
  \Prob\Bigl\{\min_{z\in Q_*}\abs{G_m(z)}\le t\Bigr\}
  \le \frac{C_{\gamma,K}\abs{Q_*}t^2}{m+1},
  \qquad
  Q_*:=\{z\in Q:\dist_\T(z,\{\pm1\})\ge\gamma/m\},
\]
and the two endpoint evaluations contribute at most \(2C_Kt/\sqrt{m+1}\).
\end{proposition}

\begin{proof}
For \(z=e^{i\theta}\), let
\[
  A_\theta:\R^{m+1}\to\R^2,\qquad
  A_\theta x=\left(\sum_{j=0}^m x_j\cos(j\theta),
                  \sum_{j=0}^m x_j\sin(j\theta)\right).
\]
The deterministic row Gram \(A_\theta A_\theta^\top\) is the same matrix
computed as the Gaussian covariance in Lemma~\ref{lem:onepoint}.  Hence, if
\(\dist_\T(z,\{\pm1\})\ge\gamma/m\), both singular values of \(A_\theta\) are at
least \(c_\gamma\sqrt{m+1}\) for all large \(m\).  The linear-image small-ball
estimate of Rudelson--Vershynin \cite{RudelsonVershynin2015}, applied to the
two-dimensional image \(A_\theta\xi\), gives uniformly in \(y\in\R^2\)
\[
  \Prob\{A_\theta\xi\in B_{\R^2}(y,t)\}
  \le
  \frac{C_K\,\operatorname{vol}(B_{\R^2}(0,t))}
       {s_1(A_\theta)s_2(A_\theta)}
  \le \frac{C_{\gamma,K}t^2}{m+1}.
\]
Taking \(y=(\Re a,\Im a)\) gives the first displayed estimate, since
\(A_\theta\xi=(\Re G_m(z),\Im G_m(z))\).

At \(z=1\) and \(z=-1\), \(G_m(z)=\sum_{j=0}^m a_j\xi_j\) with
\(\abs{a_j}=1\), so \(\norm{a}_2=\sqrt{m+1}\).  The one-dimensional
linear-image small-ball estimate gives
\[
  \sup_{u\in\R}\Prob\{\abs{\sum_j a_j\xi_j-u}\le t\}
  \le C_K\,\frac{t}{\sqrt{m+1}},
\]
and the event \(\{\abs{G_m(\pm1)-a}\le t\}\) is contained in
\(\{\abs{G_m(\pm1)-\Re a}\le t\}\).  This proves the endpoint bound.  The grid
statement is the union bound over the deterministic points.
\end{proof}

\begin{corollary}[Subquadratic deterministic grids miss the continuum scale]
\label{cor:subquadratic-grid-misses}
Let the coefficients have densities bounded by \(K\), and fix
\(\gamma>0\).  Let \(N=m+1\), and let \(Q_m\subset\T\) be deterministic nonempty finite
sets such that every point of \(Q_m\setminus\{\pm1\}\) has angular distance at
least \(\gamma/m\) from \(\{\pm1\}\).  For every fixed \(T<\infty\) there are
constants \(C_{\gamma,K,T}<\infty\) and \(m_0(\gamma)\) such that, for
\(m\ge m_0(\gamma)\),
\[
  \sup_{0\le \tau\le T}
  \Prob\Bigl\{\min_{z\in Q_m}\abs{G_m(z)}
      \le \frac{\tau}{\sqrt N}\Bigr\}
  \le C_{\gamma,K,T}\left(\frac{|Q_m|}{N^2}+\frac1N\right).
\]
In particular, if \(|Q_m|=o(N^2)\), then this probability tends to zero
uniformly for \(0\le\tau\le T\).
\end{corollary}

\begin{proof}
Apply Proposition~\ref{prop:nongaussian-grid} with
\(t=\tau/\sqrt N\).  The non-real points contribute at most
\[
  \frac{C_{\gamma,K}|Q_m|t^2}{N}
  =\frac{C_{\gamma,K}|Q_m|\tau^2}{N^2},
\]
and the two endpoint evaluations contribute at most
\[
  \frac{2C_Kt}{\sqrt N}=\frac{2C_K\tau}{N}.
\]
Taking the supremum over \(0\le\tau\le T\) gives the stated bound.
\end{proof}

\begin{corollary}[Independent randomized subquadratic grids also miss]
\label{cor:random-subquadratic-grid-misses}
Let the coefficients have densities bounded by \(K\), and fix
\(\gamma>0\).  Let \(N=m+1\), and let \(\mathcal Q_m\subset\T\) be a nonempty
finite random set, independent of the coefficients, such that every point of
\(\mathcal Q_m\setminus\{\pm1\}\) has angular distance at least
\(\gamma/m\) from \(\{\pm1\}\) almost surely.  For every fixed \(T<\infty\)
there are constants \(C_{\gamma,K,T}<\infty\) and \(m_0(\gamma)\) such that,
for \(m\ge m_0(\gamma)\),
\[
  \sup_{0\le \tau\le T}
  \Prob\Bigl\{\min_{z\in \mathcal Q_m}\abs{G_m(z)}
      \le \frac{\tau}{\sqrt N}\Bigr\}
  \le C_{\gamma,K,T}\left(\frac{\Exp|\mathcal Q_m|}{N^2}+\frac1N\right).
\]
In particular, if \(\Exp|\mathcal Q_m|=o(N^2)\), then this probability tends
to zero uniformly for \(0\le\tau\le T\).
\end{corollary}

\begin{proof}
Condition on the realized grid \(\mathcal Q_m=Q\).  The deterministic
Corollary~\ref{cor:subquadratic-grid-misses} gives
\[
  \sup_{0\le\tau\le T}
  \Prob\left\{\left.
  \min_{z\in Q}\abs{G_m(z)}\le \frac{\tau}{\sqrt N}
  \ \right|\ \mathcal Q_m=Q\right\}
  \le C_{\gamma,K,T}\left(\frac{|Q|}{N^2}+\frac1N\right).
\]
Integrating over \(Q\) gives the displayed estimate.
\end{proof}

\begin{lemma}[Quadratic grids detect under a speed cutoff]
\label{lem:grid-speed-detection}
Let \(Q\subset\T\) be an \(h\)-net in arclength: for every \(z\in\T\) there is
\(q\in Q\) with \(\dist_\T(z,q)\le h\).  For every coefficient vector, every
\(s>0\), and every \(R>0\),
\[
  \Bigl\{\min_{z\in\T}\abs{G_m(z)}\le s\Bigr\}
  \cap
  \Bigl\{\sup_{z\in\T}\abs{G_m'(z)}\le R\Bigr\}
  \subseteq
  \Bigl\{\min_{q\in Q}\abs{G_m(q)}\le s+Rh\Bigr\}.
\]
Consequently, if \(N=m+1\), \(L\ge1\), \(\tau>0\), and
\[
  h\le \frac{\tau}{2Lm^{3/2}\sqrt N},
\]
then
\[
  \Bigl\{\min_{z\in\T}\abs{G_m(z)}\le \frac{\tau}{2\sqrt N}\Bigr\}
  \cap
  \Bigl\{\sup_{z\in\T}\abs{G_m'(z)}\le Lm^{3/2}\Bigr\}
  \subseteq
  \Bigl\{\min_{q\in Q}\abs{G_m(q)}\le \frac{\tau}{\sqrt N}\Bigr\}.
\]
A uniform grid of \(M\ge C Lm^{3/2}\sqrt N/\tau\) points has this property.
For fixed \(L\) and fixed positive \(\tau\), this is \(O_L(N^2)\) points.
\end{lemma}

\begin{proof}
Choose \(z_0\in\T\) with \(\abs{G_m(z_0)}\le s\), and choose
\(q\in Q\) with \(\dist_\T(z_0,q)\le h\).  The fundamental theorem of calculus
along the shorter arc from \(z_0\) to \(q\) gives
\[
  \abs{G_m(q)}
  \le \abs{G_m(z_0)}+\abs{G_m(q)-G_m(z_0)}
  \le s+h\sup_{\T}\abs{G_m'}.
\]
This proves the first inclusion, and the second follows by putting
\(s=\tau/(2\sqrt N)\) and \(R=Lm^{3/2}\).  The assertion about uniform grids is
the elementary mesh bound for \(M\) equally spaced points on a circle of
arclength \(2\pi\), with the absolute constant \(C\) absorbing \(2\pi\).
\end{proof}

\begin{lemma}[Uniform random points form a net]\label{lem:uniform-random-net}
Let \(U_1,\ldots,U_M\) be independent Haar-uniform points of \(\T\), and put
\(\mathcal U_M=\{U_1,\ldots,U_M\}\).  There are absolute constants
\(C,c>0\) such that, for every \(0<h\le1\),
\[
  \Prob\{\mathcal U_M\text{ is not an }h\text{-net in arclength}\}
  \le \frac{C}{h}\exp(-cMh).
\]
\end{lemma}

\begin{proof}
Choose a deterministic \(h/2\)-net \(R_h\subset\T\) with
\(|R_h|\le C/h\).  If every arclength ball \(B(r,h/2)\), \(r\in R_h\),
contains at least one of the \(U_j\), then \(\mathcal U_M\) is an \(h\)-net.
Each such ball has normalized arclength measure at least \(c h\).  Hence, by
the union bound,
\[
  \Prob\{\exists r\in R_h:\mathcal U_M\cap B(r,h/2)=\emptyset\}
  \le \frac{C}{h}(1-ch)^M
  \le \frac{C}{h}\exp(-cMh).
\]
\end{proof}

\begin{corollary}[Independent uniform random grids detect with one extra logarithm]
\label{cor:uniform-random-grid-detects}
Let \(\mathcal U_m\) consist of \(M_m\) independent Haar-uniform points of
\(\T\), sampled independently of the coefficients of \(G_m\), and put
\(N=m+1\).  Fix \(L\ge1\) and \(\tau>0\).  Then
\[
\begin{aligned}
  E_{m,L,\tau}:={}&
  \Bigl\{\min_{z\in\T}\abs{G_m(z)}\le \frac{\tau}{2\sqrt N}\Bigr\}
  \cap
  \Bigl\{\sup_{z\in\T}\abs{G_m'(z)}\le Lm^{3/2}\Bigr\}.
\end{aligned}
\]
With absolute constants \(C,c>0\),
\[
  \Prob\left(
    E_{m,L,\tau}\setminus
    \Bigl\{\min_{q\in \mathcal U_m}\abs{G_m(q)}
        \le \frac{\tau}{\sqrt N}\Bigr\}
  \right)
  \le
  C\,\frac{Lm^{3/2}\sqrt N}{\tau}
  \exp\left\{-c\,\frac{\tau M_m}{Lm^{3/2}\sqrt N}\right\}.
\]
Consequently, for fixed \(L\) and fixed \(\tau>0\), independent uniform random
grids of size \(M_m\ge C_{A,L,\tau}N^2\log N\) detect continuum
\(N^{-1/2}\)-scale excursions under the speed cutoff with miss probability
\(O_A(N^{-A})\).

More generally, assume the polynomial speed-tail hypothesis in
Corollary~\ref{cor:speedtail-grid-detects} for a given \(A>0\).  Then
independent uniform random grids with
\[
  M_m\ge C_{A,\tau,L_A}N^2(\log N)^{3/2}
\]
satisfy
\[
  \Prob\left(
    \Bigl\{\min_{z\in\T}\abs{G_m(z)}\le \frac{\tau}{2\sqrt N}\Bigr\}
    \setminus
    \Bigl\{\min_{q\in \mathcal U_m}\abs{G_m(q)}
        \le \frac{\tau}{\sqrt N}\Bigr\}
  \right)
  \le C_{A,\tau,L_A}m^{-A}.
\]
\end{corollary}

\begin{proof}
Set
\[
  h=\frac{\tau}{2Lm^{3/2}\sqrt N}.
\]
On the event that \(\mathcal U_m\) is an \(h\)-net,
Lemma~\ref{lem:grid-speed-detection} shows that the displayed miss event under
the speed cutoff is empty.  The first bound is therefore
Lemma~\ref{lem:uniform-random-net}.  If \(L\) and \(\tau\) are fixed and
\(M_m\ge C_{A,L,\tau}N^2\log N\), then \(M_mh\ge C_A\log N\) while
\(h^{-1}=O_{L,\tau}(N^2)\), so the net-failure probability is \(O_A(N^{-A})\).

For the speed-tail version, put \(L=L_A\sqrt{\log m}\).  The same argument
gives an \(O_A(m^{-A})\) net-failure probability once
\(M_m\ge C_{A,\tau,L_A}m^{3/2}\sqrt N(\log m)^{3/2}\), which is the displayed
condition after changing constants.  Add the assumed \(O_A(m^{-A})\)
probability of the complementary speed event.
\end{proof}

\begin{corollary}[Stratified random grids keep the deterministic scale]
\label{cor:stratified-grid-detects}
Let \(I_{1,m},\ldots,I_{M_m,m}\) be the equal-arclength partition of \(\T\).
Choose one point \(q_{j,m}\in I_{j,m}\) for each \(j\), deterministically or
randomly, and put
\[
  \mathcal S_m:=\{q_{1,m},\ldots,q_{M_m,m}\}.
\]
Let \(N=m+1\), fix \(L\ge1\) and \(\tau>0\), and let \(E_{m,L,\tau}\) be the
event from Corollary~\ref{cor:uniform-random-grid-detects}.  If
\[
  M_m\ge C\,\frac{Lm^{3/2}\sqrt N}{\tau},
\]
then, for every realization of the stratified grid,
\[
  E_{m,L,\tau}
  \subseteq
  \Bigl\{\min_{q\in\mathcal S_m}\abs{G_m(q)}
      \le \frac{\tau}{\sqrt N}\Bigr\}.
\]
Consequently, under the polynomial speed-tail hypothesis of
Corollary~\ref{cor:speedtail-grid-detects}, if
\[
  M_m\ge C_\tau L_A m^{3/2}\sqrt N\,\sqrt{\log m},
\]
then
\[
  \Prob\left(
    \Bigl\{\min_{z\in\T}\abs{G_m(z)}\le \frac{\tau}{2\sqrt N}\Bigr\}
    \setminus
    \Bigl\{\min_{q\in \mathcal S_m}\abs{G_m(q)}
        \le \frac{\tau}{\sqrt N}\Bigr\}
  \right)
  \le C_A m^{-A}.
\]
Thus stratified one-point-per-cell sampling has the same
\(O_{A,\tau}(L_A N^2\sqrt{\log N})\) detection complexity as a deterministic
net and does not pay the coupon-collector logarithm of independent uniform
sampling.  The same conclusion holds for random shifts of a uniform grid.
\end{corollary}

\begin{proof}
The set \(\mathcal S_m\) is a \(2\pi/M_m\)-net in arclength: every point of
\(\T\) lies in one of the intervals \(I_{j,m}\), and that whole interval has
arclength \(2\pi/M_m\).  If \(M_m\) is at least a sufficiently large absolute
constant times \(Lm^{3/2}\sqrt N/\tau\), then
\[
  \frac{2\pi}{M_m}\le \frac{\tau}{2Lm^{3/2}\sqrt N}.
\]
Lemma~\ref{lem:grid-speed-detection} gives the first inclusion.  For the
speed-tail conclusion, put \(L=L_A\sqrt{\log m}\); outside the displayed speed
event the miss probability is at most \(C_A m^{-A}\), and on the event the
first inclusion applies.  A random shift of a uniform grid satisfies the same
mesh bound, so the same proof applies to it directly.
\end{proof}

\begin{remark}[What the grid estimate does not prove]\label{rem:grid-not-continuum}
Proposition~\ref{prop:nongaussian-grid} proves the density-level delocalization
at every fixed boundary evaluation, including deterministic grids.  It does not
prove Hypothesis~\ref{hyp:boundary-law}: the latter is a continuum statement,
and an excursion of \(\abs{G_m}\) below \(t\) can have angular width only
\(\asymp t\,m^{-3/2}\).  Corollary~\ref{cor:subquadratic-grid-misses} makes
the sampling gap quantitative at the continuum scale \(t\asymp m^{-1/2}\):
away from the two parity endpoints, any deterministic grid with \(o(m^2)\)
points has vanishing detection probability, and
Corollary~\ref{cor:random-subquadratic-grid-misses} gives the same conclusion
for independent randomized grids with \(o(m^2)\) expected size.
Lemma~\ref{lem:grid-speed-detection}
gives the matching deterministic sufficiency statement: once the curve speed is
bounded by \(Lm^{3/2}\), a grid with spacing \(O(t/(Lm^{3/2}))\), hence
\(O_L(m^2)\) points at \(t\asymp m^{-1/2}\), detects every continuum excursion
up to a harmless factor of two in the threshold.  For Gaussian coefficients,
Corollary~\ref{cor:gaussian-grid-detects} combines this deterministic
sufficiency with the Salem--Zygmund speed tail to give detection by
\(O(N^2\sqrt{\log N})\) grid points outside a polynomially small exceptional
event.  Corollary~\ref{cor:uniform-random-grid-detects} gives the independent
uniform random-grid analogue, with the expected extra logarithmic covering
overhead: \(O(N^2\log N)\) points under a fixed speed cutoff and
\(O(N^2(\log N)^{3/2})\) points under the natural \(\sqrt{\log N}\) speed tail.
Corollary~\ref{cor:stratified-grid-detects} shows that stratified random grids
and random shifts of a uniform lattice avoid this coupon-collector overhead and
keep the deterministic \(O(N^2\sqrt{\log N})\) natural-speed scale.
Corollary~\ref{cor:natural-speed-sampling-bracket} packages this as the
natural-speed sampling bracket.  Without such a natural speed tail, the
variance-only fallback in
Corollary~\ref{cor:variance-grid-fixed-power} gives fixed powers
\(N^\beta\), \(\beta>5/2\), with miss rate \(O_\beta(N^{5-2\beta})\).
Passing from fixed points to arcs therefore still requires a
curve-crossing, occupation-measure, or derivative-control argument.
For Gaussian coefficients away from \(z=\pm1\) this is supplied sharply by
Theorem~\ref{thm:singlearc}; the endpoint arcs are controlled at polynomial
scales by Corollary~\ref{prop:uniform-polylog}.  For general bounded-density
laws the corresponding continuum step remains the universality gap.
\end{remark}

\begin{lemma}[Salem--Zygmund derivative tail]\label{lem:derivtail}
Let $\xi_j$ be independent standard real Gaussians. There are absolute constants $C_1,c_1>0$ such that
\[
  \Prob\Bigl\{\sup_{z\in\T}\abs{G_m'(z)}\ge \lambda\,m^{3/2}\sqrt{\log m}\Bigr\}\le C_1\,m^{-c_1\lambda^2}\qquad(\lambda\ge2).
\]
In particular, $\sup_\T\abs{G_m'}\le m^{3/2}\log m$ off an event of probability $O(m^{-10})$.
\end{lemma}

\begin{proof}
$G_m'(z)=\sum_{j=1}^m j\xi_j z^{j-1}$ is a trigonometric polynomial of degree $m$ with $\Exp\abs{G_m'(z)}^2=\sum_{j=1}^m j^2\asymp m^3$ at each fixed $z\in\T$, so the standard deviation at a point is $\asymp m^{3/2}$. By Gaussian concentration $\abs{G_m'(z)}$ has Gaussian tails on that scale at each point; the Salem--Zygmund maximal inequality for random trigonometric polynomials (as used in \cite{KMS}) upgrades the pointwise bound to the supremum at the cost of the $\sqrt{\log m}$ factor, giving the stated tail. Taking $\lambda$ a large constant yields the high-probability bound.
\end{proof}

\begin{corollary}[Gaussian quadratic-log grids detect continuum excursions]
\label{cor:gaussian-grid-detects}
Let \(G_m(z)=\sum_{j=0}^m\xi_jz^j\) with independent standard real Gaussian
coefficients, and put \(N=m+1\).  Fix \(A>0\) and \(\tau>0\).  There are
constants \(C_{A,\tau},C_A<\infty\) such that, if \(Q_m\subset\T\) is an
\(h_m\)-net with
\[
  h_m\le \frac{1}{C_{A,\tau}m^{3/2}\sqrt N\,\sqrt{\log m}},
\]
then
\[
  \Prob\left(
    \Bigl\{\min_{z\in\T}\abs{G_m(z)}\le \frac{\tau}{2\sqrt N}\Bigr\}
    \setminus
    \Bigl\{\min_{q\in Q_m}\abs{G_m(q)}\le \frac{\tau}{\sqrt N}\Bigr\}
  \right)
  \le C_A m^{-A}.
\]
Equivalently, a uniform grid with
\[
  M\ge C_{A,\tau}m^{3/2}\sqrt N\,\sqrt{\log m}
\]
points detects every continuum excursion at this scale except on an
\(O_A(m^{-A})\) exceptional event.  For fixed \(\tau\), this is
\(O_A(N^2\sqrt{\log N})\) points.
\end{corollary}

\begin{proof}
Choose \(\lambda=\lambda_A\ge2\) so large that the derivative-tail bound in
Lemma~\ref{lem:derivtail} is at most \(C_A m^{-A}\).  Let
\[
  E_\lambda=\Bigl\{\sup_{\T}\abs{G_m'}\le
     \lambda m^{3/2}\sqrt{\log m}\Bigr\}.
\]
If \(h_m\le \tau/(2\lambda m^{3/2}\sqrt N\sqrt{\log m})\), then
Lemma~\ref{lem:grid-speed-detection}, with \(L=\lambda\sqrt{\log m}\), gives
\[
  \Bigl\{\min_{\T}\abs{G_m}\le \frac{\tau}{2\sqrt N}\Bigr\}
  \cap E_\lambda
  \subseteq
  \Bigl\{\min_{q\in Q_m}\abs{G_m(q)}\le \frac{\tau}{\sqrt N}\Bigr\}.
\]
Thus the miss event in the statement is contained in \(E_\lambda^c\), whose
probability is \(O_A(m^{-A})\).  The uniform-grid formulation is again the
elementary mesh estimate on a circle of arclength \(2\pi\), with constants
absorbing \(\tau^{-1}\) and \(\lambda_A\).
\end{proof}

\begin{lemma}[Kac--Rice crossing count away from the real points]\label{lem:kr-crossing}
Let \(G_m(z)=\sum_{j=0}^m\xi_jz^j\) with \(\xi_j\) independent standard real
Gaussians, and put
\[
  X_\theta=(\Re G_m(e^{i\theta}),\Im G_m(e^{i\theta})),
  \qquad Q(\theta)=|X_\theta|^2 .
\]
Fix \(\gamma>0\).  There are constants \(C_\gamma<\infty\) and
\(m_0(\gamma)<\infty\) such that, whenever \(m\ge m_0(\gamma)\), \(J\) is an
interval whose points have angular distance at least \(\gamma/m\) from \(0\)
and \(\pi\), and \(0<t\le1\),
\[
  \Exp\,\#\{\theta\in J:Q(\theta)=t^2\}
  \le C_\gamma |J|\,t\sqrt m .
\]
\end{lemma}

\begin{proof}
Write \(\dot X_\theta=dX_\theta/d\theta\), and let
\(\Sigma_\theta=\Cov(X_\theta)\).  By the covariance computation in
Lemma~\ref{lem:onepoint},
\[
  \lambda_{\min}(\Sigma_\theta)\ge c_\gamma m,
  \qquad \|\Sigma_\theta\|\le Cm
\]
uniformly on \(J\).  Also
\[
  \|\Cov(\dot X_\theta)\|\le\sum_{j=1}^m j^2\le Cm^3,
  \qquad
  \|\Cov(\dot X_\theta,X_\theta)\|
  \le \|\Cov(\dot X_\theta)\|^{1/2}\|\Sigma_\theta\|^{1/2}\le Cm^2 .
\]
Thus, conditionally on \(X_\theta=x\),
\[
  \dot X_\theta=B_\theta x+Z_\theta,\qquad
  \|B_\theta\|\le C_\gamma m,\qquad
  \|\Cov(Z_\theta)\|\le Cm^3 ,
\]
where \(Z_\theta\) is a centered Gaussian vector.  For \(|x|=t\le1\),
\[
  \Exp\bigl[|2x\cdot\dot X_\theta|\,\bigm|\,X_\theta=x\bigr]
  \le C_\gamma(mt^2+m^{3/2}t)
  \le C_\gamma m^{3/2}t .
\]
The density \(p_\theta\) of \(X_\theta\) satisfies \(p_\theta(x)\le C_\gamma/m\).
The first-moment Kac--Rice formula applied to the smooth scalar process
\(Q(\theta)-t^2\), followed by the coarea formula for \(x\mapsto |x|^2\), gives
\[
\begin{aligned}
  \Exp\,\#\{\theta\in J:Q(\theta)=t^2\}
  &=
  \int_J\int_{|x|=t}
  \frac{\Exp[|2x\cdot\dot X_\theta|\,|\,X_\theta=x]\,p_\theta(x)}
       {2|x|}\,dH(x)\,d\theta  \\
  &\le C_\gamma |J|\,t\sqrt m .
\end{aligned}
\]
\end{proof}

\begin{lemma}[Regression variance away from the parity points]\label{lem:kr-regression-lower}
Fix \(\gamma>0\).  There are constants \(\kappa_\gamma>0\) and
\(m_0(\gamma)<\infty\) such that, for \(m\ge m_0(\gamma)\),
\(\dist(\theta,\pi\mathbb Z)\ge\gamma/m\), and \(u\in S^1\),
\[
 \operatorname{dist}\!\left(
   (j(-u_1\sin j\theta+u_2\cos j\theta))_{j=0}^m,\ E_\theta
 \right)^2
 \ge \kappa_\gamma m^3,
\]
where
\[
a_\theta=(\cos j\theta)_{j=0}^m,\qquad
b_\theta=(\sin j\theta)_{j=0}^m,\qquad
E_\theta=\operatorname{span}\{a_\theta,b_\theta\}.
\]
\end{lemma}

\begin{proof}
Put \(v=(j(-u_1\sin j\theta+u_2\cos j\theta))_{j=0}^m\) and let
\(G_\theta\) be the \(2\times2\) Gram matrix of \((a_\theta,b_\theta)\).
The squared distance is the Schur complement
\[
  \|v\|_2^2
  -
  \begin{pmatrix}\langle v,a_\theta\rangle&\langle v,b_\theta\rangle\end{pmatrix}
  G_\theta^{-1}
  \begin{pmatrix}\langle v,a_\theta\rangle\\ \langle v,b_\theta\rangle\end{pmatrix}.
\]
First consider the endpoint scale \(\theta=\alpha/m\), \(\alpha\in[\gamma,C]\).
After division by \(m^3\), this Schur complement converges uniformly to the
\(L^2[0,1]\)-distance squared of
\[
  x(-u_1\sin \alpha x+u_2\cos \alpha x)
\]
from \(\operatorname{span}\{\cos\alpha x,\sin\alpha x\}\).  The limiting
distance is positive: otherwise analytic linear independence would force
\(x(-u_1\sin\alpha x+u_2\cos\alpha x)\) to be a linear combination of
\(\cos\alpha x\) and \(\sin\alpha x\) on an interval.  Compactness in
\((\alpha,u)\in[\gamma,C]\times S^1\) gives a uniform positive lower bound.
The endpoint scale near \(\pi\), namely \(\theta=\pi\pm\alpha/m\) modulo
\(2\pi\), reduces to the same case after multiplying the \(j\)-th coordinate by
\((-1)^j\), an orthogonal change of coordinates that sends \(E_\theta\) to the
corresponding endpoint space and replaces \(u\) by another unit vector.

It remains to treat the complementary region \(\dist(\theta,\pi\mathbb Z)\ge
C/m\).  Write \(\beta=u_2+i u_1\), so \(v_j=j\Re(\beta e^{ij\theta})\).  With
\(S_k=\sum_{j=0}^m j^k\), summation by parts gives, uniformly in this region,
\[
  \sum_{j=0}^m j^k e^{2ij\theta}=O_k(m^k/\dist(\theta,\pi\mathbb Z))
  =O_k(m^{k+1}/C)\qquad(k=0,1,2).
\]
Consequently
\[
  \|v\|_2^2=\frac12 S_2+O(m^3/C),\qquad
  G_\theta=\frac12 S_0 I_2+O(m/C),
\]
and
\[
  \begin{pmatrix}\langle v,a_\theta\rangle\\ \langle v,b_\theta\rangle\end{pmatrix}
  =
  \frac12 S_1\begin{pmatrix}u_2\\ -u_1\end{pmatrix}
  +O(m^2/C).
\]
Since \(G_\theta^{-1}=2S_0^{-1}I_2+O(1/(Cm))\), the Schur complement is
\[
  \frac12S_2-\frac{S_1^2}{2S_0}+O(m^3/C)
  =
  \frac{m(m+1)(m+2)}{24}+O(m^3/C).
\]
Choosing \(C=C(\gamma)\) large and then \(m_0(\gamma)\) large gives the claimed
lower bound.
\end{proof}

\begin{proposition}[Away-from-endpoint first moment]\label{prop:away-firstmoment}
Let \(G_m(z)=\sum_{j=0}^m\xi_jz^j\) with \(\xi_j\) independent standard real
Gaussians, and put \(Q(\theta)=|G_m(e^{i\theta})|^2\).  Fix \(\gamma>0\).
There are constants \(0<c_\gamma<C_\gamma<\infty\) and
\(m_0(\gamma)<\infty\) such that, whenever \(m\ge m_0(\gamma)\), \(J\) is an
interval whose points satisfy \(\dist(\theta,\pi\mathbb Z)\ge \gamma/m\), and
\(0<t\le1\),
\[
 c_\gamma |J|\,t\sqrt m
 \le
 \Exp\,\#\{\theta\in J: Q(\theta)=t^2,\ Q'(\theta)<0\}
 \le
 C_\gamma |J|\,t\sqrt m .
\]
The same two-sided bound holds for up-crossings and hence for all level
crossings.  In particular, on an arc of length \(c/m\) away from
\(\{\pm1\}\), the expected number of level-\(t\) down-crossings is
\(\asymp_{\gamma,c} t/\sqrt m\).
\end{proposition}

\begin{proof}
The upper bound follows from Lemma~\ref{lem:kr-crossing}.  For the lower bound,
write
\[
a_\theta=(\cos j\theta)_{j=0}^m,\qquad
b_\theta=(\sin j\theta)_{j=0}^m,\qquad
E_\theta=\operatorname{span}\{a_\theta,b_\theta\}.
\]
If \(u=(u_1,u_2)\in S^1\), then
\[
u\cdot \dot X_\theta
=\sum_{j=0}^m \xi_j\,j(-u_1\sin j\theta+u_2\cos j\theta).
\]
By Lemma~\ref{lem:kr-regression-lower}, taking \(u=x/t\), conditionally on
\(X_\theta=x\), with \(|x|=t\), the Gaussian regression of the normal velocity
satisfies
\[
 \operatorname{Var}(Q'(\theta)\mid X_\theta=x)
 =4\,\operatorname{Var}(x\cdot\dot X_\theta\mid X_\theta=x)
 \ge c_\gamma m^3t^2 .
\]
The regression mean is \(O_\gamma(mt^2)\), hence is \(o(m^{3/2}t)\) for
\(0<t\le1\) after increasing \(m_0(\gamma)\).  Therefore
\[
 \Exp\bigl[(-Q'(\theta))_+\,\bigm|\,X_\theta=x\bigr]
 \ge c_\gamma m^{3/2}t .
\]
The covariance bounds in Lemma~\ref{lem:onepoint} imply that the density
\(p_\theta\) of \(X_\theta\) satisfies \(p_\theta(x)\ge c_\gamma/m\) for
\(|x|\le1\).  The Kac--Rice formula for down-crossings and coarea on
\(|x|=t\) now give
\[
\begin{aligned}
 \Exp N_t^\downarrow(J)
 &=\int_J\int_{|x|=t}
   \frac{\Exp[(-Q'(\theta))_+\,|\,X_\theta=x]p_\theta(x)}
        {2t}\,dH(x)\,d\theta  \\
 &\ge c_\gamma\int_J t\sqrt m\,d\theta
 = c_\gamma |J|t\sqrt m .
\end{aligned}
\]
The up-crossing lower bound is identical, and the proposition follows.
\end{proof}

\begin{proposition}[Bulk Kac--Rice crossing constant]\label{prop:bulk-kr-constant}
Let \(G_m(z)=\sum_{j=0}^m\xi_jz^j\) with \(\xi_j\) independent standard real
Gaussians, put \(N=m+1\), and set \(Q(\theta)=|G_m(e^{i\theta})|^2\).  Fix
\(\delta>0\), and let
\[
  B_\delta:=\{\theta\in[-\pi,\pi]:\dist(\theta,\pi\mathbb Z)\ge\delta\}.
\]
If \(J\subset B_\delta\) is an interval and
\[
  N_t^\downarrow(J):=
  \#\{\theta\in J:Q(\theta)=t^2,\ Q'(\theta)<0\},
\]
then, uniformly for \(0<t\le1\) and for all such intervals \(J\),
\begin{equation}\label{eq:bulk-kr-constant}
  \Exp N_t^\downarrow(J)
  =
  (1+o_\delta(1))\,\frac{|J|\,t\sqrt N}{\sqrt{12\pi}}
  \qquad(m\to\infty).
\end{equation}
The same asymptotic holds for up-crossings, so the leading bulk all-crossing
intensity is \(2t\sqrt N/\sqrt{12\pi}\) per unit arclength.
\end{proposition}

\begin{proof}
Write
\[
  X_\theta=(\Re G_m(e^{i\theta}),\Im G_m(e^{i\theta})),\qquad
  \dot X_\theta=\frac{d}{d\theta}X_\theta .
\]
For \(k=0,1,2\), summation by parts gives, uniformly for
\(\theta\in B_\delta\),
\[
  \sum_{j=0}^{m}j^k e^{2ij\theta}=O_{\delta,k}(m^k).
\]
With \(S_1=\sum_{j=0}^m j=N(N-1)/2\) and
\(S_2=\sum_{j=0}^m j^2=N(N-1)(2N-1)/6\), the covariance matrices therefore
satisfy
\[
  \Sigma_\theta:=\Cov(X_\theta)=\frac N2 I_2+O_\delta(1),
\]
\[
  \Gamma_\theta:=\Cov(\dot X_\theta)=\frac{S_2}{2}I_2+O_\delta(N^2),
\]
and
\[
  C_\theta:=\Cov(\dot X_\theta,X_\theta)
  =
  \frac{S_1}{2}
  \begin{pmatrix}0&-1\\ 1&0\end{pmatrix}
  +O_\delta(N).
\]
Consequently, the conditional covariance of \(\dot X_\theta\) given
\(X_\theta\) is
\[
  V_\theta
  :=\Gamma_\theta-C_\theta\Sigma_\theta^{-1}C_\theta^{T}
  =
  \alpha_N I_2+O_\delta(N^2),
  \qquad
  \alpha_N:=\frac{N(N^2-1)}{24}.
\]
Also
\[
  C_\theta\Sigma_\theta^{-1}
  =
  \frac{S_1}{N}
  \begin{pmatrix}0&-1\\ 1&0\end{pmatrix}
  +O_\delta(1).
\]
Thus, for \(x=tu\), \(u\in S^1\), the conditional law of
\(u\cdot\dot X_\theta\) given \(X_\theta=x\) is Gaussian with variance
\(\alpha_N(1+O_\delta(N^{-1}))\) and mean \(O_\delta(t)\), uniformly in
\(\theta,u\), and \(0<t\le1\).  Hence
\[
  \Exp\bigl[(-u\cdot\dot X_\theta)_+\mid X_\theta=tu\bigr]
  =
  (1+o_\delta(1))\sqrt{\frac{\alpha_N}{2\pi}}.
\]
The density of \(X_\theta\) satisfies, uniformly for \(|x|\le1\),
\[
  p_\theta(x)=\frac{1+o_\delta(1)}{\pi N}.
\]

Kac--Rice for down-crossings of \(Q(\theta)-t^2\), followed by coarea on the
circle \(|x|=t\), gives the pointwise down-crossing intensity
\[
\begin{aligned}
  \mathcal I_m(\theta,t)
  &=
  \int_{|x|=t}
  \frac{\Exp[(-Q'(\theta))_+\mid X_\theta=x]p_\theta(x)}{2t}\,dH(x)  \\
  &=
  \int_{|x|=t}
  \Exp[(-u\cdot\dot X_\theta)_+\mid X_\theta=tu]p_\theta(tu)\,dH(x) \\
  &=
  (1+o_\delta(1))\,2\pi t\,
     \sqrt{\frac{\alpha_N}{2\pi}}\,
     \frac1{\pi N}
   =
  (1+o_\delta(1))\frac{t\sqrt N}{\sqrt{12\pi}}.
\end{aligned}
\]
The error is uniform on \(B_\delta\), so integration over \(J\) proves
\eqref{eq:bulk-kr-constant}.  Replacing \(Q'\) by \(-Q'\) gives the identical
up-crossing asymptotic.
\end{proof}

\begin{corollary}[Fixed-bulk continuum-scale crossing limit]
\label{cor:bulk-continuum-intensity}
With the notation of Proposition~\ref{prop:bulk-kr-constant}, fix
\(0<\delta<\pi/2\) and write
\[
  B_\delta=\{\theta\in[-\pi,\pi]:\dist(\theta,\pi\mathbb Z)\ge\delta\}.
\]
For \(0\le\tau<\infty\), set \(t_N=\tau/\sqrt N\).  Then, uniformly for
\(\tau\) in compact subsets of \([0,\infty)\),
\[
  \Exp N_{t_N}^\downarrow(B_\delta)
  \longrightarrow
  \frac{|B_\delta|\,\tau}{\sqrt{12\pi}}
  =
  \frac{(2\pi-4\delta)\tau}{\sqrt{12\pi}},
\]
and the same limit holds for up-crossings.  Hence the fixed-bulk all-crossing
first moment tends to
\[
  \frac{2(2\pi-4\delta)\tau}{\sqrt{12\pi}}.
\]
Letting \(\delta\downarrow0\) after this fixed-bulk limit gives the non-endpoint
candidate constants \(\sqrt{\pi/3}\,\tau\) for down-crossings and
\(2\sqrt{\pi/3}\,\tau\) for all crossings.  This last passage is only a
first-moment benchmark; Corollary~\ref{cor:global-prob-upper} below shows that
the endpoint neighborhoods do not change the corresponding probability upper
constant, but it does not prove the lower-tail or exponential minimum-modulus
law.
\end{corollary}

\begin{proof}
The set \(B_\delta\) is the disjoint union of the two intervals
\([-\pi+\delta,-\delta]\) and \([\delta,\pi-\delta]\), whose total length is
\(2\pi-4\delta\).  Apply Proposition~\ref{prop:bulk-kr-constant} to each
component and put \(t=t_N=\tau/\sqrt N\).  The multiplicative
\(o_\delta(1)\) error in \eqref{eq:bulk-kr-constant} is uniform for
\(0<t\le1\), so for \(\tau\) in a fixed compact interval the absolute error is
\(o_\delta(1)\tau\).  This proves the down-crossing limit, and the up-crossing
limit follows in the same way.  The displayed all-crossing limit is their sum.
\end{proof}

\begin{corollary}[Fixed-bulk continuum-scale probability upper bound]
\label{cor:bulk-prob-upper}
With the notation of Corollary~\ref{cor:bulk-continuum-intensity}, put
\[
  P_{N,\delta}(\tau):=
  \Prob\Bigl\{
    \min_{\theta\in B_\delta}|G_m(e^{i\theta})|
    \le \frac{\tau}{\sqrt N}
  \Bigr\}.
\]
For every \(K<\infty\),
\[
  \limsup_{N\to\infty}\;
  \sup_{0\le\tau\le K}
  \left(
    P_{N,\delta}(\tau)
    -\frac{|B_\delta|\,\tau}{\sqrt{12\pi}}
  \right)_+
  =0 .
\]
Equivalently, for each fixed \(\tau\),
\[
  \limsup_{N\to\infty}
  \Prob\Bigl\{
    \min_{\theta\in B_\delta}|G_m(e^{i\theta})|
    \le \frac{\tau}{\sqrt N}
  \Bigr\}
  \le
  \frac{(2\pi-4\delta)\tau}{\sqrt{12\pi}}.
\]
This is only a first-moment upper bound on a fixed bulk.  It does not supply
the matching lower bound or the exponential minimum-modulus law, which require
anti-clustering and decorrelation of crossings.
\end{corollary}

\begin{proof}
Write \(B_\delta=I_1\cup I_2\), where
\(I_1=[-\pi+\delta,-\delta]\) and \(I_2=[\delta,\pi-\delta]\).  Let
\(t_N=\tau/\sqrt N\).  If
\(\min_{\theta\in I_j}|G_m(e^{i\theta})|\le t_N\), then either one endpoint of
\(I_j\) has modulus at most \(t_N\), or the level
\(|G_m(e^{i\theta})|=t_N\) is crossed inside \(I_j\).  If both endpoints are
above level, the first entry into the sublevel set is a down-crossing of
\(Q(\theta)-t_N^2\).  The standard nondegeneracy used in the Kac--Rice formula
makes tangential level contacts a null event, so the latter contribution is
bounded by the down-crossing count.

The four interval endpoints are a fixed positive distance from \(\{\pm1\}\),
so Lemma~\ref{lem:onepoint} gives a total endpoint contribution
\[
  O_\delta(t_N^2/N)=O_{\delta,K}(N^{-2})
  \qquad(0\le\tau\le K).
\]
By Corollary~\ref{cor:bulk-continuum-intensity}, the expected down-crossing
count in \(B_\delta\) is
\[
  \left(1+o_\delta(1)\right)
  \frac{|B_\delta|\,\tau}{\sqrt{12\pi}},
\]
uniformly for \(0\le\tau\le K\).  Markov's inequality gives the displayed
uniform limsup, and the pointwise bound follows by fixing \(\tau\).
\end{proof}

\begin{proposition}[Second-moment transfer to a fixed-bulk lower tail]
\label{prop:bulk-secondmoment-transfer}
With the notation of Corollary~\ref{cor:bulk-continuum-intensity}, fix
\(0<\delta<\pi/2\), \(T<\infty\), and put
\[
  \mu_\delta(\tau):=\frac{|B_\delta|\,\tau}{\sqrt{12\pi}}
  =\frac{(2\pi-4\delta)\tau}{\sqrt{12\pi}},
  \qquad 0\le\tau\le T.
\]
For \(N=m+1\) and \(t_N=\tau/\sqrt N\), set
\[
  X_{N,\delta}(\tau):=N_{t_N}^\downarrow(B_\delta).
\]
Assume that for some constant \(L_{\delta,T}<\infty\),
\begin{equation}\label{eq:bulk-factorial-secondmoment}
  \limsup_{N\to\infty}\,
  \Exp\!\left[
    X_{N,\delta}(\tau)\bigl(X_{N,\delta}(\tau)-1\bigr)
  \right]
  \le L_{\delta,T}\,\mu_\delta(\tau)^2
  \qquad(0\le\tau\le T).
\end{equation}
Then, for every \(0\le\tau\le T\),
\begin{equation}\label{eq:bulk-lower-transfer}
  \liminf_{N\to\infty}
  \Prob\Bigl\{
    \min_{\theta\in B_\delta}|G_m(e^{i\theta})|
    \le \frac{\tau}{\sqrt N}
  \Bigr\}
  \ge
  \frac{\mu_\delta(\tau)}
       {1+L_{\delta,T}\mu_\delta(\tau)}.
\end{equation}
In particular, if the Poisson-scale decorrelation estimate
\[
  \Exp\!\left[
    X_{N,\delta}(\tau)\bigl(X_{N,\delta}(\tau)-1\bigr)
  \right]
  \le (1+o(1))\,\mu_\delta(\tau)^2
\]
holds for fixed \(\delta\) and \(\tau\), then the fixed-bulk lower tail has the
small-\(\tau\) expansion
\[
  \liminf_{N\to\infty}
  \Prob\Bigl\{
    \min_{\theta\in B_\delta}|G_m(e^{i\theta})|
    \le \frac{\tau}{\sqrt N}
  \Bigr\}
  \ge
  \mu_\delta(\tau)-O_\delta(\tau^2)
  \qquad(\tau\downarrow0).
\]
\end{proposition}

\begin{proof}
The case \(\tau=0\) is trivial, so assume \(\tau>0\).  By
Corollary~\ref{cor:bulk-continuum-intensity},
\[
  \Exp X_{N,\delta}(\tau)\longrightarrow \mu_\delta(\tau).
\]
The event \(\{X_{N,\delta}(\tau)>0\}\) is contained in
\[
  \left\{
    \min_{\theta\in B_\delta}|G_m(e^{i\theta})|\le \frac{\tau}{\sqrt N}
  \right\}.
\]
Paley--Zygmund for the nonnegative integer-valued variable
\(X_{N,\delta}(\tau)\) gives
\[
  \Prob\{X_{N,\delta}(\tau)>0\}
  \ge
  \frac{(\Exp X_{N,\delta}(\tau))^2}
       {\Exp X_{N,\delta}(\tau)^2}.
\]
Since
\[
  \Exp X_{N,\delta}(\tau)^2
  =
  \Exp X_{N,\delta}(\tau)
  +
  \Exp\!\left[
    X_{N,\delta}(\tau)(X_{N,\delta}(\tau)-1)
  \right],
\]
the convergence of the first moment and the assumed factorial second-moment
bound imply
\[
  \liminf_{N\to\infty}\Prob\{X_{N,\delta}(\tau)>0\}
  \ge
  \frac{\mu_\delta(\tau)^2}
       {\mu_\delta(\tau)+L_{\delta,T}\mu_\delta(\tau)^2},
\]
which is \eqref{eq:bulk-lower-transfer}.  Under the Poisson-scale bound, put
\(L_{\delta,T}=1+o(1)\) at the fixed value of \(\tau\) to get
\(\mu_\delta(\tau)/(1+\mu_\delta(\tau))\), which is
\(\mu_\delta(\tau)-O_\delta(\tau^2)\) as \(\tau\downarrow0\).
\end{proof}

\begin{corollary}[Fixed-bulk first-order law from a factorial moment bound]
\label{cor:bulk-firstorder-sandwich}
With the notation and hypotheses of
Proposition~\ref{prop:bulk-secondmoment-transfer}, and with
\(P_{N,\delta}(\tau)\) as in Corollary~\ref{cor:bulk-prob-upper}, one has, for
every \(0<\tau\le T\),
\[
  \frac{\mu_\delta(\tau)}
       {1+L_{\delta,T}\mu_\delta(\tau)}
  \le
  \liminf_{N\to\infty} P_{N,\delta}(\tau)
  \le
  \limsup_{N\to\infty} P_{N,\delta}(\tau)
  \le
  \mu_\delta(\tau).
\]
Consequently,
\[
  \lim_{\tau\downarrow0}
  \frac{1}{\tau}\liminf_{N\to\infty} P_{N,\delta}(\tau)
  =
  \lim_{\tau\downarrow0}
  \frac{1}{\tau}\limsup_{N\to\infty} P_{N,\delta}(\tau)
  =
  \frac{|B_\delta|}{\sqrt{12\pi}}.
\]
\end{corollary}

\begin{proof}
The lower bound is \eqref{eq:bulk-lower-transfer}, and the upper bound is
Corollary~\ref{cor:bulk-prob-upper}.  Since
\(\mu_\delta(\tau)=|B_\delta|\tau/\sqrt{12\pi}\), dividing the sandwich by
\(\tau\) and letting \(\tau\downarrow0\) gives the displayed first-order
coefficient.
\end{proof}

\begin{corollary}[Uniform bulk second moment implies a global lower tail]
\label{cor:global-secondmoment-transfer}
Fix \(T<\infty\).  Suppose there is \(L_T<\infty\) such that, for every
\(0<\delta<\pi/2\) and \(0\le\tau\le T\), the variables
\[
  X_{N,\delta}(\tau)=N_{\tau/\sqrt N}^\downarrow(B_\delta)
\]
satisfy
\[
  \limsup_{N\to\infty}
  \Exp\!\left[
    X_{N,\delta}(\tau)\bigl(X_{N,\delta}(\tau)-1\bigr)
  \right]
  \le L_T\,\mu_\delta(\tau)^2,
  \qquad
  \mu_\delta(\tau)=\frac{(2\pi-4\delta)\tau}{\sqrt{12\pi}}.
\]
Then, for every \(0\le\tau\le T\),
\begin{equation}\label{eq:global-lower-transfer}
  \liminf_{N\to\infty}
  \Prob\Bigl\{
    \min_{\theta\in[-\pi,\pi]}|G_m(e^{i\theta})|
    \le \frac{\tau}{\sqrt N}
  \Bigr\}
  \ge
  \frac{\sqrt{\pi/3}\,\tau}
       {1+L_T\sqrt{\pi/3}\,\tau}.
\end{equation}
In particular, a uniform Poisson-scale factorial second-moment estimate
\(L_T=1\) would give the small-\(\tau\) global lower-tail expansion
\[
  \liminf_{N\to\infty}
  \Prob\Bigl\{
    \min_{\T}|G_m|
    \le \frac{\tau}{\sqrt N}
  \Bigr\}
  \ge
  \sqrt{\frac{\pi}{3}}\,\tau-\frac{\pi}{3}\tau^2+O(\tau^3)
  \qquad(\tau\downarrow0).
\]
\end{corollary}

\begin{proof}
For fixed \(\delta\), Proposition~\ref{prop:bulk-secondmoment-transfer} gives
\[
  \liminf_{N\to\infty}
  \Prob\Bigl\{
    \min_{\theta\in B_\delta}|G_m(e^{i\theta})|
    \le \frac{\tau}{\sqrt N}
  \Bigr\}
  \ge
  \frac{\mu_\delta(\tau)}
       {1+L_T\mu_\delta(\tau)}.
\]
The event on \(B_\delta\) is contained in the corresponding whole-circle event.
Letting \(\delta\downarrow0\) and using
\[
  \lim_{\delta\downarrow0}\mu_\delta(\tau)
  =
  \frac{2\pi\tau}{\sqrt{12\pi}}
  =
  \sqrt{\frac{\pi}{3}}\,\tau
\]
proves \eqref{eq:global-lower-transfer}.  The final expansion is the Taylor
series of \(x/(1+x)\) at \(x=\sqrt{\pi/3}\,\tau\).
\end{proof}

\begin{proof}[Proof of Theorem~\ref{thm:singlearc}]
Parametrize \(I\) by an interval \(J=[a,b]\), and set
\(Q(\theta)=|G_m(e^{i\theta})|^2\).  If \(\min_JQ\le t^2\), then either
\(Q(a)\le t^2\), or \(Q(b)\le t^2\), or \(Q(\theta)=t^2\) for some
\(\theta\in(a,b)\).  Therefore Lemmas~\ref{lem:onepoint} and
\ref{lem:kr-crossing} give
\[
\begin{aligned}
  \Prob\{\min_JQ\le t^2\}
  &\le \Prob\{Q(a)\le t^2\}+\Prob\{Q(b)\le t^2\}
      +\Exp\,\#\{\theta\in J:Q(\theta)=t^2\}  \\
  &\le C_\gamma\frac{t^2}{m}+C_\gamma\frac{c}{m}t\sqrt m
   \le C_\gamma(1+c)\frac{t}{\sqrt m},
\end{aligned}
\]
since \(0<t\le1\).  This is \eqref{eq:singlearc}.
\end{proof}

\begin{proposition}[Log-loss single-arc boundary small ball]\label{prop:singlearc-log}
Let $G_m(z)=\sum_{j=0}^m\xi_jz^j$ with $\xi_j$ independent standard real Gaussians. For every fixed $\gamma>0$ and \(A>0\) there are constants \(C_{\gamma,A}<\infty\) and \(m_0(\gamma,A)\) such that, for every \(m\ge m_0(\gamma,A)\), every \(c\ge1\), every arc \(I\subset\T\) of arclength \(\abs I=c/m\) whose angular distance from \(z=\pm1\) is at least \(\gamma/m\), and every \(t\in(0,1]\),
\[
  \Prob\Bigl\{\min_{z\in I}\abs{G_m(z)}\le t\Bigr\}
  \le C_{\gamma,A}\left((1+c)\frac{\sqrt{\log m}\,t}{\sqrt m}+m^{-A}\right).
\]
\end{proposition}

\begin{proof}
Parametrize \(I\) by an interval \(J\) of length \(c/m\), and set
\[
  g(\theta):=\abs{G_m(e^{i\theta})},
  \qquad
  \Lambda(s):=\int_J \mathbf1\{g(\theta)\le s\}\,d\theta .
\]
By Fubini and Lemma~\ref{lem:onepoint}, since every point of \(J\) is \(\gamma/m\) away from \(0\) and \(\pi\),
\[
  \Exp\Lambda(s)
  \le \abs J\,\frac{C_\gamma s^2}{m+1}
  \le \frac{C_\gamma c s^2}{m^2}.
\]
Fix \(\lambda=\lambda(A)\ge2\) large enough that the event
\[
  E_\lambda:=\Bigl\{\sup_{z\in\T}\abs{G_m'(z)}
       \le \lambda m^{3/2}\sqrt{\log m}\Bigr\}
\]
has complement of probability at most \(C_A m^{-A}\), by Lemma~\ref{lem:derivtail}. On \(E_\lambda\), if \(g(\theta_0)\le t\) for some \(\theta_0\in J\), then
\[
  g(\theta)\le2t
  \qquad
  \text{whenever } \theta\in J,\quad
  \abs{\theta-\theta_0}\le
  \delta:=\frac{t}{\lambda m^{3/2}\sqrt{\log m}},
\]
because \(\bigl|\frac{d}{d\theta}G_m(e^{i\theta})\bigr|=\abs{e^{i\theta}G_m'(e^{i\theta})}\). For \(m\) large, \(c\ge1\) and \(t\le1\) give \(\abs J=c/m\ge\delta\), so the intersection of \(J\) with this \(\delta\)-neighborhood has length at least a constant multiple of \(\delta\). Hence
\[
  \mathbf1\Bigl\{\min_{\theta\in J}g(\theta)\le t\Bigr\}\mathbf1_{E_\lambda}
  \le C\,\frac{\Lambda(2t)}{\delta}.
\]
Taking expectations and using the preceding occupation bound with \(s=2t\) gives
\[
  \Prob\Bigl\{\min_I\abs{G_m}\le t\Bigr\}
  \le C_A m^{-A}
     + C\,\frac{\lambda m^{3/2}\sqrt{\log m}}{t}\,
       \frac{C_\gamma c t^2}{m^2}
  \le C_{\gamma,A}\left(c\,\frac{\sqrt{\log m}\,t}{\sqrt m}+m^{-A}\right).
\]
Replacing \(c\) by \(1+c\) gives the stated form.
\end{proof}

\begin{lemma}[Endpoint Dirichlet deficit]\label{lem:endpoint-dirichlet}
Fix \(C_0>0\), and put
\[
  S_m(\theta):=\sum_{j=0}^m e^{2ij\theta}.
\]
There are constants \(c_{C_0}>0\) and \(m_0(C_0)\) such that, for all
\(m\ge m_0(C_0)\) and \(0<\vartheta:=|\theta|\le C_0/m\),
\[
  (m+1)-|S_m(\theta)|\ge c_{C_0}m^3\vartheta^2.
\]
\end{lemma}

\begin{proof}
The geometric-sum formula gives
\[
  |S_m(\theta)|
  =\left|\frac{\sin((m+1)\theta)}{\sin\theta}\right|.
\]
Let \(x=(m+1)\vartheta\).  Since \(\vartheta\le C_0/m\), the variable \(x\)
stays in a fixed compact interval \((0,2C_0]\) for all large \(m\).  The
function \(x\mapsto |\sin x|/x\) is strictly less than one on every compact
subinterval of \((0,\infty)\), and near zero its Taylor expansion is
\[
  \frac{\sin x}{x}=1-\frac{x^2}{6}+O(x^4).
\]
Also
\[
  \frac{\vartheta}{\sin\vartheta}=1+O(\vartheta^2)
  =1+O\left(\frac{x^2}{m^2}\right)
\]
uniformly in the present range.  Splitting \(x\in(0,2C_0]\) into a small
neighborhood of zero and its compact complement, these two estimates give a
constant \(c'_{C_0}>0\) such that
\[
  1-\frac{|S_m(\theta)|}{m+1}
  =1-\frac{|\sin x|}{x}\frac{\vartheta}{\sin\vartheta}
  \ge c'_{C_0}x^2
\]
for all large \(m\).  Multiplying by \(m+1\) and using
\((m+1)^3\asymp m^3\) proves the claim.
\end{proof}

\begin{lemma}[Endpoint occupation bound for bounded-density laws]\label{lem:endpoint-occ}
Let \(\xi_0,\dots,\xi_m\) be independent real random variables whose densities
are bounded by \(K\).  Fix \(C_0>0\) and \(B>0\).  Let
\[
  E_{C_0}:=\{\theta\in[-\pi,\pi]:\operatorname{dist}(\theta,\{0,\pi,-\pi\})\le C_0/m\}.
\]
For \(m\) sufficiently large and \(m^{-B}\le s\le2\),
\[
  \Exp\int_{E_{C_0}}\mathbf1\{\abs{G_m(e^{i\theta})}\le s\}\,d\theta
  \le C_{B,C_0,K}\,\frac{s^2\log m}{m^2}.
\]
\end{lemma}

\begin{proof}
It is enough to work near \(0\); the neighborhoods of \(\pi\) and \(-\pi\) are identical after replacing \(\xi_j\) by \((-1)^j\xi_j\). Put \(\vartheta=|\theta|\le C_0/m\). The real-part coefficient vector has squared norm \(\sum_j\cos^2(j\theta)\ge c_{C_0}m\): indeed, with \(a=m\theta\in[-C_0,C_0]\), the normalized sums converge uniformly to \(\int_0^1\cos^2(ax)\,dx\), whose minimum on this compact interval is positive. Projection onto the real axis gives the one-dimensional bound
\[
  \Prob\{\abs{G_m(e^{i\theta})}\le s\}\le C_{C_0,K}\frac{s}{\sqrt m},
\]
by the one-dimensional linear-image density estimate.  For \(\vartheta>0\), let
\(A_\theta\) be the two-row evaluation matrix from
Proposition~\ref{prop:nongaussian-grid}.  The eigenvalues of
\(A_\theta A_\theta^\top\) are
\[
  \frac12\bigl(m+1\pm |S_m(\theta)|\bigr),
  \qquad S_m(\theta)=\sum_{j=0}^m e^{2ij\theta}.
\]
Lemma~\ref{lem:endpoint-dirichlet} gives
\[
  (m+1)-|S_m(\theta)|\ge c_{C_0}m^3\vartheta^2
  \qquad(0<\vartheta\le C_0/m),
\]
while the larger eigenvalue is \(\asymp m\). Hence
\(\sqrt{\det(A_\theta A_\theta^\top)}\ge c_{C_0}m^2\vartheta\).  The
two-dimensional linear-image density estimate gives
\[
  \Prob\{\abs{G_m(e^{i\theta})}\le s\}
  \le C_{C_0,K}\frac{s^2}{m^2\vartheta}.
\]
Combining the two estimates and integrating,
\[
\begin{aligned}
 \int_0^{C_0/m}\Prob\{\abs{G_m(e^{i\theta})}\le s\}\,d\theta
 &\le C\int_0^{s/m^{3/2}}\frac{s}{\sqrt m}\,d\theta
   +C\int_{s/m^{3/2}}^{C_0/m}\frac{s^2}{m^2\theta}\,d\theta  \\
 &\le C_{B,C_0,K}\frac{s^2\log m}{m^2},
\end{aligned}
\]
because \(s\le2\) makes \(s/m^{3/2}\le C_0/m\) for all large \(m\), and \(s\ge m^{-B}\) controls the logarithm. Multiplying by the constant number of endpoint neighborhoods gives the claim.
\end{proof}

\begin{proposition}[Endpoint single-arc boundary small ball]\label{prop:endpoint-singlearc}
Fix \(C_0<\infty\), and set
\[
  E_+(C_0)=\{e^{i\theta}:|\theta|\le C_0/m\},\qquad
  E_-(C_0)=\{e^{i\theta}:|\theta-\pi|\le C_0/m\}.
\]
Let \(G_m(z)=\sum_{j=0}^m\xi_jz^j\), where the \(\xi_j\) are independent
standard real Gaussians.  There are constants \(C_{C_0}<\infty\) and
\(m_0(C_0)<\infty\) such that, for \(m\ge m_0(C_0)\) and \(0<t\le1\),
\[
  \Prob\Bigl\{\min_{z\in E_\pm(C_0)} |G_m(z)|\le t\Bigr\}
  \le C_{C_0}\frac{t}{\sqrt m}.
\]
Consequently, any arc of length \(c/m\) which meets one of the two real points
satisfies the same bound with right-hand side \(C_{C_0}(1+c)t/\sqrt m\).
\end{proposition}

\begin{proof}
It is enough to treat \(E_+(C_0)\), since \(E_-(C_0)\) is obtained after
replacing \(\xi_j\) by \((-1)^j\xi_j\).  Put
\[
  F_m(u):=m^{-1/2}G_m(e^{iu/m}),\qquad
  X_m(u):=(\Re F_m(u),\Im F_m(u)),\qquad
  \eps:=t/\sqrt m .
\]
The endpoint value gives
\[
  \Prob\{|F_m(0)|\le\eps\}\le C\eps .
\]
Thus it remains to bound the expected number of crossings
\(\#\{u\in(0,C_0]: |F_m(u)|=\eps\}\).

For \(0<u\le u_0\), with \(u_0>0\) small depending only on \(C_0\), write
\[
  r(u)=m^{-1/2}(\cos(ju/m))_{j=0}^m,\qquad
  s(u)=m^{-1/2}(\sin(ju/m))_{j=0}^m,
\]
\[
  \beta(u)=\frac{\langle r(u),s(u)\rangle}{\|r(u)\|^2},
  \qquad w(u)=s(u)-\beta(u)r(u).
\]
Uniformly for large \(m\),
\[
  c\le\|r(u)\|^2\le C,\qquad
  cu^2\le\|w(u)\|^2\le Cu^2,\qquad
  \|r'(u)\|\le Cu,\qquad
  \|w'(u)-u^{-1}w(u)\|\le Cu,
\]
and \(|\beta(u)|\le Cu\), \(|\beta'(u)|\le C\).  These follow from Taylor
expansion and the Riemann-sum identities
\(m^{-1}\sum_{j=0}^m(j/m)^k=(k+1)^{-1}+O(m^{-1})\).

Let \(y_1=x_1\) and \(y_2=x_2-\beta(u)x_1\).  Since \(r(u)\perp w(u)\), the
density of \(X_m(u)\) satisfies
\[
  p_u(x)\le \frac{C}{u}\exp\!\left(-c\,\frac{y_2^2}{u^2}\right).
\]
The same normal form and Gaussian regression give, for \(|x|=\eps\),
\[
  \Exp\bigl[|x\cdot \dot X_m(u)|\,\bigm|\,X_m(u)=x\bigr]
  \le C\left(u\eps+\frac{y_2^2}{u}\right).
\]
Kac--Rice for \(Q(u)=|X_m(u)|^2\), followed by coarea on the circle
\(|x|=\eps\), therefore gives
\[
\begin{aligned}
  \Exp \#\{0<u\le u_0:Q(u)=\eps^2\}
  &\le C\int_0^{u_0}\int_{|x|=\eps}
     \frac{u\eps+y_2^2/u}{\eps}\,
     \frac1u e^{-c y_2^2/u^2}\,dH(x)\,du .
\end{aligned}
\]
On the circle \(|x|=\eps\),
\[
  \int e^{-c y_2^2/u^2}\,dH(x)\le C\min\{\eps,u\},\qquad
  \int y_2^2 e^{-c y_2^2/u^2}\,dH(x)\le C\min\{\eps^3,u^3\}.
\]
Hence the last display is bounded by
\[
  C\int_0^{u_0}\left(\min\{\eps,u\}
    +\frac{\min\{\eps^3,u^3\}}{u^2\eps}\right)\,du
  \le C\eps .
\]
On \(u\in[u_0,C_0]\), equivalently
\(\theta\in[u_0/m,C_0/m]\), Lemma~\ref{lem:kr-crossing} applies with
\(\gamma=u_0\) and contributes at most \(C_{C_0}\eps\).  Therefore
\[
  \Exp \#\{0<u\le C_0: |F_m(u)|=\eps\}\le C_{C_0}\eps .
\]
If the minimum over \(E_+(C_0)\) is at most \(t\), then either the endpoint
value is at most \(t\), or the level \(t\) is crossed inside the arc.  Markov's
inequality for the crossing count proves the displayed bound.  The final
sentence follows by applying this endpoint estimate on a fixed \(O(1/m)\)
neighborhood of the real point and Theorem~\ref{thm:singlearc} to the remaining
part of the arc.
\end{proof}

\begin{proposition}[Endpoint lower bound]\label{prop:endpoint-lower}
For every \(C_0>0\) there are constants \(c_{C_0},C_{C_0}>0\) and
\(m_0(C_0)<\infty\) such that, for \(m\ge m_0(C_0)\) and \(0<t\le1\),
\[
  c_{C_0}\frac{t}{\sqrt m}
  \le
  \Prob\Bigl\{\min_{z\in E_\pm(C_0)} |G_m(z)|\le t\Bigr\}
  \le
  C_{C_0}\frac{t}{\sqrt m}.
\]
In particular the endpoint estimate of Proposition~\ref{prop:endpoint-singlearc}
has the correct order in both \(t\) and \(m\).
\end{proposition}

\begin{proof}
The upper bound is Proposition~\ref{prop:endpoint-singlearc}.  For the lower
bound, \(1\in E_+(C_0)\) and \(-1\in E_-(C_0)\), while
\[
  G_m(1)=\sum_{j=0}^m\xi_j,\qquad
  G_m(-1)=\sum_{j=0}^m(-1)^j\xi_j
\]
are each real Gaussians with variance \(m+1\).  Thus, for \(0<t\le1\),
\[
  \Prob\{|G_m(\pm1)|\le t\}
  =
  \frac1{\sqrt{2\pi(m+1)}}\int_{-t}^{t}e^{-x^2/(2(m+1))}\,dx
  \ge c\,\frac{t}{\sqrt m},
\]
after decreasing \(c\) and taking \(m\) large enough. Since
\(\{\abs{G_m(\pm1)}\le t\}\subseteq
\{\min_{E_\pm(C_0)}\abs{G_m}\le t\}\), the lower bound follows.
\end{proof}

\begin{corollary}[Exact two-endpoint Gaussian formula and sharp expansion]\label{cor:endpoint-two-point}
Let \(G_m(z)=\sum_{j=0}^m\xi_jz^j\) have independent standard real Gaussian
coefficients. Assume \(m\ge1\), and put \(N=m+1\) and
\[
  \rho_N:=\sum_{j=0}^{N-1}(-1)^j\in\{0,1\}.
\]
For \(t\ge0\), set
\[
  p_N(t):=\frac1{\sqrt{2\pi N}}\int_{-t}^{t}e^{-x^2/(2N)}\,dx
\]
and
\[
  q_N(t):=
  \frac1{2\pi\sqrt{N^2-\rho_N^2}}
  \int_{-t}^{t}\int_{-t}^{t}
  \exp\!\left\{
    -\frac{N(x^2+y^2)-2\rho_Nxy}{2(N^2-\rho_N^2)}
  \right\}\,dx\,dy .
\]
Then the endpoint event is exactly
\begin{equation}\label{eq:endpoint-exact}
  \Prob\Bigl\{\min\bigl(|G_m(1)|,|G_m(-1)|\bigr)\le t\Bigr\}
  =2p_N(t)-q_N(t).
\end{equation}
In particular, if \(N\) is even then the two endpoint values are independent
and \(q_N(t)=p_N(t)^2\).  Uniformly for \(0\le t\le1\),
\[
  \begin{aligned}
  \Prob\Bigl\{\min\bigl(|G_m(1)|,|G_m(-1)|\bigr)\le t\Bigr\}
  &=
  2\sqrt{\frac2\pi}\,\frac{t}{\sqrt N}\\
  &\quad-\frac{2}{\pi}\frac{t^2}{\sqrt{N^2-\rho_N^2}}
  +O\!\left(\frac{t^3}{N^{3/2}}\right)
  \end{aligned}
  \qquad(m\to\infty).
\]
Thus the two parity endpoints contribute a sharp global lower floor of order
\(t/\sqrt m\), with leading constant \(2\sqrt{2/\pi}\) before the non-real arcs
are counted.
\end{corollary}

\begin{proof}
Set
\[
  X=G_m(1),\qquad Y=G_m(-1),\qquad
  \rho_N=\sum_{j=0}^{N-1}(-1)^j\in\{0,1\}.
\]
Then \(X\) and \(Y\) are centered Gaussians with
\[
  \Var X=\Var Y=N,\qquad \Cov(X,Y)=\rho_N .
\]
Therefore \(p_N(t)=\Prob\{|X|\le t\}=\Prob\{|Y|\le t\}\), and \(q_N(t)\) is
the bivariate Gaussian probability
\(\Prob\{|X|\le t,\ |Y|\le t\}\). Inclusion--exclusion proves
\eqref{eq:endpoint-exact}.  If \(N\) is even then \(\rho_N=0\), so the centered
Gaussian variables \(X\) and \(Y\) are independent and \(q_N(t)=p_N(t)^2\).

For the asymptotic expansion, one endpoint gives
\[
  p_N(t)
  =\frac1{\sqrt{2\pi N}}\int_{-t}^{t}e^{-x^2/(2N)}\,dx
  =
  \sqrt{\frac2\pi}\,\frac{t}{\sqrt N}
  +O\!\left(\frac{t^3}{N^{3/2}}\right),
\]
uniformly for \(0\le t\le1\).  The joint density of \((X,Y)\) is bounded by
\(C/N\) for all \(m\ge1\), since
\(\det\begin{psmallmatrix}N&\rho_N\\ \rho_N&N\end{psmallmatrix}
=N^2-\rho_N^2\asymp N^2\).  More precisely, on the square
\([-t,t]^2\),
\[
  \exp\!\left\{
    -\frac{N(x^2+y^2)-2\rho_Nxy}{2(N^2-\rho_N^2)}
  \right\}
  =1+O\!\left(\frac{t^2}{N}\right),
\]
uniformly for \(0\le t\le1\).  Therefore
\[
  q_N(t)
  =\frac{2}{\pi}\frac{t^2}{\sqrt{N^2-\rho_N^2}}
  +O\!\left(\frac{t^4}{N^2}\right).
\]
Combining this with \eqref{eq:endpoint-exact} and the one-point expansion gives
the displayed formula, since \(t^4/N^2\le t^3/N^{3/2}\) on \(0\le t\le1\) and
\(N\ge2\).
\end{proof}

\begin{corollary}[Endpoint scale separation at \(t=\tau/\sqrt N\)]
\label{cor:endpoint-continuum-negligible}
Let \(G_m(z)=\sum_{j=0}^m\xi_jz^j\) have independent standard real Gaussian
coefficients, and put \(N=m+1\).  For every \(K<\infty\), uniformly for
\(0\le \tau\le K\),
\[
  \begin{aligned}
  &\Prob\Bigl\{
      \min\bigl(|G_m(1)|,|G_m(-1)|\bigr)\le \frac{\tau}{\sqrt N}
    \Bigr\}  \\
  &\qquad =
  2\sqrt{\frac2\pi}\,\frac{\tau}{N}
  -\frac{2}{\pi}\frac{\tau^2}{N\sqrt{N^2-\rho_N^2}}
  +O_K(N^{-3})
  =
  2\sqrt{\frac2\pi}\,\frac{\tau}{N}+O_K(N^{-2}),
  \end{aligned}
\]
where \(\rho_N=\sum_{j=0}^{N-1}(-1)^j\in\{0,1\}\).  Consequently
\[
  N\,\Prob\Bigl\{
      \min\bigl(|G_m(1)|,|G_m(-1)|\bigr)\le \frac{\tau}{\sqrt N}
    \Bigr\}
  \longrightarrow 2\sqrt{\frac2\pi}\,\tau
\]
uniformly for \(\tau\) in compact subsets of \([0,\infty)\).  Thus the two
parity endpoints have vanishing unscaled probability at the continuum
boundary-minimum threshold \(t=\tau/\sqrt N\); the \(O(1)\) candidate global
law must come from the non-real arcs.
\end{corollary}

\begin{proof}
Substitute \(t=\tau/\sqrt N\) in the expansion of
Corollary~\ref{cor:endpoint-two-point}.  For \(0\le\tau\le K\), the condition
\(t\le1\) holds for all sufficiently large \(N\), and
\[
  \frac{t^2}{\sqrt{N^2-\rho_N^2}}
  =\frac{\tau^2}{N\sqrt{N^2-\rho_N^2}},
  \qquad
  \frac{t^3}{N^{3/2}}=O_K(N^{-3}).
\]
This gives the displayed asymptotic.  Multiplication by \(N\) and
\(\rho_N\in\{0,1\}\) give the uniform compact limit.
\end{proof}

\begin{corollary}[Endpoint neighborhoods are negligible at continuum scale]
\label{cor:endpoint-neighborhood-continuum}
Fix \(C_0<\infty\), and put
\[
  E_{C_0}:=E_+(C_0)\cup E_-(C_0).
\]
Let \(G_m(z)=\sum_{j=0}^m\xi_jz^j\) have independent standard real Gaussian
coefficients, and put \(N=m+1\).  There are constants
\(C_{C_0}<\infty\) and \(m_0(C_0)<\infty\) such that, for
\(m\ge m_0(C_0)\) and \(0<\tau\le \sqrt N\),
\[
  \Prob\Bigl\{
    \min_{z\in E_{C_0}} |G_m(z)|\le \frac{\tau}{\sqrt N}
  \Bigr\}
  \le C_{C_0}\frac{\tau}{N}.
\]
In particular, for every fixed \(\tau\) the natural \(O(1/m)\) neighborhoods
of the two real endpoints have vanishing probability at the continuum
boundary-minimum threshold.  They can contribute only on the \(N^{-1}\)
probability scale, not to an \(O(1)\) whole-circle limiting law.
\end{corollary}

\begin{proof}
Apply Proposition~\ref{prop:endpoint-singlearc} to \(E_+(C_0)\) and
\(E_-(C_0)\), with \(t=\tau/\sqrt N\).  Since \(0<\tau\le\sqrt N\), this
threshold satisfies \(0<t\le1\), and the union bound gives
\[
  \Prob\Bigl\{\min_{z\in E_{C_0}} |G_m(z)|\le \frac{\tau}{\sqrt N}\Bigr\}
  \le 2C_{C_0}\frac{\tau}{\sqrt{mN}}.
\]
Because \(m=N-1\), the last quantity is at most \(C'_{C_0}\tau/N\) after
increasing \(m_0(C_0)\) and the constant.
\end{proof}

\begin{corollary}[Sharp global continuum-scale probability upper bound]
\label{cor:global-prob-upper}
Let \(G_m(z)=\sum_{j=0}^m\xi_jz^j\) have independent standard real Gaussian
coefficients, put \(N=m+1\), and define
\[
  P_N^{\T}(\tau):=
  \Prob\Bigl\{
    \min_{\theta\in[-\pi,\pi]} |G_m(e^{i\theta})|
    \le \frac{\tau}{\sqrt N}
  \Bigr\}.
\]
Then, for every \(K<\infty\),
\[
  \limsup_{N\to\infty}\;
  \sup_{0\le\tau\le K}
  \left(
    P_N^{\T}(\tau)-\sqrt{\frac{\pi}{3}}\,\tau
  \right)_+
  =0 .
\]
Equivalently, for every fixed \(\tau\ge0\),
\[
  \limsup_{N\to\infty} P_N^{\T}(\tau)
  \le \sqrt{\frac{\pi}{3}}\,\tau .
\]
\end{corollary}

\begin{proof}
Fix \(K<\infty\), \(0<\eta<\pi/4\), and a fixed \(C_0\ge4\).  Let
\[
  A_\eta:=\{\theta\in[-\pi,\pi]:\dist(\theta,\pi\mathbb Z)<\eta\}.
\]
The complement of \(A_\eta\) is \(B_\eta\).  Corollary~\ref{cor:bulk-prob-upper}
gives, uniformly for \(0\le\tau\le K\),
\[
  \limsup_{N\to\infty}\;
  \sup_{0\le\tau\le K}
  \left(
    \Prob\Bigl\{
      \min_{\theta\in B_\eta}|G_m(e^{i\theta})|
      \le \frac{\tau}{\sqrt N}
    \Bigr\}
    -\frac{|B_\eta|\,\tau}{\sqrt{12\pi}}
  \right)_+
  =0 .
\]
It remains to bound \(A_\eta\).  Remove from \(A_\eta\) the natural endpoint
neighborhoods \(E_{C_0}\).  For all large \(m\), the remaining set is covered by
at most \(C\eta m\) arcs \(J_\ell\) of arclength \(c_\ell/m\), with
\(1\le c_\ell\le2\), each at angular distance at least \((C_0/2)/m\) from
\(\{\pm1\}\).  Theorem~\ref{thm:singlearc}, with
\(\gamma=C_0/2\) and \(t=\tau/\sqrt N\), gives
\[
  \Prob\Bigl\{
    \min_{e^{i\theta}\in J_\ell}|G_m(e^{i\theta})|
    \le \frac{\tau}{\sqrt N}
  \Bigr\}
  \le C_{C_0,K}\frac{\tau}{N}
  \qquad(0\le\tau\le K),
\]
after increasing \(m\) if necessary.  Summing over the \(O(\eta m)\) arcs gives
\[
  \limsup_{N\to\infty}\sup_{0\le\tau\le K}
  \Prob\Bigl\{
    \min_{\theta\in A_\eta,\ e^{i\theta}\notin E_{C_0}}
      |G_m(e^{i\theta})|
    \le \frac{\tau}{\sqrt N}
  \Bigr\}
  \le C_{C_0}\eta K .
\]
By Corollary~\ref{cor:endpoint-neighborhood-continuum}, the natural endpoint
neighborhoods \(E_{C_0}\) contribute \(o_N(1)\), uniformly for
\(0\le\tau\le K\).  Therefore
\[
  \limsup_{N\to\infty}
  \sup_{0\le\tau\le K}
  \left(
    P_N^{\T}(\tau)
    -\frac{|B_\eta|\,\tau}{\sqrt{12\pi}}
  \right)_+
  \le C_{C_0}\eta K .
\]
Since \(|B_\eta|=2\pi-4\eta\), the coefficient
\(|B_\eta|/\sqrt{12\pi}\) is at most \(2\pi/\sqrt{12\pi}=\sqrt{\pi/3}\).
Thus the same bound holds with
\(\sqrt{\pi/3}\,\tau\) subtracted in place of
\(|B_\eta|\tau/\sqrt{12\pi}\).  Letting \(\eta\downarrow0\) proves the uniform
upper bound, and the pointwise bound follows by fixing \(\tau\).
\end{proof}

\begin{corollary}[Global first-order law from a uniform factorial moment bound]
\label{cor:global-firstorder-sandwich}
Under the hypotheses of Corollary~\ref{cor:global-secondmoment-transfer}, with
\(P_N^{\T}\) as in Corollary~\ref{cor:global-prob-upper}, one has, for every
\(0<\tau\le T\),
\[
  \frac{\sqrt{\pi/3}\,\tau}
       {1+L_T\sqrt{\pi/3}\,\tau}
  \le
  \liminf_{N\to\infty}P_N^{\T}(\tau)
  \le
  \limsup_{N\to\infty}P_N^{\T}(\tau)
  \le
  \sqrt{\frac{\pi}{3}}\,\tau .
\]
Consequently,
\[
  \lim_{\tau\downarrow0}
  \frac{1}{\tau}\liminf_{N\to\infty}P_N^{\T}(\tau)
  =
  \lim_{\tau\downarrow0}
  \frac{1}{\tau}\limsup_{N\to\infty}P_N^{\T}(\tau)
  =
  \sqrt{\frac{\pi}{3}}.
\]
\end{corollary}

\begin{proof}
The lower bound is \eqref{eq:global-lower-transfer}, and the upper bound is
Corollary~\ref{cor:global-prob-upper}.  Divide by \(\tau\) and let
\(\tau\downarrow0\).
\end{proof}

\begin{corollary}[Negative moments below the boundary exponent]
\label{cor:global-negative-moments}
Let
\[
  M_N:=\sqrt N\,\min_{z\in\T}|G_m(z)|,\qquad N=m+1,
\]
for the standard real Gaussian polynomial \(G_m\).  For every \(0<q<1\),
\[
  \sup_{N\ge2}\Exp M_N^{-q}<\infty .
\]
More sharply, for every \(a>0\),
\[
  \limsup_{N\to\infty}
  \Exp\!\left[M_N^{-q}\mathbf1\{M_N\le a\}\right]
  \le
  \frac{\sqrt{\pi/3}}{1-q}\,a^{1-q}.
\]
\end{corollary}

\begin{proof}
First \(M_N>0\) almost surely.  Indeed, a unit-circle zero of \(G_m\) forces
\(G_m\) and the reflected polynomial \(z^mG_m(1/z)\) to have a common root, so
the coefficient vector lies in the zero set of a non-trivial resultant; the
Gaussian coefficient law gives that algebraic hypersurface probability zero.

Put \(F_N(s)=\Prob\{M_N\le s\}\).  Theorem~\ref{thm:uniform}, with
\(t=s/\sqrt N\), gives \(F_N(s)\le C s\) for \(0<s\le1\), after increasing
\(C\) to cover the finitely many small \(N\).  Therefore, for \(0<q<1\),
\[
  \Exp M_N^{-q}
  =
  q\int_0^\infty u^{q-1}\Prob\{M_N^{-1}>u\}\,du
  \le
  1+qC\int_1^\infty u^{q-2}\,du,
\]
uniformly in \(N\).

For the sharp small-window estimate, fix \(0<\rho<a\).  Stieltjes integration
by parts gives
\[
  \Exp\!\left[M_N^{-q}\mathbf1\{M_N\le a\}\right]
  =
  a^{-q}F_N(a)+q\int_0^a s^{-q-1}F_N(s)\,ds .
\]
On \((0,\rho)\) use the uniform bound \(F_N(s)\le Cs\).  On \([\rho,a]\) use
Corollary~\ref{cor:global-prob-upper}, which gives
\[
  \limsup_{N\to\infty}\sup_{\rho\le s\le a}
  \bigl(F_N(s)-\sqrt{\pi/3}\,s\bigr)_+=0 .
\]
Taking the limsup in \(N\) yields
\[
  \limsup_{N\to\infty}
  \Exp\!\left[M_N^{-q}\mathbf1\{M_N\le a\}\right]
  \le
  \sqrt{\frac{\pi}{3}}\,a^{1-q}
  +\frac{qC}{1-q}\rho^{1-q}
  +\frac{q\sqrt{\pi/3}}{1-q}\left(a^{1-q}-\rho^{1-q}\right).
\]
Letting \(\rho\downarrow0\) gives the claimed bound.
\end{proof}

\begin{corollary}[Weak endpoint bound for the inverse scaled minimum]
\label{cor:global-inverse-weak}
With \(M_N\) as in Corollary~\ref{cor:global-negative-moments}, put
\(Y_N=M_N^{-1}\).  Then
\[
  \sup_{N\ge2}\sup_{u\ge1}
  u\,\Prob\{Y_N>u\}<\infty .
\]
Moreover, for every \(1\le R<\infty\),
\[
  \limsup_{N\to\infty}\;
  \sup_{1\le u\le R}
  \left(
    u\,\Prob\{Y_N>u\}-\sqrt{\frac{\pi}{3}}
  \right)_+
  =0 .
\]
Thus the endpoint exponent for \(Y_N\) is weak \(L^1\), and the sharp global
upper constant controls every fixed inverse scale.
\end{corollary}

\begin{proof}
Let \(F_N(s)=\Prob\{M_N\le s\}\).  The proof of
Corollary~\ref{cor:global-negative-moments}, using Theorem~\ref{thm:uniform},
gives a constant \(C\) such that \(F_N(s)\le Cs\) for all \(N\ge2\) and
\(0<s\le1\).  Therefore, for \(u\ge1\),
\[
  \Prob\{Y_N>u\}=F_N(1/u)\le \frac{C}{u},
\]
which proves the uniform weak bound.

For the sharp finite-window constant, Corollary~\ref{cor:global-prob-upper}
with \(K=1\) gives
\[
  \varepsilon_N:=
  \sup_{0\le s\le1}
  \left(F_N(s)-\sqrt{\frac{\pi}{3}}\,s\right)_+
  \longrightarrow0 .
\]
If \(1\le u\le R\), then
\[
  u\,\Prob\{Y_N>u\}
  =uF_N(1/u)
  \le \sqrt{\frac{\pi}{3}}+R\varepsilon_N .
\]
Taking the supremum over \(1\le u\le R\) and then the limsup in \(N\) proves the
claim.
\end{proof}

\begin{corollary}[Quantiles of the inverse scaled minimum]
\label{cor:global-inverse-quantiles}
With \(M_N\) and \(Y_N=M_N^{-1}\) as in
Corollary~\ref{cor:global-inverse-weak}, define, for \(0<\alpha\le1\),
\[
  Q_N(\alpha):=\inf\{u\ge1:\Prob\{Y_N>u\}\le\alpha\}.
\]
Then there is an absolute constant \(C<\infty\) such that
\[
  Q_N(\alpha)\le \frac{C}{\alpha}
  \qquad(N\ge2,\ 0<\alpha\le1).
\]
More sharply, for every fixed \(0<\alpha\le1\),
\[
  \limsup_{N\to\infty} Q_N(\alpha)
  \le \frac{1}{\alpha}\sqrt{\frac{\pi}{3}}.
\]
In lower-quantile form, if
\[
  q_N(p):=\inf\{s\ge0:\Prob\{M_N\le s\}\ge p\},
\]
then there is \(c>0\) such that
\[
  q_N(p)\ge cp
  \qquad(N\ge2,\ 0<p<1),
\]
and, for every fixed \(0<p<1\),
\[
  \liminf_{N\to\infty} q_N(p)\ge p\sqrt{\frac{3}{\pi}}.
\]
Thus target tail probabilities for the inverse minimum translate into linear
clearance thresholds for the scaled minimum, with the same sharp global upper
constant in the limit.
\end{corollary}

\begin{proof}
Enlarge the constant in the uniform weak bound of
Corollary~\ref{cor:global-inverse-weak} so that \(C\ge1\).  Then
\[
  \Prob\{Y_N>C/\alpha\}\le\alpha,
\]
and hence \(Q_N(\alpha)\le C/\alpha\).

For the sharp limsup, put \(K_0=\sqrt{\pi/3}\), fix \(\eps>0\), and set
\(u_\eps=(K_0+\eps)/\alpha\).  The finite-window part of
Corollary~\ref{cor:global-inverse-weak}, applied with \(R=u_\eps\), gives
\[
  \limsup_{N\to\infty}\Prob\{Y_N>u_\eps\}
  \le \frac{K_0}{u_\eps}
  <\alpha .
\]
Thus \(Q_N(\alpha)\le u_\eps\) for all sufficiently large \(N\), and
\(\eps\downarrow0\) gives the displayed bound.

For the lower quantiles of \(M_N\), let \(F_N(s)=\Prob\{M_N\le s\}\).  The
proof of Corollary~\ref{cor:global-inverse-weak} gives \(F_N(s)\le C_0s\) for
all \(N\ge2\) and \(0<s\le1\).  Choosing \(c\le(2C_0)^{-1}\) gives
\(F_N(cp)<p\), hence \(q_N(p)\ge cp\).  For the sharp asymptotic statement,
fix \(a<p\sqrt{3/\pi}\).  Corollary~\ref{cor:global-prob-upper} gives
\[
  \limsup_{N\to\infty}F_N(a)
  \le \sqrt{\frac{\pi}{3}}\,a
  <p .
\]
Therefore \(q_N(p)\ge a\) for all sufficiently large \(N\).  Letting
\(a\uparrow p\sqrt{3/\pi}\) proves the last claim.
\end{proof}

\begin{corollary}[Finite-family weak endpoint bound]
\label{cor:global-inverse-family-weak}
Fix \(L\ge1\).  For each \(N\), let
\(Y_{N,1},\ldots,Y_{N,L}\) be positive random variables whose marginals have the
same law as \(Y_N\) in Corollary~\ref{cor:global-inverse-weak}; no independence
between them is required.  Put
\[
  Y_{N,L}^*:=\max_{1\le\ell\le L}Y_{N,\ell}.
\]
Then
\[
  \sup_{N\ge2}\sup_{u\ge1}
  u\,\Prob\{Y_{N,L}^*>u\}\le C L
\]
with an absolute constant \(C<\infty\).  More sharply, for every fixed
\(L\ge1\) and \(1\le R<\infty\),
\[
  \limsup_{N\to\infty}\;
  \sup_{1\le u\le R}
  \left(
    u\,\Prob\{Y_{N,L}^*>u\}
    -L\sqrt{\frac{\pi}{3}}
  \right)_+
  =0 .
\]
Thus finite simultaneous testing preserves the weak \(L^1\) endpoint, paying
only the linear union-bound factor in the sharp fixed-window constant.
\end{corollary}

\begin{proof}
Let \(C_1\) be the uniform weak constant in
Corollary~\ref{cor:global-inverse-weak}.  For \(u\ge1\), the union bound gives
\[
  \Prob\{Y_{N,L}^*>u\}
  \le \sum_{\ell=1}^L\Prob\{Y_{N,\ell}>u\}
  \le \frac{LC_1}{u},
\]
which proves the uniform estimate.

For the sharp fixed-window statement, put \(K_0=\sqrt{\pi/3}\).  The
finite-window part of Corollary~\ref{cor:global-inverse-weak} gives
\[
  \varepsilon_N(R):=
  \sup_{1\le u\le R}
  \left(u\Prob\{Y_N>u\}-K_0\right)_+
  \longrightarrow0 .
\]
Hence, uniformly for \(1\le u\le R\),
\[
  u\,\Prob\{Y_{N,L}^*>u\}
  \le L\left(K_0+\varepsilon_N(R)\right).
\]
Taking the supremum in \(u\) and the limsup in \(N\) proves the claim.
\end{proof}

\begin{corollary}[Finite-family inverse-minimum quantiles]
\label{cor:global-inverse-family-quantiles}
Fix \(L\ge1\).  For each \(N\), let
\(Y_{N,1},\ldots,Y_{N,L}\) be positive random variables whose marginals have the
same law as \(Y_N\) in Corollary~\ref{cor:global-inverse-weak}; no independence
between them is required.  Put
\[
  Y_{N,L}^*:=\max_{1\le \ell\le L}Y_{N,\ell}
\]
and, for \(0<\alpha\le1\),
\[
  Q_{N,L}(\alpha):=\inf\{u\ge1:\Prob\{Y_{N,L}^*>u\}\le\alpha\}.
\]
Then there is an absolute constant \(C<\infty\) such that
\[
  Q_{N,L}(\alpha)\le \frac{CL}{\alpha}
  \qquad(N\ge2,\ L\ge1,\ 0<\alpha\le1).
\]
Moreover, for every fixed \(L\) and fixed \(0<\alpha\le1\),
\[
  \limsup_{N\to\infty} Q_{N,L}(\alpha)
  \le \frac{L}{\alpha}\sqrt{\frac{\pi}{3}}.
\]
Equivalently, if \(M_{N,\ell}:=Y_{N,\ell}^{-1}\),
\[
  M_{N,L,*}:=\min_{1\le\ell\le L}M_{N,\ell},
  \qquad
  q_{N,L}(p):=\inf\{s\ge0:\Prob\{M_{N,L,*}\le s\}\ge p\},
\]
then there is \(c>0\) such that
\[
  q_{N,L}(p)\ge c\,\frac{p}{L}
  \qquad(N\ge2,\ L\ge1,\ 0<p<1),
\]
and, for every fixed \(L\) and fixed \(0<p<1\),
\[
  \liminf_{N\to\infty}q_{N,L}(p)
  \ge \frac{p}{L}\sqrt{\frac{3}{\pi}}.
\]
Thus a simultaneous clearance requirement for \(L\) Gaussian symbols costs only
the union-bound factor \(L\) at the inverse-minimum endpoint, with the sharp
single-symbol constant preserved for fixed \(L\).
\end{corollary}

\begin{proof}
Let \(C_1\) be the absolute constant in
Corollary~\ref{cor:global-inverse-family-weak}, enlarged so that \(C_1\ge1\).
That corollary gives \(Q_{N,L}(\alpha)\le C_1L/\alpha\).

For the sharp fixed-\(L\) bound, put \(K_0=\sqrt{\pi/3}\), fix \(\eps>0\), and
set \(u_\eps=(LK_0+\eps)/\alpha\).  The finite-window part of
Corollary~\ref{cor:global-inverse-family-weak}, applied with \(R=u_\eps\), gives
\[
  \limsup_{N\to\infty}\Prob\{Y_{N,L}^*>u_\eps\}
  \le L\,\frac{K_0}{u_\eps}
  <\alpha .
\]
Hence \(Q_{N,L}(\alpha)\le u_\eps\) for all sufficiently large \(N\), and
\(\eps\downarrow0\) proves the displayed limsup estimate.

For the lower-quantile form, write \(F_N(s)=\Prob\{Y_N^{-1}\le s\}\).  The
proof of Corollary~\ref{cor:global-inverse-weak} gives \(F_N(s)\le C_0s\) for
all \(N\ge2\) and \(0<s\le1\).  Therefore
\[
  \Prob\{M_{N,L,*}\le s\}\le L F_N(s)\le LC_0s .
\]
Taking \(s=cp/L\) with \(c\le(2C_0)^{-1}\) proves the uniform lower bound.  For
the sharp bound, if \(a<p\sqrt{3/\pi}/L\), then
Corollary~\ref{cor:global-prob-upper} and the union bound give
\[
  \limsup_{N\to\infty}\Prob\{M_{N,L,*}\le a\}
  \le L\sqrt{\frac{\pi}{3}}\,a
  <p .
\]
Thus \(q_{N,L}(p)\ge a\) eventually, and
\(a\uparrow p\sqrt{3/\pi}/L\) completes the proof.
\end{proof}

\begin{corollary}[Finite-family bad-clearance counts]
\label{cor:global-inverse-family-counts}
In the setting of Corollary~\ref{cor:global-inverse-family-quantiles}, define,
for \(s\ge0\),
\[
  Z_{N,L}(s):=\sum_{\ell=1}^L \mathbf1\{M_{N,\ell}\le s\}.
\]
Then
\[
  \Exp Z_{N,L}(s)\le C Ls
  \qquad(N\ge2,\ L\ge1,\ s\ge0)
\]
with an absolute constant \(C<\infty\).  More sharply, for every fixed
\(L\ge1\) and \(K<\infty\),
\[
  \limsup_{N\to\infty}\;
  \sup_{0\le s\le K}
  \left(
    \Exp Z_{N,L}(s)-L\sqrt{\frac{\pi}{3}}\,s
  \right)_+
  =0 .
\]
Consequently, for every integer \(r\ge1\),
\[
  \Prob\{Z_{N,L}(s)\ge r\}\le \frac{CLs}{r},
\]
and, for fixed \(L,K,r\),
\[
  \limsup_{N\to\infty}\;
  \sup_{0\le s\le K}
  \left(
    r\,\Prob\{Z_{N,L}(s)\ge r\}
    -L\sqrt{\frac{\pi}{3}}\,s
  \right)_+
  =0 .
\]
Thus finite simultaneous testing not only controls the first bad clearance but
also the expected number of bad clearances, with the same sharp linear
coefficient and no independence assumption.
\end{corollary}

\begin{proof}
Let \(F_N(s)=\Prob\{Y_N^{-1}\le s\}\).  Since each \(M_{N,\ell}\) has the same
law as \(Y_N^{-1}\),
\[
  \Exp Z_{N,L}(s)=\sum_{\ell=1}^L \Prob\{M_{N,\ell}\le s\}=L F_N(s).
\]
The proof of Corollary~\ref{cor:global-inverse-weak} gives \(F_N(s)\le C_0s\)
for \(0<s\le1\), and the case \(s=0\) is trivial since \(F_N(0)=0\).  Also
\(F_N(s)\le1\le s\) for \(s\ge1\).  This proves the uniform bound after
increasing \(C\).

For the sharp fixed-window statement, Corollary~\ref{cor:global-prob-upper}
gives, for every \(K<\infty\),
\[
  \limsup_{N\to\infty}\;
  \sup_{0\le s\le K}
  \left(F_N(s)-\sqrt{\frac{\pi}{3}}\,s\right)_+
  =0 .
\]
Multiplying by \(L\) gives the expectation estimate.  The two probability
bounds follow from Markov's inequality
\(r\Prob\{Z_{N,L}(s)\ge r\}\le \Exp Z_{N,L}(s)\).
\end{proof}

\begin{corollary}[Finite-family clearance order statistics]
\label{cor:global-inverse-family-order-quantiles}
In the setting of Corollary~\ref{cor:global-inverse-family-quantiles}, let
\[
  M_{N,L}^{(1)}\le M_{N,L}^{(2)}\le\cdots\le M_{N,L}^{(L)}
\]
be the increasing rearrangement of \(M_{N,1},\ldots,M_{N,L}\).  For
\(1\le r\le L\) and \(0<p<1\), define
\[
  q_{N,L,r}(p):=
  \inf\{s\ge0:\Prob\{M_{N,L}^{(r)}\le s\}\ge p\}.
\]
Then there is an absolute constant \(c>0\) such that
\[
  q_{N,L,r}(p)\ge c\,\frac{pr}{L}
  \qquad(N\ge2,\ 1\le r\le L,\ 0<p<1).
\]
Moreover, for every fixed \(L\ge1\), \(1\le r\le L\), and \(0<p<1\),
\[
  \liminf_{N\to\infty}q_{N,L,r}(p)
  \ge \frac{pr}{L}\sqrt{\frac{3}{\pi}}.
\]
Equivalently, the threshold at which \(r\) out of \(L\) simultaneous symbols
can have bad boundary clearance is at least \(r/L\) times the one-symbol
continuum scale, up to the same fixed-window constant, without any independence
assumption among the \(L\) symbols.
\end{corollary}

\begin{proof}
For every \(s\ge0\),
\[
  \{M_{N,L}^{(r)}\le s\}=\{Z_{N,L}(s)\ge r\}.
\]
Corollary~\ref{cor:global-inverse-family-counts} gives
\[
  \Prob\{M_{N,L}^{(r)}\le s\}\le \frac{CLs}{r}.
\]
Taking \(s=cp r/L\), with \(c\le(2C)^{-1}\), proves the uniform bound.

For the fixed-\(L,r\) asymptotic statement, put \(K_0=\sqrt{\pi/3}\).  If
\(0<a<pr/(LK_0)\), then Corollary~\ref{cor:global-inverse-family-counts}, used
with \(K=a\), gives
\[
  \limsup_{N\to\infty}\Prob\{M_{N,L}^{(r)}\le a\}
  \le \frac{LK_0a}{r}
  <p .
\]
Thus \(q_{N,L,r}(p)\ge a\) eventually, and
\(a\uparrow pr/(LK_0)\) completes the proof.
\end{proof}

\begin{corollary}[Finite-family truncated and logarithmic inverse moments]
\label{cor:global-inverse-family-moments}
In the setting of Corollary~\ref{cor:global-inverse-family-quantiles}, for
every \(R\ge1\),
\[
  \sup_{N\ge2}\Exp\bigl[(Y_{N,L}^*\wedge R)-1\bigr]_+
  \le C L\log R
\]
with an absolute constant \(C<\infty\).  More sharply, for every fixed
\(L\ge1\) and fixed \(1\le R<\infty\),
\[
  \limsup_{N\to\infty}
  \Exp\bigl[(Y_{N,L}^*\wedge R)-1\bigr]_+
  \le L\sqrt{\frac{\pi}{3}}\,\log R .
\]
Also, for every \(p>0\),
\[
  \sup_{N\ge2}\Exp\bigl[(\log_+Y_{N,L}^*)^p\bigr]
  \le C_p L,
\]
and, for every fixed \(L\ge1\),
\[
  \limsup_{N\to\infty}\Exp\bigl[(\log_+Y_{N,L}^*)^p\bigr]
  \le L\sqrt{\frac{\pi}{3}}\,\Gamma(p+1).
\]
Thus finite simultaneous testing preserves the endpoint \(q=1\) obstruction:
the price of \(L\) marginals is only the linear union-bound factor in the
critical truncation and logarithmic-integrability constants.
\end{corollary}

\begin{proof}
For any nonnegative random variable \(Y\),
\[
  \Exp\bigl[(Y\wedge R)-1\bigr]_+
  =\int_1^R\Prob\{Y>u\}\,du .
\]
The union-bound tail estimate from
Corollary~\ref{cor:global-inverse-family-quantiles},
\[
  \Prob\{Y_{N,L}^*>u\}\le LC_1/u\qquad(u\ge1),
\]
gives the uniform truncated bound.  For the sharp fixed-window coefficient, the
finite-window part of Corollary~\ref{cor:global-inverse-weak} gives, for fixed
\(R\),
\[
  \varepsilon_N(R):=
  \sup_{1\le u\le R}
  \left(u\Prob\{Y_N>u\}-\sqrt{\frac{\pi}{3}}\right)_+
  \longrightarrow0 .
\]
Hence, uniformly for \(1\le u\le R\),
\[
  \Prob\{Y_{N,L}^*>u\}
  \le L\left(\sqrt{\frac{\pi}{3}}+\varepsilon_N(R)\right)u^{-1}.
\]
Integrating over \(1\le u\le R\) and taking the limsup proves the truncated
estimate.

For logarithmic moments, use layer cake:
\[
  \Exp\bigl[(\log_+Y_{N,L}^*)^p\bigr]
  =
  p\int_0^\infty t^{p-1}\Prob\{Y_{N,L}^*>e^t\}\,dt .
\]
The uniform tail bound gives
\[
  \Prob\{Y_{N,L}^*>e^t\}\le LC_1 e^{-t},
\]
so the uniform \(C_pL\) estimate follows.  For the sharp limsup, fix
\(T<\infty\).  The preceding finite-window bound with \(R=e^T\) yields
\[
  \limsup_{N\to\infty}
  p\int_0^T t^{p-1}\Prob\{Y_{N,L}^*>e^t\}\,dt
  \le
  L\sqrt{\frac{\pi}{3}}\,
  p\int_0^T t^{p-1}e^{-t}\,dt .
\]
The tail \(t\ge T\) is at most
\[
  LC_1p\int_T^\infty t^{p-1}e^{-t}\,dt
\]
uniformly in \(N\).  Letting \(T\to\infty\) gives
\(L\sqrt{\pi/3}\,\Gamma(p+1)\).
\end{proof}

\begin{corollary}[Finite-family subcritical inverse powers]
\label{cor:global-inverse-family-subcritical}
In the setting of Corollary~\ref{cor:global-inverse-family-quantiles}, for
every \(0<q<1\),
\[
  \sup_{N\ge2}\Exp\bigl[(Y_{N,L}^*)^q\bigr]
  =
  \sup_{N\ge2}\Exp M_{N,L,*}^{-q}
  \le C_q L,
\]
where \(C_q<\infty\) depends only on \(q\).  More sharply, for every fixed
\(L\ge1\) and every \(a>0\),
\[
  \limsup_{N\to\infty}
  \Exp\!\left[M_{N,L,*}^{-q}\mathbf1\{M_{N,L,*}\le a\}\right]
  \le
  \frac{L\sqrt{\pi/3}}{1-q}\,a^{1-q}.
\]
Thus finite simultaneous clearance preserves the same subcritical endpoint
exponent; \(L\) marginals only multiply the sharp small-clearance coefficient by
\(L\).
\end{corollary}

\begin{proof}
Put
\[
  F_{N,L}(s)=\Prob\{M_{N,L,*}\le s\},
  \qquad F_N(s)=\Prob\{M_N\le s\}.
\]
The proof of Corollary~\ref{cor:global-inverse-family-quantiles} gives a
constant \(C_0\) such that \(F_N(s)\le C_0s\) for all \(N\ge2\) and
\(0<s\le1\).  Hence, by the union bound,
\[
  F_{N,L}(s)\le L F_N(s)\le LC_0s\qquad(0<s\le1).
\]
For \(u\ge1\), this gives
\[
  \Prob\{Y_{N,L}^*>u\}
  \le F_{N,L}(1/u)
  \le LC_0/u.
\]
Therefore
\[
  \Exp\bigl[(Y_{N,L}^*)^q\bigr]
  =
  q\int_0^\infty u^{q-1}\Prob\{Y_{N,L}^*>u\}\,du
  \le
  1+qLC_0\int_1^\infty u^{q-2}\,du
  \le C_q L .
\]

For the sharp small-window estimate, fix
\(0<\rho<\min\{a,1\}\).  Stieltjes integration by parts gives
\[
  \Exp\!\left[M_{N,L,*}^{-q}\mathbf1\{M_{N,L,*}\le a\}\right]
  =
  a^{-q}F_{N,L}(a)+q\int_0^a s^{-q-1}F_{N,L}(s)\,ds .
\]
On \((0,\rho)\) use \(F_{N,L}(s)\le LC_0s\).  On \([\rho,a]\), the union bound
and Corollary~\ref{cor:global-prob-upper} give
\[
  \limsup_{N\to\infty}\sup_{\rho\le s\le a}
  \bigl(F_{N,L}(s)-L\sqrt{\pi/3}\,s\bigr)_+=0 .
\]
Taking the limsup in \(N\) yields
\[
\begin{aligned}
  \limsup_{N\to\infty}
  \Exp\!\left[M_{N,L,*}^{-q}\mathbf1\{M_{N,L,*}\le a\}\right]
  &\le
  L\sqrt{\frac{\pi}{3}}\,a^{1-q}
  +\frac{qLC_0}{1-q}\rho^{1-q}
  \\
  &\quad
  +\frac{qL\sqrt{\pi/3}}{1-q}
  \left(a^{1-q}-\rho^{1-q}\right).
\end{aligned}
\]
Letting \(\rho\downarrow0\) proves the claim.
\end{proof}

\begin{corollary}[Summable subcritical inverse-clearance mass]
\label{cor:global-inverse-family-subcritical-sum}
For each \(N\ge2\), let \(L_N\ge1\), and let
\(Y_{N,1},\ldots,Y_{N,L_N}\) be defined on a common probability space, each with
the same marginal law as \(Y_N\) from
Corollary~\ref{cor:global-inverse-weak}.  Put \(M_{N,\ell}=Y_{N,\ell}^{-1}\).
For \(s\ge0\), put
\[
  Z_N(s):=\sum_{\ell=1}^{L_N}\mathbf1\{M_{N,\ell}\le s\}.
\]
For \(0<q<1\) and \(a\ge0\), define the subcritical inverse-clearance mass
\[
  S_N^{(q)}(a):=
  \sum_{\ell=1}^{L_N}M_{N,\ell}^{-q}\mathbf1\{M_{N,\ell}\le a\}.
\]
Then, for \(0\le a\le1\),
\[
  \Exp S_N^{(q)}(a)\le C_q L_Na^{1-q},
\]
where \(C_q<\infty\) depends only on \(q\).  More sharply, if \(L_N\equiv L\),
then for every fixed \(L\ge1\) and \(0<a<\infty\),
\[
  \limsup_{N\to\infty}\Exp S_N^{(q)}(a)
  \le
  \frac{L\sqrt{\pi/3}}{1-q}\,a^{1-q}.
\]
Consequently, for any deterministic \(b_N>0\), if \(0\le a_N\le1\) and
\[
  \sum_{N=2}^\infty \frac{L_Na_N^{1-q}}{b_N}<\infty,
\]
then
\[
  \Exp\sum_{N=2}^\infty \frac{S_N^{(q)}(a_N)}{b_N}<\infty,
  \qquad
  \sum_{N=2}^\infty \frac{S_N^{(q)}(a_N)}{b_N}<\infty
  \quad\text{almost surely}.
\]
In particular \(S_N^{(q)}(a_N)=o(b_N)\) almost surely.  If
\(L_N\le C_LN^\beta\) and \(b_N\ge c_bN^{-\delta}\), then for every \(\eps>0\),
\[
  S_N^{(q)}\!\left(
    N^{-(\beta+\delta+1)/(1-q)-\eps}
  \right)
  =
  o\!\left(N^{-\delta}\right)
  \qquad\text{almost surely}.
\]
For actual Gaussian symbols this says that the total subcritical inverse mass of
symbols with
\[
  \min_{z\in\T}|G_m^{(\ell)}(z)|
  \le
  N^{-(\beta+\delta+1)/(1-q)-1/2-\eps}
\]
is \(o(N^{-\delta})\) after the natural \(N^{-q/2}\) scaling.  No independence
is required within a batch or across \(N\).
\end{corollary}

\begin{proof}
The case \(a=0\) is trivial.  For \(x>0\) and \(a>0\),
\[
  x^{-q}\mathbf1\{x\le a\}
  =
  a^{-q}\mathbf1\{x\le a\}
  +q\int_0^a s^{-q-1}\mathbf1\{x\le s\}\,ds,
\]
with the evident zero value when \(a=0\).  Therefore Tonelli's theorem gives
\[
  \Exp S_N^{(q)}(a)
  =
  a^{-q}\Exp Z_N(a)
  +q\int_0^a s^{-q-1}\Exp Z_N(s)\,ds .
\]
For \(a\le1\), Corollary~\ref{cor:global-inverse-family-counts} gives
\(\Exp Z_N(s)\le CL_Ns\), hence
\[
  \Exp S_N^{(q)}(a)
  \le
  CL_Na^{1-q}+CqL_N\int_0^a s^{-q}\,ds
  \le C_qL_Na^{1-q}.
\]

For the sharp fixed-\(L\) bound, put \(K_0=\sqrt{\pi/3}\) and fix
\(0<\rho<\min\{a,1\}\).  On \((0,\rho)\), use the uniform bound
\(\Exp Z_N(s)\le CLs\).  On \([\rho,a]\), Corollary~\ref{cor:global-inverse-family-counts}
gives
\[
  \limsup_{N\to\infty}
  \sup_{\rho\le s\le a}
  \left(\Exp Z_N(s)-LK_0s\right)_+
  =0 .
\]
Taking the limsup in \(N\) in the preceding layer-cake formula yields
\[
\begin{aligned}
  \limsup_{N\to\infty}\Exp S_N^{(q)}(a)
  &\le
  LK_0a^{1-q}
  +\frac{qCL}{1-q}\rho^{1-q}  \\
  &\quad
  +\frac{qLK_0}{1-q}\left(a^{1-q}-\rho^{1-q}\right).
\end{aligned}
\]
Letting \(\rho\downarrow0\) gives the sharp coefficient.

The normalized summability assertion follows by monotone convergence:
\[
  \Exp\sum_{N=2}^\infty \frac{S_N^{(q)}(a_N)}{b_N}
  =
  \sum_{N=2}^\infty \frac{\Exp S_N^{(q)}(a_N)}{b_N}
  \le
  C_q\sum_{N=2}^\infty \frac{L_Na_N^{1-q}}{b_N}
  <\infty .
\]
Thus the nonnegative normalized sum is finite almost surely, so its terms tend
to zero.  The polynomial specialization takes
\[
  a_N=N^{-(\beta+\delta+1)/(1-q)-\eps}.
\]
Finally \(M_{N,\ell}=\sqrt N\min_{\T}|G_m^{(\ell)}|\) gives the stated unscaled
interpretation.
\end{proof}

\begin{corollary}[Summable Stieltjes inverse-clearance gauge mass]
\label{cor:global-inverse-family-gauge-sum}
For each \(N\ge2\), let \(L_N\ge1\), and let
\(Y_{N,1},\ldots,Y_{N,L_N}\) be defined on a common probability space, each with
the same marginal law as \(Y_N\) from
Corollary~\ref{cor:global-inverse-weak}.  Put \(M_{N,\ell}=Y_{N,\ell}^{-1}\).
Let \(a>0\), and let \(\psi:(0,a]\to[0,\infty)\) be nonincreasing and
right-continuous.  Let \(\nu_\psi\) be the positive Stieltjes measure on
\((0,a]\) determined by
\[
  \nu_\psi((u,v])=\psi(u)-\psi(v)\qquad(0<u<v\le a),
\]
and suppose
\[
  B_\psi(a):=a\psi(a)+\int_{(0,a]}s\,d\nu_\psi(s)<\infty .
\]
Define
\[
  T_N^\psi(a):=
  \sum_{\ell=1}^{L_N}\psi(M_{N,\ell})\mathbf1\{M_{N,\ell}\le a\}.
\]
Then
\[
  \Exp T_N^\psi(a)\le C L_N B_\psi(a),
\]
with an absolute constant \(C<\infty\).  More sharply, if \(L_N\equiv L\),
then for every fixed \(L\ge1\), \(a<\infty\), and such \(\psi\),
\[
  \limsup_{N\to\infty}\Exp T_N^\psi(a)
  \le
  L\sqrt{\frac{\pi}{3}}\,B_\psi(a).
\]
Consequently, for deterministic \(b_N>0\), if
\(\psi_N:(0,a_N]\to[0,\infty)\) are nonincreasing right-continuous gauges, with
associated Stieltjes measures \(\nu_{\psi_N}\) and \(B_{\psi_N}(a_N)<\infty\),
and
\[
  \sum_{N=2}^\infty
  \frac{L_NB_{\psi_N}(a_N)}{b_N}<\infty,
\]
then
\[
  \Exp\sum_{N=2}^\infty \frac{T_N^{\psi_N}(a_N)}{b_N}<\infty,
  \qquad
  \sum_{N=2}^\infty \frac{T_N^{\psi_N}(a_N)}{b_N}<\infty
  \quad\text{almost surely}.
\]
In particular \(T_N^{\psi_N}(a_N)=o(b_N)\) almost surely.  In the absolutely
continuous case \(d\nu_\psi(s)=(-\psi'(s))\,ds\).  The choice
\(\psi(s)=s^{-q}\), \(0<q<1\), gives
\(B_\psi(a)=a^{1-q}/(1-q)\) and recovers the subcritical power scale.  The
choices \(\psi(s)=s^{-q}(\log(e/s))^\gamma\), \(0<q<1\), \(\gamma\ge0\), and
\(\psi(s)=(\log(e/s))^\gamma\), \(\gamma>0\), give the corresponding logarithmic
refinements with \(B_\psi(a)\asymp_{q,\gamma}a^{1-q}(\log(e/a))^\gamma\) and
\(B_\psi(a)\asymp_\gamma a(\log(e/a))^\gamma\), respectively, for
\(0<a\le1\).
\end{corollary}

\begin{proof}
For \(s\ge0\), put
\[
  Z_N(s):=\sum_{\ell=1}^{L_N}\mathbf1\{M_{N,\ell}\le s\}.
\]
For \(x>0\),
\[
  \psi(x)\mathbf1\{x\le a\}
  =
  \psi(a)\mathbf1\{x\le a\}
  +
  \int_{(0,a]}\mathbf1\{x<s\}\,d\nu_\psi(s) .
\]
Tonelli's theorem therefore gives
\[
  \Exp T_N^\psi(a)
  \le
  \psi(a)\Exp Z_N(a)
  +
  \int_{(0,a]}\Exp Z_N(s)\,d\nu_\psi(s) .
\]
The uniform count estimate in Corollary~\ref{cor:global-inverse-family-counts}
gives \(\Exp Z_N(s)\le CL_Ns\), hence
\[
  \Exp T_N^\psi(a)
  \le
  CL_N\left(a\psi(a)+\int_{(0,a]}s\,d\nu_\psi(s)\right).
\]

For the sharp fixed-\(L\) estimate, set \(K_0=\sqrt{\pi/3}\) and fix
\(0<\rho<a\).  Corollary~\ref{cor:global-inverse-family-counts} gives
\[
  \varepsilon_N(\rho,a):=
  \sup_{\rho\le s\le a}
  \left(\Exp Z_N(s)-LK_0s\right)_+\longrightarrow0 .
\]
Since \(\nu_\psi([\rho,a])<\infty\), the error term contributes at most
\(\varepsilon_N(\rho,a)\nu_\psi([\rho,a])\), and the preceding bound yields
\[
\begin{aligned}
  \limsup_{N\to\infty}\Exp T_N^\psi(a)
  &\le
  LK_0a\psi(a)
  +CL\int_{(0,\rho)}s\,d\nu_\psi(s)  \\
  &\quad
  +LK_0\int_{[\rho,a]}s\,d\nu_\psi(s) .
\end{aligned}
\]
Letting \(\rho\downarrow0\) and using \(B_\psi(a)<\infty\) proves the sharp
bound.  Monotone convergence gives the normalized summability statement exactly
as in Corollary~\ref{cor:global-inverse-family-subcritical-sum}.  The displayed
examples follow by direct one-variable integration.
\end{proof}

\begin{corollary}[Inverse-gauge confidence budgets]
\label{cor:global-inverse-family-gauge-confidence}
In the setting of Corollary~\ref{cor:global-inverse-family-gauge-sum}, let
\[
  Q_N^\psi(\alpha)
  :=
  \inf\{r\ge0:\Prob\{T_N^\psi(a)>r\}\le\alpha\},
  \qquad 0<\alpha\le1.
\]
Then
\[
  Q_N^\psi(\alpha)\le \frac{C L_NB_\psi(a)}{\alpha}
  \qquad(N\ge2,\ 0<\alpha\le1),
\]
with an absolute constant \(C<\infty\).  More sharply, if \(L_N\equiv L\), then
for every fixed \(L\ge1\), \(a<\infty\), gauge \(\psi\), and
\(0<\alpha\le1\),
\[
  \limsup_{N\to\infty}Q_N^\psi(\alpha)
  \le
  \frac{L\sqrt{\pi/3}\,B_\psi(a)}{\alpha}.
\]
Consequently, for a triangular family \(\psi_N,a_N\) as in
Corollary~\ref{cor:global-inverse-family-gauge-sum} and any
\(0<\eta_N\le1\),
\[
  \Prob\left\{
    T_N^{\psi_N}(a_N)>
    \frac{C L_NB_{\psi_N}(a_N)}{\eta_N}
  \right\}\le \eta_N .
\]
For every finite index set \(I\), the preceding bounds hold simultaneously for
all \(N\in I\) outside an event of probability at most \(\sum_{N\in I}\eta_N\).
If \(\sum_N\eta_N<\infty\), then almost surely, for all sufficiently large
\(N\),
\[
  T_N^{\psi_N}(a_N)\le
  \frac{C L_NB_{\psi_N}(a_N)}{\eta_N}.
\]
No independence is required in the batch or across \(N\).
\end{corollary}

\begin{proof}
Markov's inequality and Corollary~\ref{cor:global-inverse-family-gauge-sum}
give
\[
  \Prob\{T_N^\psi(a)>r\}
  \le
  \frac{C L_NB_\psi(a)}{r},
\]
which proves the uniform quantile estimate after taking
\(r=C L_NB_\psi(a)/\alpha\), with \(C\) enlarged if necessary.  If \(L_N\equiv L\),
the sharp expectation bound in
Corollary~\ref{cor:global-inverse-family-gauge-sum} gives, for every
\(\eps>0\),
\[
  \limsup_{N\to\infty}
  \Prob\left\{
    T_N^\psi(a)>
    \frac{L\sqrt{\pi/3}\,B_\psi(a)+\eps}{\alpha}
  \right\}
  <\alpha .
\]
Thus \(Q_N^\psi(\alpha)\le
(L\sqrt{\pi/3}\,B_\psi(a)+\eps)/\alpha\) for all sufficiently large \(N\);
letting \(\eps\downarrow0\) proves the fixed-\(L\) quantile coefficient.

The displayed one-step confidence bound is the same Markov estimate with
\(r=C L_NB_{\psi_N}(a_N)/\eta_N\).  The finite-horizon statement is the union
bound, and the almost-sure eventual statement is the first Borel--Cantelli
lemma.
\end{proof}

\begin{corollary}[Simultaneous inverse-gauge menu budgets]
\label{cor:global-inverse-family-gauge-menu}
For each \(N\ge2\), let \(L_N\ge1\), and let
\(Y_{N,1},\ldots,Y_{N,L_N}\) have the common marginal law from
Corollary~\ref{cor:global-inverse-weak}, with no independence assumption.  Put
\(M_{N,\ell}=Y_{N,\ell}^{-1}\).  Let \(J_N\) be a finite or countable
deterministic index set.  For each \(j\in J_N\), let
\(\psi_{N,j}:(0,a_{N,j}]\to[0,\infty)\) be a nonincreasing right-continuous
gauge with Stieltjes budget \(B_{N,j}:=B_{\psi_{N,j}}(a_{N,j})<\infty\), and
put
\[
  T_{N,j}:=
  \sum_{\ell=1}^{L_N}
  \psi_{N,j}(M_{N,\ell})\mathbf1\{M_{N,\ell}\le a_{N,j}\}.
\]
There is an absolute constant \(C<\infty\) such that, for any finite set
\(\mathcal I\subset\{(N,j):N\ge2,\ j\in J_N\}\) and any
\(0<\eta_{N,j}\le1\),
\[
  \Prob\left\{
    \exists (N,j)\in\mathcal I:
    T_{N,j}>\frac{C L_NB_{N,j}}{\eta_{N,j}}
  \right\}
  \le
  \sum_{(N,j)\in\mathcal I}\eta_{N,j}.
\]
If
\[
  \sum_{N=2}^\infty\sum_{j\in J_N}\eta_{N,j}<\infty ,
\]
then almost surely, for all but finitely many pairs \((N,j)\),
\[
  T_{N,j}\le \frac{C L_NB_{N,j}}{\eta_{N,j}}.
\]
In particular, for a single \(N\), a finite menu \(J_N\), and
\(0<\alpha\le1\), choosing \(\eta_{N,j}=\alpha/|J_N|\) gives simultaneous
control of the whole menu with failure probability at most \(\alpha\):
\[
  T_{N,j}\le \frac{C |J_N| L_NB_{N,j}}{\alpha}
  \qquad(j\in J_N).
\]
Thus power, logarithmic, and jump gauges may be tested on the same batch with
only the explicit union-budget cost.
\end{corollary}

\begin{proof}
Apply Corollary~\ref{cor:global-inverse-family-gauge-confidence} to each
\((N,j)\), with confidence parameter \(\eta_{N,j}\), and sum the resulting
probability bounds over \(\mathcal I\).  This proves the finite-menu assertion.
For the countable almost-sure statement, enumerate the countable set of pairs
\((N,j)\) with \(\eta_{N,j}>0\).  The displayed summability makes the expected
number of failures finite, equivalently the sum of their probabilities finite;
the first Borel--Cantelli lemma gives only finitely many failures almost surely.
The final finite-menu specialization is the same bound with
\(\eta_{N,j}=\alpha/|J_N|\).
\end{proof}

\begin{corollary}[Aggregate inverse-gauge menu mass]
\label{cor:global-inverse-family-gauge-menu-sum}
In the setting of Corollary~\ref{cor:global-inverse-family-gauge-menu}, set
\[
  \mathcal T_N:=\sum_{j\in J_N}T_{N,j},
  \qquad
  \mathcal B_N:=\sum_{j\in J_N}B_{N,j}.
\]
If \(\mathcal B_N<\infty\), then
\[
  \Exp \mathcal T_N\le C L_N\mathcal B_N
\]
with an absolute constant \(C<\infty\).  More sharply, if \(L_N\equiv L\) and
the menu \(J\), gauges \(\psi_j\), windows \(a_j\), and budgets \(B_j\) are
fixed with \(\sum_{j\in J}B_j<\infty\), then
\[
  \limsup_{N\to\infty}\Exp\sum_{j\in J}T_{N,j}
  \le
  L\sqrt{\frac{\pi}{3}}\sum_{j\in J}B_j .
\]
Consequently, for any deterministic \(b_N>0\), if
\[
  \sum_{N=2}^\infty \frac{L_N\mathcal B_N}{b_N}<\infty,
\]
then
\[
  \Exp\sum_{N=2}^\infty\frac{\mathcal T_N}{b_N}<\infty,
  \qquad
  \sum_{N=2}^\infty\frac{\mathcal T_N}{b_N}<\infty
  \quad\text{almost surely}.
\]
In particular \(\mathcal T_N=o(b_N)\) almost surely.  Thus a countable menu of
power, logarithmic, and jump gauges may be aggregated as a single endpoint
diagnostic whenever its Stieltjes budgets are summable.
\end{corollary}

\begin{proof}
Apply Corollary~\ref{cor:global-inverse-family-gauge-sum} to each \(j\) and
sum.  Monotone convergence justifies the countable sum:
\[
  \Exp\mathcal T_N
  =
  \sum_{j\in J_N}\Exp T_{N,j}
  \le
  C L_N\sum_{j\in J_N}B_{N,j}.
\]
For the sharp fixed-\(L\) statement, enumerate \(J\) as
\(\{j_1,j_2,\ldots\}\), with the evident finite truncation when \(J\) is finite.
For a finite initial segment, the sharp bound in
Corollary~\ref{cor:global-inverse-family-gauge-sum} gives
\[
  \limsup_{N\to\infty}\Exp\sum_{m=1}^M T_{N,j_m}
  \le
  L\sqrt{\frac{\pi}{3}}\sum_{m=1}^M B_{j_m}.
\]
The remaining tail has expectation at most
\[
  C L\sum_{m>M}B_{j_m}
\]
uniformly in \(N\).  Letting \(M\to\infty\) proves the sharp countable-menu
coefficient.  The normalized summability assertion follows from the first
display and monotone convergence, and finiteness of the nonnegative series
implies its terms tend to zero.
\end{proof}

\begin{corollary}[Summable finite-family clearance schedules]
\label{cor:global-inverse-family-borel}
For each \(N\ge2\), let \(L_N\ge1\), and let
\(Y_{N,1},\ldots,Y_{N,L_N}\) be defined on a common probability space, each with
the same marginal law as \(Y_N\) from
Corollary~\ref{cor:global-inverse-weak}.  Put \(M_{N,\ell}=Y_{N,\ell}^{-1}\),
\[
  M_{N,L_N,*}:=\min_{1\le\ell\le L_N}M_{N,\ell},
  \qquad
  Y_{N,L_N}^*:=\max_{1\le\ell\le L_N}Y_{N,\ell}.
\]
If \(0<a_N\le1\) and
\[
  \sum_{N=2}^\infty L_N a_N<\infty,
\]
then
\[
  \Prob\{M_{N,L_N,*}\le a_N\ \text{for infinitely many }N\}=0.
\]
Equivalently, almost surely \(Y_{N,L_N}^*<a_N^{-1}\) for all sufficiently large
\(N\).  In particular, if \(L_N\le C_LN^\beta\), then for every
\(\eps>0\),
\[
  M_{N,L_N,*}>N^{-\beta-1-\eps}
\]
eventually almost surely.  For actual Gaussian symbols this means
\[
  \min_{1\le\ell\le L_N}\min_{z\in\T}|G_m^{(\ell)}(z)|
  > N^{-\beta-3/2-\eps}
\]
eventually almost surely, with no independence required among the symbols or
across \(N\).
\end{corollary}

\begin{proof}
As in Corollary~\ref{cor:global-inverse-family-subcritical}, there is an
absolute \(C_0\) such that, for \(0<s\le1\),
\[
  \Prob\{M_{N,L_N,*}\le s\}\le C_0L_Ns .
\]
Thus
\[
  \sum_{N=2}^\infty
  \Prob\{M_{N,L_N,*}\le a_N\}
  \le
  C_0\sum_{N=2}^\infty L_Na_N
  <\infty .
\]
The first Borel--Cantelli lemma gives the summable-schedule statement.  The
polynomial example follows by taking \(a_N=N^{-\beta-1-\eps}\), since then
\(\sum_N L_Na_N\le C_L\sum_NN^{-1-\eps}<\infty\).  Dividing the scaled
clearance \(M_{N,L_N,*}\) by \(\sqrt N\) gives the final unscaled form.
\end{proof}

\begin{corollary}[Summable total bad-clearance mass]
\label{cor:global-inverse-family-count-sum}
For each \(N\ge2\), let \(L_N\ge1\), and let
\(Y_{N,1},\ldots,Y_{N,L_N}\) be defined on a common probability space, each with
the same marginal law as \(Y_N\) from
Corollary~\ref{cor:global-inverse-weak}.  Put \(M_{N,\ell}=Y_{N,\ell}^{-1}\) and
\[
  Z_N(s):=\sum_{\ell=1}^{L_N}\mathbf1\{M_{N,\ell}\le s\}.
\]
If \(0<a_N\le1\) and
\[
  \sum_{N=2}^\infty L_Na_N<\infty,
\]
then
\[
  \Exp\sum_{N=2}^\infty Z_N(a_N)<\infty
  \qquad\text{and hence}\qquad
  \sum_{N=2}^\infty Z_N(a_N)<\infty
  \quad\text{almost surely}.
\]
Equivalently, only finitely many individual pairs \((N,\ell)\) have
\(M_{N,\ell}\le a_N\), almost surely.  In particular, if
\(L_N\le C_LN^\beta\), then for every \(\eps>0\),
\[
  \sum_{N=2}^\infty
  \sum_{\ell=1}^{L_N}
  \mathbf1\{M_{N,\ell}\le N^{-\beta-1-\eps}\}<\infty
  \qquad\text{almost surely}.
\]
For actual Gaussian symbols, this says that only finitely many individual
\((N,\ell)\) satisfy
\[
  \min_{z\in\T}|G_m^{(\ell)}(z)|
  \le N^{-\beta-3/2-\eps}.
\]
No independence is required within a batch or across \(N\).
\end{corollary}

\begin{proof}
Corollary~\ref{cor:global-inverse-family-counts} gives an absolute constant
\(C\) such that, for \(0<s\le1\),
\[
  \Exp Z_N(s)\le CL_Ns .
\]
Therefore
\[
  \Exp\sum_{N=2}^\infty Z_N(a_N)
  =
  \sum_{N=2}^\infty \Exp Z_N(a_N)
  \le
  C\sum_{N=2}^\infty L_Na_N
  <\infty ,
\]
where monotone convergence justifies the exchange.  A nonnegative random
variable with finite expectation is finite almost surely.  The polynomial
specialization takes \(a_N=N^{-\beta-1-\eps}\), and the unscaled Gaussian form
divides the scaled clearance by \(\sqrt N\).
\end{proof}

\begin{corollary}[Summable clearance-deficit mass]
\label{cor:global-inverse-family-deficit-sum}
In the setting of Corollary~\ref{cor:global-inverse-family-count-sum}, define
the total deficit below level \(a\ge0\) by
\[
  D_N(a):=\sum_{\ell=1}^{L_N}(a-M_{N,\ell})_+ .
\]
For \(0\le a\le1\),
\[
  \Exp D_N(a)\le C L_Na^2
\]
with an absolute constant \(C<\infty\).  More sharply, for every fixed
\(L\ge1\) and \(K<\infty\), if \(L_N\equiv L\), then
\[
  \limsup_{N\to\infty}\;
  \sup_{0\le a\le K}
  \left(
    \Exp D_N(a)-\frac{L}{2}\sqrt{\frac{\pi}{3}}\,a^2
  \right)_+
  =0 .
\]
Consequently, if \(0\le a_N\le1\) and
\[
  \sum_{N=2}^\infty L_Na_N^2<\infty,
\]
then
\[
  \Exp\sum_{N=2}^\infty D_N(a_N)<\infty
  \qquad\text{and hence}\qquad
  \sum_{N=2}^\infty D_N(a_N)<\infty
  \quad\text{almost surely}.
\]
In particular, if \(L_N\le C_LN^\beta\), then for every \(\eps>0\),
\[
  \sum_{N=2}^\infty
  \sum_{\ell=1}^{L_N}
  \left(N^{-(\beta+1)/2-\eps}-M_{N,\ell}\right)_+
  <\infty
  \qquad\text{almost surely}.
\]
For actual Gaussian symbols this is equivalently
\[
  \sum_{N=2}^\infty
  \sum_{\ell=1}^{L_N}
  \sqrt N\left(
    N^{-(\beta+2)/2-\eps}
    -\min_{z\in\T}|G_m^{(\ell)}(z)|
  \right)_+
  <\infty
  \qquad\text{almost surely}.
\]
Thus one can sum shortfall sizes below a much larger clearance threshold than
the threshold for summing the bad-clearance indicators themselves.  No
independence is required within a batch or across \(N\).
\end{corollary}

\begin{proof}
For each \(a\ge0\),
\[
  (a-M_{N,\ell})_+=\int_0^a \mathbf1\{M_{N,\ell}\le s\}\,ds,
\]
and hence Tonelli's theorem gives
\[
  \Exp D_N(a)=\int_0^a \Exp Z_N(s)\,ds .
\]
Corollary~\ref{cor:global-inverse-family-counts} gives
\(\Exp Z_N(s)\le CL_Ns\) for \(0\le s\le1\), which proves
\(\Exp D_N(a)\le CL_Na^2\) for \(0\le a\le1\), after changing \(C\).

For the sharp fixed-\(L\) statement, put \(K_0=\sqrt{\pi/3}\).  By
Corollary~\ref{cor:global-inverse-family-counts},
\[
  \varepsilon_N(K):=
  \sup_{0\le s\le K}
  \left(\Exp Z_N(s)-LK_0s\right)_+
  \longrightarrow0 .
\]
Therefore, uniformly for \(0\le a\le K\),
\[
  \Exp D_N(a)
  \le \int_0^a (LK_0s+\varepsilon_N(K))\,ds
  \le \frac{LK_0}{2}a^2+K\varepsilon_N(K).
\]
Taking the limsup proves the displayed sharp bound.

If \(\sum_N L_Na_N^2<\infty\), then monotone convergence and the uniform
deficit estimate give
\[
  \Exp\sum_{N=2}^\infty D_N(a_N)
  =
  \sum_{N=2}^\infty \Exp D_N(a_N)
  \le C\sum_{N=2}^\infty L_Na_N^2
  <\infty .
\]
Thus the nonnegative sum is finite almost surely.  The polynomial example takes
\(a_N=N^{-(\beta+1)/2-\eps}\), so
\(L_Na_N^2\le C_LN^{-1-2\eps}\).  Finally
\[
  M_{N,\ell}
  =
  \sqrt N\,\min_{z\in\T}|G_m^{(\ell)}(z)|
\]
turns the scaled deficit into the displayed unscaled Gaussian form.
\end{proof}

\begin{corollary}[Normalized clearance-deficit budgets]
\label{cor:global-inverse-family-deficit-budget}
In the setting of Corollary~\ref{cor:global-inverse-family-deficit-sum}, let
\(b_N>0\) be deterministic.  If \(0\le a_N\le1\) and
\[
  \sum_{N=2}^\infty \frac{L_Na_N^2}{b_N}<\infty,
\]
then
\[
  \Exp\sum_{N=2}^\infty \frac{D_N(a_N)}{b_N}<\infty,
  \qquad
  \sum_{N=2}^\infty \frac{D_N(a_N)}{b_N}<\infty
  \quad\text{almost surely}.
\]
In particular,
\[
  D_N(a_N)=o(b_N)
  \qquad\text{almost surely}.
\]
Equivalently, the total shortfall below \(a_N\) is eventually negligible
relative to any deterministic budget \(b_N\) whose reciprocal still leaves
\(\sum_N L_Na_N^2/b_N\) summable.  If \(L_N\le C_LN^\beta\) and
\(b_N\ge c_bN^{-\delta}\) with \(\delta\ge0\), then for every \(\eps>0\),
\[
  D_N\!\left(N^{-(\beta+\delta+1)/2-\eps}\right)
  =o\!\left(N^{-\delta}\right)
  \qquad\text{almost surely}.
\]
For actual Gaussian symbols this becomes
\[
  \sum_{\ell=1}^{L_N}
  \sqrt N\left(
    N^{-(\beta+\delta+2)/2-\eps}
    -\min_{z\in\T}|G_m^{(\ell)}(z)|
  \right)_+
  =
  o\!\left(N^{-\delta}\right)
  \qquad\text{almost surely}.
\]
No independence is required within a batch or across \(N\).
\end{corollary}

\begin{proof}
Corollary~\ref{cor:global-inverse-family-deficit-sum} gives
\(\Exp D_N(a_N)\le CL_Na_N^2\).  Hence, by monotone convergence,
\[
  \Exp\sum_{N=2}^\infty \frac{D_N(a_N)}{b_N}
  =
  \sum_{N=2}^\infty \frac{\Exp D_N(a_N)}{b_N}
  \le
  C\sum_{N=2}^\infty \frac{L_Na_N^2}{b_N}
  <\infty .
\]
The nonnegative normalized sum is therefore finite almost surely, which implies
\(D_N(a_N)/b_N\to0\).  For the polynomial specialization take
\[
  a_N=N^{-(\beta+\delta+1)/2-\eps}.
\]
Then
\[
  \frac{L_Na_N^2}{b_N}
  \le \frac{C_L}{c_b}N^{-1-2\eps},
\]
which is summable.  Finally
\[
  M_{N,\ell}=\sqrt N\,\min_{z\in\T}|G_m^{(\ell)}(z)|
\]
gives the unscaled Gaussian form.
\end{proof}

\begin{corollary}[Powered clearance-deficit budgets]
\label{cor:global-inverse-family-deficit-powers}
In the setting of Corollary~\ref{cor:global-inverse-family-deficit-sum}, let
\(p>0\) and define
\[
  D_N^{(p)}(a):=
  \sum_{\ell=1}^{L_N}(a-M_{N,\ell})_+^p .
\]
For \(0\le a\le1\),
\[
  \Exp D_N^{(p)}(a)\le C_p L_Na^{p+1}
\]
with a constant \(C_p<\infty\) depending only on \(p\).  More sharply, for
every fixed \(L\ge1\) and \(K<\infty\), if \(L_N\equiv L\), then
\[
  \limsup_{N\to\infty}\;
  \sup_{0\le a\le K}
  \left(
    \Exp D_N^{(p)}(a)
    -\frac{L}{p+1}\sqrt{\frac{\pi}{3}}\,a^{p+1}
  \right)_+
  =0 .
\]
Consequently, for any deterministic \(b_N>0\), if \(0\le a_N\le1\) and
\[
  \sum_{N=2}^\infty \frac{L_Na_N^{p+1}}{b_N}<\infty,
\]
then
\[
  \Exp\sum_{N=2}^\infty \frac{D_N^{(p)}(a_N)}{b_N}<\infty,
  \qquad
  \sum_{N=2}^\infty \frac{D_N^{(p)}(a_N)}{b_N}<\infty
  \quad\text{almost surely},
\]
and hence \(D_N^{(p)}(a_N)=o(b_N)\) almost surely.  If
\(L_N\le C_LN^\beta\) and \(b_N\ge c_bN^{-\delta}\), then for every
\(\eps>0\),
\[
  D_N^{(p)}\!\left(N^{-(\beta+\delta+1)/(p+1)-\eps}\right)
  =o\!\left(N^{-\delta}\right)
  \qquad\text{almost surely}.
\]
For actual Gaussian symbols this is equivalently
\[
  \sum_{\ell=1}^{L_N}
  N^{p/2}\left(
    N^{-(\beta+\delta+1)/(p+1)-1/2-\eps}
    -\min_{z\in\T}|G_m^{(\ell)}(z)|
  \right)_+^p
  =
  o\!\left(N^{-\delta}\right)
  \qquad\text{almost surely}.
\]
The case \(p=1\) recovers Corollary~\ref{cor:global-inverse-family-deficit-budget}.
No independence is required within a batch or across \(N\).
\end{corollary}

\begin{proof}
For \(p>0\) and \(a\ge0\),
\[
  (a-M_{N,\ell})_+^p
  =
  p\int_0^a (a-s)^{p-1}\mathbf1\{M_{N,\ell}\le s\}\,ds .
\]
Thus Tonelli's theorem gives
\[
  \Exp D_N^{(p)}(a)
  =
  p\int_0^a (a-s)^{p-1}\Exp Z_N(s)\,ds .
\]
For \(0\le a\le1\), Corollary~\ref{cor:global-inverse-family-counts} gives
\(\Exp Z_N(s)\le CL_Ns\).  Therefore
\[
  \Exp D_N^{(p)}(a)
  \le
  CL_Np\int_0^a (a-s)^{p-1}s\,ds
  =
  \frac{CL_N}{p+1}a^{p+1},
\]
which is the uniform estimate after absorbing constants into \(C_p\).

For the sharp fixed-\(L\) statement, put \(K_0=\sqrt{\pi/3}\).  As in
Corollary~\ref{cor:global-inverse-family-deficit-sum}, let
\[
  \varepsilon_N(K):=
  \sup_{0\le s\le K}
  \left(\Exp Z_N(s)-LK_0s\right)_+
  \longrightarrow0 .
\]
Then, uniformly for \(0\le a\le K\),
\[
\begin{aligned}
  \Exp D_N^{(p)}(a)
  &\le
  p\int_0^a (a-s)^{p-1}(LK_0s+\varepsilon_N(K))\,ds  \\
  &=
  \frac{LK_0}{p+1}a^{p+1}+\varepsilon_N(K)a^p
  \le
  \frac{LK_0}{p+1}a^{p+1}+\varepsilon_N(K)K^p .
\end{aligned}
\]
Taking the limsup proves the sharp coefficient.

The normalized summability assertion follows by monotone convergence:
\[
  \Exp\sum_{N=2}^\infty \frac{D_N^{(p)}(a_N)}{b_N}
  =
  \sum_{N=2}^\infty \frac{\Exp D_N^{(p)}(a_N)}{b_N}
  \le
  C_p\sum_{N=2}^\infty \frac{L_Na_N^{p+1}}{b_N}
  <\infty .
\]
Hence the nonnegative normalized sum is finite almost surely, so its terms tend
to zero.  For the polynomial specialization take
\[
  a_N=N^{-(\beta+\delta+1)/(p+1)-\eps}.
\]
Then
\[
  \frac{L_Na_N^{p+1}}{b_N}
  \le \frac{C_L}{c_b}N^{-1-(p+1)\eps},
\]
which is summable.  Finally
\[
  M_{N,\ell}=\sqrt N\,\min_{z\in\T}|G_m^{(\ell)}(z)|
\]
gives the unscaled Gaussian form.
\end{proof}

\begin{corollary}[Jump-weighted clearance-deficit functionals]
\label{cor:global-inverse-family-deficit-weighted}
In the setting of Corollary~\ref{cor:global-inverse-family-deficit-sum}, let
\(\Phi:[0,K]\to[0,\infty)\) be nondecreasing, right-continuous on \((0,K]\), and
satisfy \(\Phi(0)=0\).  Put
\[
  \Delta_0\Phi:=\lim_{u\downarrow0}\Phi(u),
\]
and let \(\mu_\Phi\) be the Stieltjes measure on \((0,K]\) determined by
\(\mu_\Phi((0,v])=\Phi(v)-\Delta_0\Phi\) and
\(\mu_\Phi((u,v])=\Phi(v)-\Phi(u)\) for \(0<u<v\).  For \(0\le a\le K\), define
\[
  W_N^\Phi(a):=
  \sum_{\ell=1}^{L_N}\Phi\bigl((a-M_{N,\ell})_+\bigr)
\]
and
\[
  H_\Phi(a):=
  a\,\Delta_0\Phi+\int_{(0,a]} (a-u)\,d\mu_\Phi(u).
\]
If \(0\le a\le\min\{K,1\}\), then
\[
  \Exp W_N^\Phi(a)\le C L_N H_\Phi(a)
\]
with an absolute constant \(C<\infty\).  More sharply, for every fixed
\(L\ge1\) and \(K<\infty\), if \(L_N\equiv L\), then
\[
  \limsup_{N\to\infty}\;
  \sup_{0\le a\le K}
  \left(
    \Exp W_N^\Phi(a)
    -L\sqrt{\frac{\pi}{3}}\,H_\Phi(a)
  \right)_+
  =0 .
\]
Consequently, let \(0<a_N\le1\), \(b_N>0\), and let
\(\Phi_N:[0,a_N]\to[0,\infty)\) be nondecreasing, right-continuous on
\((0,a_N]\), and satisfy \(\Phi_N(0)=0\).  Let
\(\Delta_N=\lim_{u\downarrow0}\Phi_N(u)\), let \(\mu_N\) be the corresponding
Stieltjes measure on \((0,a_N]\), and put
\[
  H_N:=a_N\Delta_N+\int_{(0,a_N]}(a_N-u)\,d\mu_N(u).
\]
If
\[
  \sum_{N=2}^\infty
  \frac{L_NH_N}{b_N}
  <\infty ,
\]
then
\[
  \sum_{N=2}^\infty
  \frac{1}{b_N}
  \sum_{\ell=1}^{L_N}
  \Phi_N\bigl((a_N-M_{N,\ell})_+\bigr)
  <\infty
  \qquad\text{almost surely}.
\]
Thus any summable weighted shortfall budget follows from the same one-dimensional
occupation integral, now allowing an entrance jump at zero as well as later
jump penalties and smooth weights.  If \(\Delta_0\Phi=0\) and \(\Phi\) is
absolutely continuous, then
\[
  H_\Phi(a)=\int_0^a\Phi'(a-s)s\,ds .
\]
Taking \(\Phi(u)=u^p\) recovers
Corollary~\ref{cor:global-inverse-family-deficit-powers}; taking \(\Phi(u)=u\)
recovers the first-power deficit estimates; and taking
\(\Phi(u)=\mathbf1_{\{u>0\}}\) recovers the strict bad-clearance count scale.
\end{corollary}

\begin{proof}
For \(0\le x\le K\),
\[
  \Phi(x)=\Delta_0\Phi\,\mathbf1\{x>0\}
  +\int_{(0,K]}\mathbf1\{u\le x\}\,d\mu_\Phi(u).
\]
Therefore, for \(0\le a\le K\),
\[
  \Phi\bigl((a-M_{N,\ell})_+\bigr)
  =
  \Delta_0\Phi\,\mathbf1\{M_{N,\ell}<a\}
  +\int_{(0,a]}\mathbf1\{M_{N,\ell}\le a-u\}\,d\mu_\Phi(u).
\]
Summing over \(\ell\) and applying Tonelli's theorem gives
\[
  \Exp W_N^\Phi(a)
  \le
  \Delta_0\Phi\,\Exp Z_N(a)+
  \int_{(0,a]}\Exp Z_N(a-u)\,d\mu_\Phi(u).
\]
For \(a\le1\), Corollary~\ref{cor:global-inverse-family-counts} gives
\(\Exp Z_N(s)\le CL_Ns\), which proves the uniform estimate.

For the sharp fixed-\(L\) statement, put \(K_0=\sqrt{\pi/3}\) and
\[
  \varepsilon_N(K):=
  \sup_{0\le s\le K}
  \left(\Exp Z_N(s)-LK_0s\right)_+
  \longrightarrow0 .
\]
Then, uniformly for \(0\le a\le K\),
\[
\begin{aligned}
  \Exp W_N^\Phi(a)
  &\le
  LK_0H_\Phi(a)+\varepsilon_N(K)\bigl(\Delta_0\Phi+\mu_\Phi((0,a])\bigr) \\
  &\le
  LK_0H_\Phi(a)+\varepsilon_N(K)\Phi(K).
\end{aligned}
\]
Taking the limsup proves the sharp bound.

The schedule statement follows by applying the uniform estimate with
\(\Phi=\Phi_N\) and \(a=a_N\), then using monotone convergence:
\[
\begin{aligned}
  \Exp\sum_{N=2}^\infty
  \frac{1}{b_N}
  \sum_{\ell=1}^{L_N}\Phi_N\bigl((a_N-M_{N,\ell})_+\bigr)
  &\le
  C\sum_{N=2}^\infty
  \frac{L_NH_N}{b_N}  \\
  &<\infty .
\end{aligned}
\]
The nonnegative sum is therefore finite almost surely.  The smooth
specializations follow by evaluating \(H_\Phi\); the indicator specialization has
\(\Delta_0\Phi=1\) and \(\mu_\Phi=0\).
\end{proof}

\begin{corollary}[Summable bad-count schedules]
\label{cor:global-inverse-family-count-borel}
For each \(N\ge2\), let \(L_N\ge1\), and let
\(Y_{N,1},\ldots,Y_{N,L_N}\) be defined on a common probability space, each with
the same marginal law as \(Y_N\) from
Corollary~\ref{cor:global-inverse-weak}.  Put \(M_{N,\ell}=Y_{N,\ell}^{-1}\) and
\[
  Z_N(s):=\sum_{\ell=1}^{L_N}\mathbf1\{M_{N,\ell}\le s\}.
\]
Let \(1\le r_N\le L_N\) be deterministic integers.  If \(0<a_N\le1\) and
\[
  \sum_{N=2}^\infty \frac{L_Na_N}{r_N}<\infty,
\]
then
\[
  \Prob\{Z_N(a_N)\ge r_N\ \text{for infinitely many }N\}=0.
\]
Equivalently, the \(r_N\)-th smallest clearance among the \(L_N\) symbols is
larger than \(a_N\) for all sufficiently large \(N\), almost surely.  In
particular, if \(L_N\le C_LN^\beta\) and \(r_N\ge c_rN^\gamma\) with
\(\gamma\ge0\), then for every \(\eps>0\),
\[
  Z_N\!\left(N^{-\beta+\gamma-1-\eps}\right)<r_N
\]
eventually almost surely.  Thus a polynomially growing allowed bad-count budget
raises the eventual simultaneous clearance threshold by the factor
\(r_N/L_N\), again without any independence assumption.
\end{corollary}

\begin{proof}
By Corollary~\ref{cor:global-inverse-family-counts}, there is an absolute
constant \(C\) such that, for \(0<s\le1\),
\[
  \Prob\{Z_N(s)\ge r_N\}\le \frac{CL_Ns}{r_N}.
\]
Therefore
\[
  \sum_{N=2}^\infty \Prob\{Z_N(a_N)\ge r_N\}
  \le
  C\sum_{N=2}^\infty \frac{L_Na_N}{r_N}
  <\infty .
\]
The first Borel--Cantelli lemma gives the general assertion.  For the
polynomial consequence, take
\(a_N=N^{-\beta+\gamma-1-\eps}\).  Then
\[
  \frac{L_Na_N}{r_N}
  \le \frac{C_L}{c_r}N^{-1-\eps},
\]
which is summable.
\end{proof}

\begin{corollary}[Critical truncated inverse moment]
\label{cor:global-inverse-truncated}
With \(Y_N=M_N^{-1}\) as in Corollary~\ref{cor:global-inverse-weak}, for every
\(R\ge1\),
\[
  \sup_{N\ge2}\Exp\bigl[(Y_N\wedge R)-1\bigr]_+
  \le C\log R
\]
with an absolute constant \(C<\infty\).  More sharply, for every fixed
\(1\le R<\infty\),
\[
  \limsup_{N\to\infty}
  \Exp\bigl[(Y_N\wedge R)-1\bigr]_+
  \le \sqrt{\frac{\pi}{3}}\,\log R .
\]
Equivalently, the critical \(q=1\) inverse moment grows at most logarithmically
under truncation, while Corollary~\ref{cor:global-negative-moments} gives strong
\(L^q\) control for every \(q<1\).
\end{corollary}

\begin{proof}
For a nonnegative random variable \(Y\),
\[
  \Exp\bigl[(Y\wedge R)-1\bigr]_+
  =\int_1^R \Prob\{Y>u\}\,du .
\]
The uniform weak bound in Corollary~\ref{cor:global-inverse-weak} gives the
first assertion.  For the sharp finite-window coefficient, the second assertion
of Corollary~\ref{cor:global-inverse-weak} implies that, for fixed \(R\),
\[
  \varepsilon_N(R):=
  \sup_{1\le u\le R}
  \left(
    u\,\Prob\{Y_N>u\}-\sqrt{\frac{\pi}{3}}
  \right)_+
  \longrightarrow0 .
\]
Hence
\[
  \Exp\bigl[(Y_N\wedge R)-1\bigr]_+
  \le
  \left(\sqrt{\frac{\pi}{3}}+\varepsilon_N(R)\right)\log R,
\]
and the limsup in \(N\) gives the claim.
\end{proof}

\begin{corollary}[Logarithmic moments at the inverse endpoint]
\label{cor:global-inverse-logmoments}
With \(Y_N=M_N^{-1}\) as in Corollary~\ref{cor:global-inverse-weak}, for every
\(p>0\),
\[
  \sup_{N\ge2}\Exp\bigl[(\log_+Y_N)^p\bigr]<\infty ,
\]
where \(\log_+u:=\max\{\log u,0\}\).  More sharply,
\[
  \limsup_{N\to\infty}\Exp\bigl[(\log_+Y_N)^p\bigr]
  \le \sqrt{\frac{\pi}{3}}\,\Gamma(p+1).
\]
Thus the weak-\(L^1\) endpoint still gives full logarithmic integrability of
the inverse scaled minimum, with the same sharp upper-tail constant.
\end{corollary}

\begin{proof}
For any positive random variable \(Y\),
\[
  \Exp\bigl[(\log_+Y)^p\bigr]
  =
  p\int_0^\infty t^{p-1}\Prob\{Y>e^t\}\,dt .
\]
The uniform weak bound in Corollary~\ref{cor:global-inverse-weak} gives
\(\Prob\{Y_N>e^t\}\le Ce^{-t}\) for all \(t\ge0\), uniformly in \(N\), which
proves the first assertion.

For the sharp limsup, fix \(T<\infty\).  The finite-window part of
Corollary~\ref{cor:global-inverse-weak}, applied with \(R=e^T\), gives
\[
  \limsup_{N\to\infty}
  p\int_0^T t^{p-1}\Prob\{Y_N>e^t\}\,dt
  \le
  \sqrt{\frac{\pi}{3}}\,
  p\int_0^T t^{p-1}e^{-t}\,dt .
\]
The tail \(t\ge T\) is bounded uniformly in \(N\) by
\[
  Cp\int_T^\infty t^{p-1}e^{-t}\,dt ,
\]
again by the weak bound.  Letting \(T\to\infty\) gives the displayed estimate,
since \(p\int_0^\infty t^{p-1}e^{-t}\,dt=\Gamma(p+1)\).
\end{proof}

\begin{corollary}[Endpoint floor for the global minimum]
\label{cor:endpoint-global-floor}
For standard real Gaussian coefficients there are constants \(c>0\) and
\(m_0<\infty\) such that, for \(m\ge m_0\) and \(0<t\le1\),
\[
  \Prob\Bigl\{\min_{z\in\T}|G_m(z)|\le t\Bigr\}
  \ge c\,\frac{t}{\sqrt m}.
\]
Equivalently,
\[
  \Prob\Bigl\{\min_{z\in\T}|G_m(z)|\le \frac{\tau}{\sqrt m}\Bigr\}
  \ge c\,\frac{\tau}{m}
  \qquad(0<\tau\le \sqrt m).
\]
Thus the two parity endpoints already force a nonzero global lower floor, but
only at order \(t/\sqrt m\); Corollaries~\ref{cor:endpoint-two-point},
\ref{cor:endpoint-continuum-negligible}, and
\ref{cor:endpoint-neighborhood-continuum} give the exact bivariate formula,
sharp two-endpoint expansion, and continuum-scale \(O_\tau(N^{-1})\)
contribution of both the endpoint values and their natural neighborhoods.  The
conjectural whole-circle lower tail at order \(\sqrt m\,t\) must come from the
\(O(m)\) non-real arcs and their decorrelation.
\end{corollary}

\begin{proof}
The event in Corollary~\ref{cor:endpoint-two-point} is contained in
\(\{\min_{\T}\abs{G_m}\le t\}\), and its lower asymptotic is
\(\gg t/\sqrt m\).  The rescaled statement is the same estimate with
\(t=\tau/\sqrt m\).
\end{proof}

\begin{proof}[Proof of Theorem~\ref{thm:uniform}]
Fix endpoint arcs \(E_{C_0}\) with \(C_0\) large enough that the complement can
be partitioned into \(O(m)\) arcs \(I_\ell\) of lengths \(c_\ell/m\), with
\(1\le c_\ell\le2\), each at angular distance at least \(1/m\) from
\(\{\pm1\}\).  Applying Theorem~\ref{thm:singlearc}, then summing over the
\(O(m)\) arcs, gives
\[
  \Prob\Bigl\{\min_{\theta\notin E_{C_0}}\abs{G_m(e^{i\theta})}\le t\Bigr\}
  \le C\sqrt m\,t.
\]
By Proposition~\ref{prop:endpoint-singlearc},
\[
  \Prob\Bigl\{\min_{\theta\in E_{C_0}}\abs{G_m(e^{i\theta})}\le t\Bigr\}
  \le C\frac{t}{\sqrt m}.
\]
Combining the endpoint and non-endpoint estimates proves \eqref{eq:uniform}.
\end{proof}

\begin{proof}[Proof of Corollary~\ref{prop:uniform-polylog}]
This is immediate from Theorem~\ref{thm:uniform}, after increasing
\(C_{A,B}\) and \(m_0(A,B)\) if necessary.
\end{proof}

\begin{lemma}[Arc sampling under a speed cutoff]\label{lem:arcsampling}
Let $I\subset\T$ be an arc of arclength $\abs I=c/m$, sampled at $L:=\lceil 4c\rceil$ equally spaced nodes $z_1,\dots,z_L\in I$. For every $z\in I$ there is a node $z_\ell$ with
\[
  \bigl|\,\abs{G_m(z)}-\abs{G_m(z_\ell)}\,\bigr|\;\le\;\abs{G_m(z)-G_m(z_\ell)}\;\le\;\frac{\abs I}{L}\,\sup_{w\in I}\abs{G_m'(w)}.
\]
\end{lemma}

\begin{proof}
This is the fundamental theorem of calculus along \(I\), since each \(z\in I\)
lies within arclength \(\abs I/L\) of a node and
\(\abs{\,\abs a-\abs b\,}\le\abs{a-b}\).
\end{proof}

\begin{proposition}[Bounded-density local bounds under a speed cutoff]\label{prop:bd-speed}
Let \(\xi_0,\ldots,\xi_m\) be independent real random variables with densities
bounded by \(K\), and let \(G_m(z)=\sum_{j=0}^m\xi_jz^j\).

\emph{(i) Away from the parity points.}  Fix \(\gamma>0\).  If \(I\subset\T\)
is an arc of arclength \(c/m\), \(c\ge1\), with
\(\dist_\T(I,\{\pm1\})\ge\gamma/m\), then for \(L\ge1\) and \(0<t\le1\),
\[
  \Prob\Bigl(
    \Bigl\{\min_{z\in I}\abs{G_m(z)}\le t\Bigr\}
    \cap
    \Bigl\{\sup_{z\in I}\abs{G_m'(z)}\le Lm^{3/2}\Bigr\}
  \Bigr)
  \le C_{\gamma,K}(1+c)\frac{Lt}{\sqrt m}.
\]

\emph{(ii) Endpoint arcs.}  Fix \(C_0,B>0\), and put
\[
  E_{C_0}:=\{\theta\in[-\pi,\pi]:
  \dist(\theta,\{0,\pi,-\pi\})\le C_0/m\}.
\]
For \(m^{-B}\le t\le1\) and \(L\ge1\),
\[
  \Prob\Bigl(
    \Bigl\{\min_{\theta\in E_{C_0}}\abs{G_m(e^{i\theta})}\le t\Bigr\}
    \cap
    \Bigl\{\sup_{\T}\abs{G_m'}\le Lm^{3/2}\Bigr\}
  \Bigr)
  \le C_{B,C_0,K}\frac{Lt\log m}{\sqrt m}.
\]
\end{proposition}

\begin{proof}
For (i), parametrize \(I\) by an interval \(J\) of length \(c/m\) and set
\(\Lambda(s)=\int_J\mathbf1\{\abs{G_m(e^{i\theta})}\le s\}\,d\theta\).  By
Fubini and Proposition~\ref{prop:nongaussian-grid},
\[
  \Exp\Lambda(s)\le \frac{C_{\gamma,K}c s^2}{m^2}.
\]
On the speed event, one visit to \(\{\abs{G_m}\le t\}\) forces
\(\abs{G_m}\le2t\) on an interval of length at least a constant multiple of
\(\delta=t/(Lm^{3/2})\), truncated only by a harmless endpoint factor because
\(\abs J=c/m\ge\delta\) for large \(m\).  Hence
\[
  \mathbf1\{\min_I\abs{G_m}\le t\}\mathbf1_{\{\sup_I\abs{G_m'}\le Lm^{3/2}\}}
  \le C\,\frac{\Lambda(2t)}{\delta}.
\]
Taking expectations gives \(C_{\gamma,K}cLt/\sqrt m\), and \(c\le1+c\).

For (ii), the same argument uses the endpoint occupation estimate
Lemma~\ref{lem:endpoint-occ} instead of the away-from-endpoint occupation
estimate.  On the speed event a visit contributes length
\(\delta=t/(Lm^{3/2})\) inside \(E_{C_0}\), up to a constant one-sided loss at
the ends of the endpoint arcs, and Lemma~\ref{lem:endpoint-occ} with \(s=2t\)
gives
\[
  C\,\frac{Lm^{3/2}}{t}\cdot
  C_{B,C_0,K}\frac{t^2\log m}{m^2}
  \le C_{B,C_0,K}\frac{Lt\log m}{\sqrt m}.
\]
\end{proof}

\begin{corollary}[Uniform bounded-density boundary bound under a speed cutoff]
\label{cor:bd-uniform-speed}
Let \(\xi_0,\ldots,\xi_m\) be independent real random variables with densities
bounded by \(K\), and let \(G_m(z)=\sum_{j=0}^m\xi_jz^j\).  For every
\(B>0\) there are constants \(C_{B,K}<\infty\) and \(m_0(B,K)\) such that, for
\(m\ge m_0(B,K)\), \(m^{-B}\le t\le1\), and \(L\ge1\),
\[
  \Prob\Bigl(
    \Bigl\{\min_{z\in\T}\abs{G_m(z)}\le t\Bigr\}
    \cap
    \Bigl\{\sup_{z\in\T}\abs{G_m'(z)}\le Lm^{3/2}\Bigr\}
  \Bigr)
  \le C_{B,K}L\sqrt m\,t .
\]
\end{corollary}

\begin{proof}
Choose endpoint arcs \(E_{C_0}\) with fixed \(C_0\) large enough that the
complement can be partitioned into \(O(m)\) arcs of lengths \(c_\ell/m\), with
\(1\le c_\ell\le2\), all at angular distance at least \(1/m\) from
\(\{\pm1\}\).  Sum Proposition~\ref{prop:bd-speed}(i) over those arcs, using
\(\gamma=1\), to obtain \(C_KL\sqrt m\,t\) on the complement.  On the endpoint
arcs, Proposition~\ref{prop:bd-speed}(ii) gives
\[
  C_{B,C_0,K}\frac{Lt\log m}{\sqrt m}\le C_{B,K}L\sqrt m\,t
\]
for \(m\) large.  Combining the two estimates proves the claim.
\end{proof}

\begin{lemma}[Polynomial derivative cutoff from variance]\label{lem:variance-speed}
Let \(\xi_0,\ldots,\xi_m\) be independent mean-zero real random variables with
\(\Exp|\xi_j|^2=1\), and let \(G_m(z)=\sum_{j=0}^m\xi_jz^j\). For every
\(\delta>1/2\), there is a constant \(C_\delta<\infty\) such that
\[
  \Prob\Bigl\{\sup_{z\in\T}|G_m'(z)|>m^{3/2+\delta}\Bigr\}
  \le C_\delta m^{1-2\delta}.
\]
\end{lemma}

\begin{proof}
Put
\[
  H_m(\theta):=G_m'(e^{i\theta})
  =\sum_{j=1}^m j\xi_j e^{i(j-1)\theta}.
\]
Bernstein's inequality for trigonometric polynomials gives a grid
\(\{\theta_\ell\}_{\ell=1}^N\) with \(N\le Cm\) and mesh at most \(1/(4m)\)
such that
\[
  \sup_\theta |H_m(\theta)|
  \le 2\max_{1\le\ell\le N}|H_m(\theta_\ell)|.
\]
Indeed, at a maximum point choose the nearest grid point and use
\(\|H_m'\|_\infty\le m\|H_m\|_\infty\).
For each fixed grid point, independence and centering give
\[
  \Exp |H_m(\theta_\ell)|^2
  =\sum_{j=1}^m j^2\Exp|\xi_j|^2
  \le C m^3 .
\]
Chebyshev's inequality therefore gives
\[
  \Prob\{|H_m(\theta_\ell)|>\tfrac12 m^{3/2+\delta}\}
  \le C m^{-2\delta}.
\]
The union bound over \(N\le Cm\) grid points yields
\[
  \Prob\{\sup_\theta |H_m(\theta)|>m^{3/2+\delta}\}
  \le C m^{1-2\delta}=o(1),
\]
because \(\delta>1/2\).
\end{proof}

\begin{corollary}[Finite-confidence derivative cutoff from variance]
\label{cor:variance-speed-confidence}
Let \(\xi_0,\ldots,\xi_m\) be independent mean-zero real random variables with
\(\Exp|\xi_j|^2=1\), and let \(G_m(z)=\sum_{j=0}^m\xi_jz^j\).  There is an
absolute constant \(C<\infty\) such that, for every \(R>0\),
\[
  \Prob\Bigl\{\sup_{z\in\T}|G_m'(z)|>R\Bigr\}
  \le C\min\left\{1,\frac{m^4}{R^2}\right\}.
\]
Consequently, for every \(0<\eta\le1\),
\[
  \Prob\Bigl\{\sup_{z\in\T}|G_m'(z)|
  > C m^2\eta^{-1/2}\Bigr\}\le \eta .
\]
\end{corollary}

\begin{proof}
The proof of Lemma~\ref{lem:variance-speed}, before the final substitution
\(R=m^{3/2+\delta}\), gives the quantitative estimate.  Indeed, Bernstein's
inequality for trigonometric polynomials reduces the supremum of
\[
  H_m(\theta)=\frac{d}{d\theta}G_m(e^{i\theta})
\]
to a grid of \(O(m)\) points, and for each grid point
\(\Exp |H_m(\theta_\ell)|^2\le Cm^3\).  Chebyshev's inequality and the union
bound therefore give
\[
  \Prob\{\sup_\theta |H_m(\theta)|>R\}
  \le \frac{Cm^4}{R^2}.
\]
Since \(|H_m(\theta)|=|G_m'(e^{i\theta})|\) on \(\T\), this is the first
display, after taking the minimum with the trivial bound \(1\).  The
finite-confidence form follows by taking \(R=Cm^2\eta^{-1/2}\) with \(C\)
large enough.
\end{proof}

\begin{corollary}[Tail-free polynomial boundary separation]
\label{cor:variance-boundary-separation}
Let \(\xi_0,\ldots,\xi_m\) be independent mean-zero real random variables whose
densities are bounded by \(K\) and satisfy \(\Exp|\xi_j|^2=1\).  If
\(G_m(z)=\sum_{j=0}^m\xi_jz^j\), then, for every \(A>1\),
\[
  \Prob\Bigl\{\min_{z\in\T}\abs{G_m(z)}\le m^{-A}\Bigr\}
  \le C_{A,K}\,m^{-2(A-1)/3}.
\]
Equivalently, every bounded-density variance-one law has a polynomial
continuum separation floor \(m^{-A}\), for any prescribed exponent \(A>1\),
without any tail assumption, and with an explicit polynomial exceptional
probability.
\end{corollary}

\begin{proof}
Choose
\[
  \delta=\frac{A+1/2}{3}.
\]
Since \(A>1\), we have \(1/2<\delta<A-1/2\).  Let
\[
  E_\delta=\Bigl\{\sup_{\T}\abs{G_m'}\le m^{3/2+\delta}\Bigr\}.
\]
By Lemma~\ref{lem:variance-speed},
\[
  \Prob(E_\delta^c)\le C_A m^{1-2\delta}.
\]
On
\(E_\delta\), apply Corollary~\ref{cor:bd-uniform-speed} with
\(L=m^\delta\), \(B=A\), and \(t=m^{-A}\).  This gives
\[
  \Prob\Bigl(
    \Bigl\{\min_{\T}\abs{G_m}\le m^{-A}\Bigr\}\cap E_\delta
  \Bigr)
  \le C_{A,K}m^{\delta+1/2-A},
\]
and the chosen \(\delta\) makes
\[
  1-2\delta=\delta+\frac12-A=-\frac{2(A-1)}{3}.
\]
Adding the two estimates gives the displayed bound.
\end{proof}

\begin{corollary}[Finite-confidence boundary separation from variance]
\label{cor:variance-boundary-confidence}
Let \(\xi_0,\ldots,\xi_m\) be independent mean-zero real random variables whose
densities are bounded by \(K\) and satisfy \(\Exp|\xi_j|^2=1\), and let
\(G_m(z)=\sum_{j=0}^m\xi_jz^j\).  Fix \(B>0\).  There are constants
\(C_{B,K}<\infty\) and \(m_0(B,K)\) such that, for all \(m\ge m_0(B,K)\),
all \(0<\eta\le1\), and all \(m^{-B}\le t\le1\),
\[
  \Prob\Bigl\{\min_{z\in\T}\abs{G_m(z)}\le t\Bigr\}
  \le
  \eta+C_{B,K}\frac{mt}{\sqrt\eta}.
\]
In particular, if \(m^{-B}\le t\le c_{B,K}\eta^{3/2}/m\), then
\[
  \Prob\Bigl\{\min_{z\in\T}\abs{G_m(z)}\le t\Bigr\}\le2\eta .
\]
\end{corollary}

\begin{proof}
By Corollary~\ref{cor:variance-speed-confidence}, outside an event of
probability at most \(\eta\),
\[
  \sup_{\T}|G_m'|\le C m^2\eta^{-1/2}
  =
  Lm^{3/2},
  \qquad
  L=Cm^{1/2}\eta^{-1/2}.
\]
Let
\[
  E_\eta:=\left\{\sup_{\T}|G_m'|\le C m^2\eta^{-1/2}\right\}.
\]
On this speed event, Corollary~\ref{cor:bd-uniform-speed} gives
\[
  \Prob\left(\left\{\min_{\T}\abs{G_m}\le t\right\}\cap E_\eta\right)
  \le C_{B,K}Lm^{1/2}t
  \le C_{B,K}\frac{mt}{\sqrt\eta}.
\]
Adding the speed-exception probability proves the first display.  The second
display follows by choosing \(c_{B,K}\) so that
\(C_{B,K}mt/\sqrt\eta\le\eta\).
\end{proof}

\begin{corollary}[Variance-only polynomial grids detect continuum excursions]
\label{cor:variance-grid-detects}
Let \(\xi_0,\ldots,\xi_m\) be independent mean-zero real random variables with
\(\Exp|\xi_j|^2=1\), let \(G_m(z)=\sum_{j=0}^m\xi_jz^j\), and put \(N=m+1\).
Fix \(\delta>1/2\) and \(\tau>0\).  If \(Q_m\subset\T\) is an \(h_m\)-net in
arclength with
\[
  h_m\le \frac{\tau}{2m^{3/2+\delta}\sqrt N},
\]
then there is a constant \(C_\delta<\infty\) such that
\[
  \Prob\left(
    \Bigl\{\min_{z\in\T}\abs{G_m(z)}\le \frac{\tau}{2\sqrt N}\Bigr\}
    \setminus
    \Bigl\{\min_{q\in Q_m}\abs{G_m(q)}\le \frac{\tau}{\sqrt N}\Bigr\}
  \right)
  \le C_\delta m^{1-2\delta}.
\]
Thus a uniform grid with \(M\ge C m^{3/2+\delta}\sqrt N/\tau\) points, in
particular \(M=O_\tau(N^{2+\delta})\), has this detection property with the
same exceptional-probability bound.
\end{corollary}

\begin{proof}
Let
\[
  E_m=\Bigl\{\sup_{\T}\abs{G_m'}\le m^{3/2+\delta}\Bigr\}.
\]
On \(E_m\), Lemma~\ref{lem:grid-speed-detection}, with \(L=m^\delta\), gives
\[
  \Bigl\{\min_{\T}\abs{G_m}\le \frac{\tau}{2\sqrt N}\Bigr\}
  \subseteq
  \Bigl\{\min_{q\in Q_m}\abs{G_m(q)}\le \frac{\tau}{\sqrt N}\Bigr\}.
\]
The miss event is therefore contained in \(E_m^c\), whose probability is
at most \(C_\delta m^{1-2\delta}\) by Lemma~\ref{lem:variance-speed}.  The
uniform-grid assertion is the same elementary mesh estimate as in
Lemma~\ref{lem:grid-speed-detection}.
\end{proof}

\begin{corollary}[Fixed-power variance-only grids]
\label{cor:variance-grid-fixed-power}
Under the hypotheses of Corollary~\ref{cor:variance-grid-detects}, fix
\(\beta>5/2\) and \(\tau>0\).  There are constants
\(C_{\beta,\tau},C_\beta<\infty\) such that, if \(Q_m\) is a uniform grid on
\(\T\) with
\[
  |Q_m|\ge C_{\beta,\tau}N^\beta ,
\]
then
\[
  \Prob\left(
    \Bigl\{\min_{z\in\T}\abs{G_m(z)}\le \frac{\tau}{2\sqrt N}\Bigr\}
    \setminus
    \Bigl\{\min_{q\in Q_m}\abs{G_m(q)}\le \frac{\tau}{\sqrt N}\Bigr\}
  \right)
  \le C_\beta N^{5-2\beta}.
\]
Thus every fixed power strictly above \(N^{5/2}\) gives a tail-free
variance-only detector with an explicit polynomial miss rate.
\end{corollary}

\begin{proof}
Apply Corollary~\ref{cor:variance-grid-detects} with
\(\delta=\beta-2>1/2\).  A uniform grid with \(M\) points has mesh at most
\(C/M\), so \(M\ge C_{\beta,\tau}N^\beta\) implies
\[
  h_m\le \frac{\tau}{2m^{3/2+\delta}\sqrt N}
\]
after increasing \(C_{\beta,\tau}\), because \(N=m+1\).  The miss probability
is at most
\[
  C_\delta m^{1-2\delta}=C_\beta m^{5-2\beta}\le C_\beta N^{5-2\beta}.
\]
\end{proof}

\begin{corollary}[Tail-free polynomial-confidence grids]
\label{cor:variance-grid-polynomial-confidence}
Under the hypotheses of Corollary~\ref{cor:variance-grid-detects}, fix
\(A>0\) and \(\tau>0\).  If \(Q_m\subset\T\) is an \(h_m\)-net in arclength
with
\[
  h_m\le \frac{\tau}{2m^{2+A/2}\sqrt N},
\]
then there is a constant \(C_A<\infty\) such that
\[
  \Prob\left(
    \Bigl\{\min_{z\in\T}\abs{G_m(z)}\le \frac{\tau}{2\sqrt N}\Bigr\}
    \setminus
    \Bigl\{\min_{q\in Q_m}\abs{G_m(q)}\le \frac{\tau}{\sqrt N}\Bigr\}
  \right)
  \le C_A m^{-A}.
\]
Equivalently, a uniform grid of size
\[
  M=O_{A,\tau}\bigl(N^{(A+5)/2}\bigr)
\]
detects continuum \(N^{-1/2}\)-scale excursions with failure probability
\(O_A(m^{-A})\), using only centering and unit variance.
\end{corollary}

\begin{proof}
Apply Corollary~\ref{cor:variance-grid-detects} with
\(\delta=(A+1)/2\).  Then \(1-2\delta=-A\), and
\[
  m^{3/2+\delta}=m^{2+A/2}.
\]
The uniform-grid size follows from the same mesh estimate as in
Corollary~\ref{cor:variance-grid-detects}, using \(N=m+1\).
\end{proof}

\begin{corollary}[Finite-confidence variance-only grids]
\label{cor:variance-grid-confidence}
Let \(\xi_0,\ldots,\xi_m\) be independent mean-zero real random variables with
\(\Exp|\xi_j|^2=1\), let \(G_m(z)=\sum_{j=0}^m\xi_jz^j\), and put
\(N=m+1\).  Fix \(\tau>0\) and \(0<\eta\le1\).  If \(Q_m\subset\T\) is an
\(h_m\)-net in arclength with
\[
  h_m\le c\,\frac{\tau\sqrt\eta}{m^2\sqrt N},
\]
then
\[
  \Prob\left(
    \Bigl\{\min_{z\in\T}\abs{G_m(z)}\le \frac{\tau}{2\sqrt N}\Bigr\}
    \setminus
    \Bigl\{\min_{q\in Q_m}\abs{G_m(q)}\le \frac{\tau}{\sqrt N}\Bigr\}
  \right)
  \le \eta .
\]
Equivalently, a uniform grid of size
\[
  M\ge C_\tau m^2\sqrt N\,\eta^{-1/2}
\]
detects continuum \(N^{-1/2}\)-scale excursions with miss probability at most
\(\eta\), using only centering and unit variance.  The same conclusion holds
for stratified one-point-per-equal-arc grids and for random shifts of a uniform
grid with the same mesh bound.
\end{corollary}

\begin{proof}
By Corollary~\ref{cor:variance-speed-confidence}, with probability at least
\(1-\eta\),
\[
  \sup_{\T}|G_m'|\le R_\eta:=C m^2\eta^{-1/2}.
\]
If \(h_m\le\tau/(2R_\eta\sqrt N)\), then
Lemma~\ref{lem:grid-speed-detection} gives
\[
  \left\{\min_{\T}\abs{G_m}\le \frac{\tau}{2\sqrt N}\right\}
  \cap
  \left\{\sup_{\T}|G_m'|\le R_\eta\right\}
  \subseteq
  \left\{\min_{q\in Q_m}\abs{G_m(q)}\le \frac{\tau}{\sqrt N}\right\}.
\]
Thus the miss event is contained in the derivative-speed exceptional event.
The uniform-grid and stratified-grid assertions are the elementary circle mesh
bound.
\end{proof}

\begin{corollary}[Continuum-scale sampling bracket]
\label{cor:grid-sampling-bracket}
Let \(\xi_0,\ldots,\xi_m\) be independent mean-zero real random variables whose
densities are bounded by \(K\) and satisfy \(\Exp|\xi_j|^2=1\), and put
\(N=m+1\).
\begin{enumerate}[label=\textup{(\roman*)}]
\item Fix \(\gamma>0\), \(T<\infty\), and deterministic finite sets
\(Q_m^-\subset\T\) such that every point of \(Q_m^-\setminus\{\pm1\}\) has
angular distance at least \(\gamma/m\) from \(\{\pm1\}\).  If
\(|Q_m^-|=o(N^2)\), then
\[
  \sup_{0\le\tau\le T}
  \Prob\Bigl\{\min_{q\in Q_m^-}\abs{G_m(q)}
      \le \frac{\tau}{\sqrt N}\Bigr\}\to0 .
\]
The same conclusion holds for independent random grids \(\mathcal Q_m^-\)
satisfying the same endpoint-separation condition almost surely and
\(\Exp|\mathcal Q_m^-|=o(N^2)\).
\item Fix \(A>0\) and \(\tau>0\).  If \(Q_m^+\subset\T\) is an \(h_m\)-net with
\[
  h_m\le \frac{\tau}{2m^{2+A/2}\sqrt N},
\]
then
\[
  \Prob\left(
    \Bigl\{\min_{z\in\T}\abs{G_m(z)}\le \frac{\tau}{2\sqrt N}\Bigr\}
    \setminus
    \Bigl\{\min_{q\in Q_m^+}\abs{G_m(q)}\le \frac{\tau}{\sqrt N}\Bigr\}
  \right)
  \le C_A m^{-A}.
\]
\end{enumerate}
Consequently, sampling at the continuum threshold has a genuine scale
separation: deterministic and independent-random subquadratic grids see only
the pointwise small-ball scale,
whereas any fixed power above \(N^{5/2}\) is complete with a polynomial miss
rate under variance alone, and sufficiently higher polynomial powers give
arbitrary polynomial failure.
\end{corollary}

\begin{proof}
Part (i) is Corollaries~\ref{cor:subquadratic-grid-misses} and
\ref{cor:random-subquadratic-grid-misses}.  Part (ii) is
Corollary~\ref{cor:variance-grid-polynomial-confidence}.  The final sentence
only restates these two estimates at the common threshold \(N^{-1/2}\).
\end{proof}

\begin{corollary}[Boundary law from a polynomial speed tail]
\label{cor:bd-speedtail-to-boundary}
Assume the hypotheses of Corollary~\ref{cor:bd-uniform-speed}.  Suppose in
addition that, for this fixed coefficient law, for every \(A>0\) there are
constants \(L_A,C_A<\infty\) and \(m_0(A)\) such that
\[
  \Prob\Bigl\{\sup_{z\in\T}\abs{G_m'(z)}>L_A m^{3/2}\sqrt{\log m}\Bigr\}
  \le C_A m^{-A}
  \qquad(m\ge m_0(A)).
\]
Then Hypothesis~\ref{hyp:boundary-law} holds for this coefficient law, with
constants depending on \(K\) and on the displayed speed-tail constants.
\end{corollary}

\begin{proof}
Fix \(A,B>0\), let
\[
  E_A:=\Bigl\{\sup_{\T}\abs{G_m'}\le L_A m^{3/2}\sqrt{\log m}\Bigr\},
\]
and put \(L=L_A\sqrt{\log m}\) in
Corollary~\ref{cor:bd-uniform-speed}.  For \(m^{-B}\le t\le1\),
\[
\begin{aligned}
  \Prob\{\min_{\T}\abs{G_m}\le t\}
  &\le
  \Prob\bigl(\{\min_{\T}\abs{G_m}\le t\}\cap E_A\bigr)
  +\Prob(E_A^c) \\
  &\le C_{B,K}L_A\sqrt{m\log m}\,t+C_A m^{-A}.
\end{aligned}
\]
Absorbing \(L_A\), \(C_A\), and the threshold into the constants in
Hypothesis~\ref{hyp:boundary-law} gives the claim.  The conclusion is therefore
uniform over any class of coefficient laws for which the speed-tail constants
are uniform.
\end{proof}

\begin{corollary}[Grid detection from a polynomial speed tail]
\label{cor:speedtail-grid-detects}
Let \(G_m(z)=\sum_{j=0}^m\xi_jz^j\), put \(N=m+1\), and suppose that for a
given coefficient law and a given \(A>0\) there are constants \(L_A,C_A<\infty\)
and \(m_0(A)\) such that
\[
  \Prob\Bigl\{\sup_{z\in\T}\abs{G_m'(z)}>L_A m^{3/2}\sqrt{\log m}\Bigr\}
  \le C_A m^{-A}
  \qquad(m\ge m_0(A)).
\]
Fix \(\tau>0\).  If \(Q_m\subset\T\) is an \(h_m\)-net with
\[
  h_m\le \frac{\tau}{2L_A m^{3/2}\sqrt N\,\sqrt{\log m}},
\]
then, for \(m\ge m_0(A)\),
\[
  \Prob\left(
    \Bigl\{\min_{z\in\T}\abs{G_m(z)}\le \frac{\tau}{2\sqrt N}\Bigr\}
    \setminus
    \Bigl\{\min_{q\in Q_m}\abs{G_m(q)}\le \frac{\tau}{\sqrt N}\Bigr\}
  \right)
  \le C_A m^{-A}.
\]
Thus a uniform grid with \(M\ge C L_A m^{3/2}\sqrt N\sqrt{\log m}/\tau\)
points has the same property.  In particular, the conclusion holds for the
Gaussian law, for uniformly subexponential bounded-density laws, and for
bounded-density laws with a finite \(p\)-moment satisfying \(p>2A+4\), using
the speed-tail inputs proved in Lemma~\ref{lem:derivtail} and
Corollaries~\ref{cor:bd-polymoment-boundary}--\ref{cor:bd-subexp-boundary}.
\end{corollary}

\begin{proof}
Let
\[
  E_A=\Bigl\{\sup_{\T}\abs{G_m'}\le L_A m^{3/2}\sqrt{\log m}\Bigr\}.
\]
On \(E_A\), Lemma~\ref{lem:grid-speed-detection}, with
\(L=L_A\sqrt{\log m}\), gives
\[
  \Bigl\{\min_{\T}\abs{G_m}\le \frac{\tau}{2\sqrt N}\Bigr\}
  \subseteq
  \Bigl\{\min_{q\in Q_m}\abs{G_m(q)}\le \frac{\tau}{\sqrt N}\Bigr\}.
\]
Therefore the miss event in the statement is contained in \(E_A^c\), whose
probability is bounded by the assumed speed tail.  The uniform-grid assertion
is again the elementary mesh estimate on a circle of arclength \(2\pi\).
The listed coefficient classes have the required speed tails by the cited
results, with the polynomial-moment case choosing \(p\) according to the
requested exponent \(A\).
\end{proof}

\begin{corollary}[Natural-speed sampling bracket]
\label{cor:natural-speed-sampling-bracket}
Let \(\xi_0,\ldots,\xi_m\) be independent real random variables whose densities
are bounded by \(K\), let \(G_m(z)=\sum_{j=0}^m\xi_jz^j\), and put
\(N=m+1\).
\begin{enumerate}[label=\textup{(\roman*)}]
\item Fix \(\gamma>0\), \(T<\infty\), and deterministic finite sets
\(Q_m^-\subset\T\) such that every point of \(Q_m^-\setminus\{\pm1\}\) has
angular distance at least \(\gamma/m\) from \(\{\pm1\}\).  If
\(|Q_m^-|=o(N^2)\), then
\[
  \sup_{0\le\tau\le T}
  \Prob\Bigl\{\min_{q\in Q_m^-}\abs{G_m(q)}
      \le \frac{\tau}{\sqrt N}\Bigr\}\to0 .
\]
The same conclusion holds for independent random grids \(\mathcal Q_m^-\)
satisfying the same endpoint-separation condition almost surely and
\(\Exp|\mathcal Q_m^-|=o(N^2)\).
\item Suppose, in addition, that for a given \(A>0\) there are constants
\(L_A,C_A<\infty\) and \(m_0(A)\) such that
\[
  \Prob\Bigl\{\sup_{z\in\T}\abs{G_m'(z)}
      >L_A m^{3/2}\sqrt{\log m}\Bigr\}
  \le C_A m^{-A}
  \qquad(m\ge m_0(A)).
\]
Fix \(\tau>0\).  If \(Q_m^+\subset\T\) is an \(h_m\)-net with
\[
  h_m\le \frac{\tau}{2L_A m^{3/2}\sqrt N\,\sqrt{\log m}},
\]
then, for \(m\ge m_0(A)\),
\[
  \Prob\left(
    \Bigl\{\min_{z\in\T}\abs{G_m(z)}\le \frac{\tau}{2\sqrt N}\Bigr\}
    \setminus
    \Bigl\{\min_{q\in Q_m^+}\abs{G_m(q)}\le \frac{\tau}{\sqrt N}\Bigr\}
  \right)
  \le C_A m^{-A}.
\]
\end{enumerate}
Thus, under a natural polynomial derivative-speed tail, sampling at the
continuum threshold is squeezed between deterministic or independent random
subquadratic grids, which see only the pointwise small-ball scale, and grids of
size
\[
  O_{A,\tau}\bigl(L_A N^2\sqrt{\log N}\bigr),
\]
which detect every continuum excursion up to \(O_A(m^{-A})\) miss probability.
Independent uniform random grids also detect, with the coupon-collector
overhead
\[
  O_{A,\tau,L_A}\bigl(N^2(\log N)^{3/2}\bigr).
\]
Stratified one-point-per-equal-arc grids, including random shifts of a uniform
lattice, detect at the deterministic
\[
  O_{A,\tau}\bigl(L_A N^2\sqrt{\log N}\bigr)
\]
scale without that extra logarithm.
The upper side applies to Gaussian coefficients, uniformly subexponential
bounded-density laws, and bounded-density laws with a finite \(p\)-moment
satisfying \(p>2A+4\).
\end{corollary}

\begin{proof}
Part (i) is Corollaries~\ref{cor:subquadratic-grid-misses} and
\ref{cor:random-subquadratic-grid-misses}.  Part (ii) is
Corollary~\ref{cor:speedtail-grid-detects}.  A uniform grid of
\(M\) points has mesh \(O(1/M)\), so the displayed mesh condition is satisfied
once \(M\ge C_\tau L_A m^{3/2}\sqrt N\sqrt{\log m}\), which is
\(O_\tau(L_A N^2\sqrt{\log N})\).  The independent-uniform statement is
Corollary~\ref{cor:uniform-random-grid-detects}, and the stratified statement
is Corollary~\ref{cor:stratified-grid-detects}.  The listed coefficient classes have the
required speed tails by Lemma~\ref{lem:derivtail} and
Corollaries~\ref{cor:bd-polymoment-boundary}--\ref{cor:bd-subexp-boundary},
choosing \(p\) according to the requested exponent in the polynomial-moment
case.
\end{proof}

\begin{corollary}[Polynomial-moment bounded-density boundary law]
\label{cor:bd-polymoment-boundary}
Let \(\xi_0,\ldots,\xi_m\) be independent, mean-zero, variance-one real random
variables whose densities are bounded by \(K\).  Fix \(A_\ast>0\), and suppose
that for some \(p>2A_\ast+4\),
\[
  M_p:=\sup_{m\ge1}\max_{0\le j\le m}\Exp\abs{\xi_j}^{p}<\infty .
\]
Then there are constants \(L_{A_\ast,p,M_p},C_{A_\ast,p,M_p}<\infty\) and
\(m_0(A_\ast,p,M_p)\) such that
\[
  \Prob\Bigl\{\sup_{z\in\T}\abs{G_m'(z)}
      > L_{A_\ast,p,M_p}m^{3/2}\sqrt{\log m}\Bigr\}
  \le C_{A_\ast,p,M_p}m^{-A_\ast}
  \qquad(m\ge m_0).
\]
Consequently, if the displayed moment bound holds for every \(p<\infty\), then
Hypothesis~\ref{hyp:boundary-law} holds for this coefficient law, with
constants depending on \(K\) and on the moment bounds used for the requested
error exponent.
\end{corollary}

\begin{proof}
By Corollary~\ref{cor:bd-speedtail-to-boundary}, it suffices to prove the
displayed speed tail.  Put
\[
  H_m(\theta):=G_m'(e^{i\theta})
  =\sum_{j=1}^m j\xi_j e^{i(j-1)\theta}.
\]
As in the proof of Corollary~\ref{cor:bd-subexp-boundary}, Bernstein's
inequality for trigonometric polynomials reduces the supremum to a grid
\(\{\theta_\ell\}_{\ell=1}^N\) with \(N\le C m\):
\[
  \sup_\theta\abs{H_m(\theta)}
  \le 2\max_{1\le\ell\le N}\abs{H_m(\theta_\ell)}.
\]
Fix a grid point and write either the real or imaginary part as a sum
\[
  S_\ell=\sum_{j=1}^m a_{j,\ell}\xi_j,
  \qquad \abs{a_{j,\ell}}\le j .
\]
The Fuk--Nagaev inequality \cite{FukNagaev} gives, for \(u>0\),
\[
  \Prob\{\abs{S_\ell}>u\}
  \le
  C_pu^{-p}\sum_{j=1}^m\Exp\abs{a_{j,\ell}\xi_j}^{p}
  +2\exp\left\{-c\frac{u^2}{\sum_{j=1}^m\Exp(a_{j,\ell}^2\xi_j^2)}\right\}.
\]
Since \(\sum_j j^p\le C_pm^{p+1}\) and \(\sum_j j^2\le C m^3\), this implies
\[
  \Prob\{\abs{H_m(\theta_\ell)}>u\}
  \le C_{p,M_p}u^{-p}m^{p+1}
     +4\exp\{-c u^2/m^3\}.
\]
With \(u=Lm^{3/2}\sqrt{\log m}\), the first term is at most
\[
  C_{p,M_p,L}m^{1-p/2}(\log m)^{-p/2},
\]
and the exponential term is at most \(4m^{-cL^2}\).  After the union bound over
the \(O(m)\)-point grid, the polynomial term is \(O(m^{2-p/2})\), which is
\(O(m^{-A_\ast})\) because \(p>2A_\ast+4\).  Choosing \(L\) large enough in
terms of \(A_\ast\) controls the exponential term as well.  This proves the
speed-tail estimate, and the final assertion follows from
Corollary~\ref{cor:bd-speedtail-to-boundary}, choosing \(p\) separately for
each requested error exponent.
\end{proof}

\begin{corollary}[Subexponential bounded-density boundary law]
\label{cor:bd-subexp-boundary}
Let \(\xi_0,\ldots,\xi_m\) be independent, mean-zero, variance-one real random
variables whose densities are bounded by \(K\), and assume in addition that
\(\norm{\xi_j}_{\psi_1}\le\kappa\) uniformly in \(j\) and \(m\).  For every
\(A>0\) there are constants \(L_{A,\kappa},C_{A,\kappa}<\infty\) and
\(m_0(A,\kappa)\) such that
\[
  \Prob\Bigl\{\sup_{z\in\T}\abs{G_m'(z)}
    >L_{A,\kappa}m^{3/2}\sqrt{\log m}\Bigr\}
  \le C_{A,\kappa}m^{-A}
  \qquad(m\ge m_0(A,\kappa)).
\]
Hence Hypothesis~\ref{hyp:boundary-law} holds, with constants depending also
on \(\kappa\).  Consequently, if Hypothesis~\ref{hyp:clearance} holds for this
coefficient law, the conclusion of Theorem~\ref{thm:main} holds.
\end{corollary}

\begin{proof}
We first verify the displayed polynomial speed tail.  Put
\[
  H_m(\theta):=G_m'(e^{i\theta})
  =\sum_{j=1}^m j\xi_j e^{i(j-1)\theta}.
\]
We first reduce the supremum to a grid.  If \(h\) is a trigonometric polynomial
of degree at most \(m\), Bernstein's inequality and a grid
\(\{\theta_\ell\}_{\ell=1}^N\) with \(N\asymp m\) and mesh at most \(1/(4m)\)
give
\[
  \sup_\theta\abs{h(\theta)}
  \le 2\max_{1\le\ell\le N}\abs{h(\theta_\ell)}.
\]
Indeed, at a maximum point choose the nearest grid point and use
\(\norm{h'}_\infty\le m\norm{h}_\infty\).

For each grid point, the real and imaginary parts of \(H_m(\theta_\ell)\) are
centered subexponential linear forms.  The Bernstein inequality for
subexponential sums \cite[Thm.~2.8.1]{VershyninHDP} gives
\[
  \Prob\{\abs{H_m(\theta_\ell)}>u\}
  \le
  4\exp\left\{-c_\kappa
  \min\left(\frac{u^2}{m^3},\frac{u}{m}\right)\right\}.
\]
Here \(\sum_{j=1}^m j^2=O(m^3)\) and \(\max_j j\le m\) are the only coefficient
norms used.  Taking \(u=Lm^{3/2}\sqrt{\log m}\), the quadratic term controls the
minimum for all large \(m\), and therefore
\[
  \Prob\{\abs{H_m(\theta_\ell)}>Lm^{3/2}\sqrt{\log m}\}
  \le 4m^{-c_\kappa L^2}.
\]
A union bound over the \(O(m)\) grid points gives, after choosing \(L=L_A\) large
in terms of \(A\) and \(\kappa\), constants \(L_A,C_A<\infty\) such that
\[
  \Prob\Bigl\{\sup_{z\in\T}\abs{G_m'(z)}
    >L_A m^{3/2}\sqrt{\log m}\Bigr\}
  \le C_A m^{-A}.
\]
Corollary~\ref{cor:bd-speedtail-to-boundary} now gives
Hypothesis~\ref{hyp:boundary-law}.  The final assertion is
Theorem~\ref{thm:main}.
\end{proof}

\begin{remark}[Speed cutoffs and the full boundary law]\label{rem:speed-extra}
Corollary~\ref{cor:bd-uniform-speed} should not be read as an unconditional
bounded-density boundary law under the assumptions of Theorem~\ref{thm:main}.
Even for standard Gaussian coefficients, no fixed constant \(L\) makes
\(\{\sup_{\T}|G_m'|\le Lm^{3/2}\}\) a high-probability event, since
\[
  G_m'(1)=\sum_{j=1}^m j\xi_j,\qquad
  m^{-3/2}G_m'(1)\Longrightarrow N(0,1/3).
\]
Moreover, bounded density and unit variance alone do not give the polynomially
small exceptional probabilities needed to combine a growing speed cutoff with
Hypothesis~\ref{hyp:boundary-law}: heavy-tailed absolutely continuous laws with
bounded density can make the complement of
\(\{\sup_{\T}|G_m'|\le L_m m^{3/2}\}\) only logarithmically small at the natural
Gaussian scale \(L_m\asymp\sqrt{\log m}\).  Removing the cutoff therefore
requires a genuine continuum crossing or local-limit argument, not just a
derivative tail estimate.  Theorem~\ref{thm:main} uses only the weaker
tail-free separation in Corollary~\ref{cor:variance-boundary-separation}: a
polynomially small exceptional probability at a deliberately polynomial speed
cutoff is enough for the single boundary-minimum threshold needed by
finite-section conditioning.
\end{remark}

\subsection{Occupation-measure route to the single-arc bound}

\begin{remark}[The occupation-measure proof and its logarithmic loss]
Fix an arc $I$ of arclength $\abs I=c/m$, $c\ge1$, whose closure is at angular distance $\ge\gamma_0/m$ from $\{\pm1\}$ for a fixed $\gamma_0>0$. Parametrize $I$ by $\theta\in J$, an interval of length $c/m$, and write $g(\theta):=\abs{G_m(e^{i\theta})}$. The event $\{\min_I g\le t\}$ is the event that the planar curve $\theta\mapsto G_m(e^{i\theta})$ enters the disk $B(0,t)$ for some $\theta\in J$.

\emph{Occupation measure.} Let $\Lambda(s):=\int_J\mathbf1\{g(\theta)\le s\}\,d\theta$ be the time the curve spends in $B(0,s)$. By Fubini and Lemma~\ref{lem:onepoint} (the away-from-$\pm1$ branch applies, since $J$ is $\gamma_0/m$ away from $0,\pi$),
\begin{equation}\label{eq:occ}
  \Exp\Lambda(s)=\int_J\Prob\{g(\theta)\le s\}\,d\theta\le\abs J\cdot\frac{C_{\gamma_0}s^2}{m+1}=\frac{C_{\gamma_0}\,c\,s^2}{m(m+1)}.
\end{equation}

\emph{Sojourn lower bound per visit.} If one had the localized speed event
\[
  G_I:=\{\sup_I\abs{G_m'}\le C_2\,m^{3/2}\}
\]
with probability \(1-o(t/\sqrt m)\), then the curve speed would be
\(\bigl|\tfrac{d}{d\theta}G_m(e^{i\theta})\bigr|=\abs{iz\,G_m'(z)}\le C_2m^{3/2}\).
Hence whenever $g(\theta_0)\le t$, the curve would stay in $B(0,2t)$ for all $\theta$ with $\abs{\theta-\theta_0}\le\delta$, where $\delta:=\tfrac{t}{C_2 m^{3/2}}$, because $g$ changes by at most $\delta\,C_2 m^{3/2}=t$. Since $\abs J=c/m\ge\delta$ for $t\le1$ and $m$ large, each visit of $\{g\le t\}$ would contribute at least $\delta$ to $\Lambda(2t)$, so on $G_I$
\[
  \mathbf1\Bigl\{\min_J g\le t\Bigr\}\le \frac{\Lambda(2t)}{\delta}.
\]

\emph{Combine.} Taking expectations and using \eqref{eq:occ} with $s=2t$ would give
\[
  \Prob\Bigl\{\min_I g\le t\Bigr\}
  \le\Prob(G_I^c)+\frac{1}{\delta}\,\Exp\Lambda(2t)
  \le\Prob(G_I^c)+\frac{C_2 m^{3/2}}{t}\cdot\frac{4C_{\gamma_0}\,c\,t^2}{m(m+1)}
  \le \frac{C\,c\,t}{\sqrt m}.
\]
This is the desired rate \eqref{eq:singlearc}.  The obstruction to proving it by
this elementary occupation route is the speed/level-crossing step: the available
global Salem--Zygmund tail in Lemma~\ref{lem:derivtail} carries the logarithmic
loss recorded in Proposition~\ref{prop:singlearc-log}.  Lemma~\ref{lem:kr-crossing}
removes that loss in the Gaussian case away from \(z=\pm1\) by counting actual
level crossings instead of enforcing a deterministic speed cutoff.
\end{remark}

\begin{proposition}[Exact Gaussian root-grid law]\label{prop:gaussian-root-grid}
Let \(N=m+1\ge3\), let \(\omega=\exp(2\pi i/N)\), and let
\[
  G_m(z)=\sum_{j=0}^{N-1}\xi_jz^j
\]
with \(\xi_j\) independent standard real Gaussians.  Put
\[
  \mathcal R_N:=\{\omega^k:1\le k\le \lfloor(N-1)/2\rfloor\}.
\]
Then the random variables \(N^{-1/2}G_m(\omega^k)\), \(1\le k\le
\lfloor(N-1)/2\rfloor\), are independent standard complex Gaussians
\((X+iY)/\sqrt2\), with \(X,Y\) independent \(N(0,1)\).  Consequently, for
every \(t\ge0\),
\begin{equation}\label{eq:gaussian-root-grid}
  \Prob\Bigl\{\min_{z\in\mathcal R_N}\abs{G_m(z)}\le t\Bigr\}
  =1-\exp\!\left(-\frac{\lfloor(N-1)/2\rfloor}{N}\,t^2\right).
\end{equation}
In particular, for \(N\ge3\) and \(0\le t\le1\), this probability is
\(\asymp t^2\), with absolute constants.  If
\(\mathcal U_N:=\{\omega^k:0\le k<N\}\) is the full \(N\)-th-root grid, set
\[
  s_N:=1+\mathbf1_{\{2\mid N\}},
  \qquad
  \alpha(x):=\Prob\{|Z|\le x\},\quad Z\sim N(0,1).
\]
Then the full grid has the exact law
\begin{equation}\label{eq:gaussian-full-root-grid-exact}
  \Prob\Bigl\{\min_{z\in\mathcal U_N}\abs{G_m(z)}>t\Bigr\}
  =
  \bigl(1-\alpha(t/\sqrt N)\bigr)^{s_N}
  \exp\!\left(-\frac{\lfloor(N-1)/2\rfloor}{N}\,t^2\right).
\end{equation}
Consequently,
\begin{equation}\label{eq:gaussian-full-root-grid}
  \Prob\Bigl\{\min_{z\in\mathcal U_N}\abs{G_m(z)}\le t\Bigr\}
  \le C\left(\frac{t}{\sqrt N}+t^2\right)
  \qquad(0\le t\le1),
\end{equation}
with an absolute constant \(C\).
\end{proposition}

\begin{proof}
Let \(M=\lfloor(N-1)/2\rfloor\).  For \(1\le k\le M\), set
\[
  A_k=\sqrt{\frac2N}\sum_{j=0}^{N-1}\xi_j\cos(2\pi kj/N),
  \qquad
  B_k=\sqrt{\frac2N}\sum_{j=0}^{N-1}\xi_j\sin(2\pi kj/N).
\]
The vectors
\[
  \sqrt{2/N}\,(\cos(2\pi kj/N))_{j=0}^{N-1},\qquad
  \sqrt{2/N}\,(\sin(2\pi kj/N))_{j=0}^{N-1},
  \qquad 1\le k\le M,
\]
are orthonormal in \(\R^N\).  Since \((\xi_0,\dots,\xi_{N-1})\) is a standard
Gaussian vector in \(\R^N\), the coordinates \(A_k,B_k\) are independent
standard real Gaussians.  Moreover
\[
  N^{-1/2}G_m(\omega^k)=\frac{A_k+iB_k}{\sqrt2}.
\]
Thus the displayed variables are independent standard complex Gaussians, and
\(\abs{N^{-1/2}G_m(\omega^k)}^2\) are independent exponential variables of
mean \(1\).  This gives \eqref{eq:gaussian-root-grid}.

For the full root grid, the conjugate roots do not add new moduli:
\(G_m(\omega^{N-k})=\overline{G_m(\omega^k)}\).  The only additional points are
the real DFT coordinates \(1\), and also \(-1\) when \(N\) is even.  The
corresponding normalized real coordinates are
\[
  C_0=N^{-1/2}\sum_{j=0}^{N-1}\xi_j,\qquad
  C_{N/2}=N^{-1/2}\sum_{j=0}^{N-1}(-1)^j\xi_j\quad(N\ {\rm even}).
\]
Together with the cosine and sine vectors above, the relevant real vectors form
an orthonormal family in \(\R^N\).  Hence these \(s_N\) real \(N(0,1)\)
coordinates are independent of the non-real complex coordinates.  The event
\(\min_{\mathcal U_N}\abs{G_m}>t\) is therefore the intersection of the
non-real event from \eqref{eq:gaussian-root-grid} with the \(s_N\) independent
events \(\{|C|>t/\sqrt N\}\), proving
\eqref{eq:gaussian-full-root-grid-exact}.  The upper bound
\eqref{eq:gaussian-full-root-grid} follows from
\(\alpha(x)\le Cx\), \(1-e^{-x}\le x\), and a union bound applied to the exact
survival formula.
\end{proof}

\begin{corollary}[The DFT root-grid minimum has pointwise scale]\label{cor:gaussian-root-grid-rayleigh}
Under the hypotheses and notation of Proposition~\ref{prop:gaussian-root-grid},
let
\[
  M_N^{\rm nr}:=\min_{z\in\mathcal R_N}\abs{G_m(z)},
  \qquad
  M_N^{\rm full}:=\min_{z\in\mathcal U_N}\abs{G_m(z)}.
\]
For every fixed \(T<\infty\) there is \(C_T<\infty\) such that
\[
  \sup_{0\le t\le T}
  \left|\Prob\{M_N^{\rm nr}>t\}-e^{-t^2/2}\right|
  \le \frac{C_T}{N}
\]
and
\[
  \sup_{0\le t\le T}
  \left|\Prob\{M_N^{\rm full}>t\}-e^{-t^2/2}\right|
  \le \frac{C_T}{\sqrt N}.
\]
Thus both root-grid minima converge in distribution to the Rayleigh law with
survival function \(e^{-t^2/2}\).  In particular, the DFT grid minimum is
typically of constant pointwise size, not of the continuum scale \(N^{-1/2}\).
Moreover, for every fixed \(p>0\) and \(N\ge3\), with
\(M=\lfloor(N-1)/2\rfloor\), \(s_N=1+\mathbf1_{\{2\mid N\}}\), and
\(R_p:=2^{p/2}\Gamma(1+p/2)\),
\[
  \Exp (M_N^{\rm nr})^p
  =\Gamma\!\left(1+\frac p2\right)\left(\frac{N}{M}\right)^{p/2}
  =R_p+\frac{s_Np}{2N}R_p+O_p(N^{-2}),
\]
and
\[
\begin{aligned}
  \Exp (M_N^{\rm full})^p
  &=R_p
   -\frac{s_Np\,2^{p/2}\Gamma\!\left((p+1)/2\right)}{\sqrt{\pi N}}
   \\
  &\quad
   +\frac{pR_p}{N}\left(\frac{s_N}{2}+\frac{s_N(s_N-1)}{\pi}\right)
   +O_p(N^{-3/2}).
\end{aligned}
\]
In particular,
\[
  \Exp M_N^{\rm nr}=\frac{\sqrt{\pi N}}{2\sqrt M}
  =\sqrt{\frac\pi2}+\frac{s_N}{2N}\sqrt{\frac\pi2}+O(N^{-2}),
\]
and
\[
  \Exp M_N^{\rm full}=\sqrt{\frac\pi2}
  -s_N\sqrt{\frac{2}{\pi N}}
  +\frac1N\sqrt{\frac\pi2}
     \left(\frac{s_N}{2}+\frac{s_N(s_N-1)}{\pi}\right)
  +O(N^{-3/2}).
\]
\end{corollary}

\begin{proof}
Put \(M=\lfloor(N-1)/2\rfloor\).  From \eqref{eq:gaussian-root-grid},
\[
  \Prob\{M_N^{\rm nr}>t\}=\exp\!\left(-\frac{M}{N}t^2\right),
  \qquad
  \left|\frac{M}{N}-\frac12\right|\le \frac1N.
\]
The first bound follows uniformly for \(0\le t\le T\) by the mean-value
theorem applied to \(u\mapsto e^{-ut^2}\).  For the full grid, use
\eqref{eq:gaussian-full-root-grid-exact}.  Uniformly for \(0\le t\le T\),
\[
  0\le \alpha(t/\sqrt N)\le C_TN^{-1/2},
  \qquad
  \left|(1-\alpha(t/\sqrt N))^{s_N}-1\right|\le C_TN^{-1/2},
\]
while the non-real exponential factor differs from \(e^{-t^2/2}\) by
\(O_T(N^{-1})\).  Multiplying the two factors gives the second bound.

For \(p>0\), the non-real moment follows by integrating the exact survival
function:
\[
  \Exp (M_N^{\rm nr})^p
  =p\int_0^\infty t^{p-1}\exp\!\left(-\frac{M}{N}t^2\right)\,dt
  =\Gamma\!\left(1+\frac p2\right)\left(\frac{N}{M}\right)^{p/2}.
\]
Since
\[
  \frac{M}{N}=\frac12-\frac{s_N}{2N},
\]
this gives the displayed non-real expansion.
For the full grid, use the exact survival formula
\eqref{eq:gaussian-full-root-grid-exact}.  Split the moment integral at
\(t=\sqrt N\); the tail \(t\ge\sqrt N\) is exponentially small.  On
\(0\le t\le\sqrt N\), with \(c_0:=\sqrt{2/\pi}\),
\[
  \alpha(t/\sqrt N)=c_0\frac{t}{\sqrt N}
  +O\!\left(\frac{t^3}{N^{3/2}}\right),
\]
and
\[
  (1-\alpha(t/\sqrt N))^{s_N}
  =
  1-s_Nc_0\frac{t}{\sqrt N}
  +\binom{s_N}{2}c_0^2\frac{t^2}{N}
  +O\!\left(\frac{t^3+t^4}{N^{3/2}}\right).
\]
Also,
\[
  e^{-Mt^2/N}
  =e^{-t^2/2}\left(1+\frac{s_Nt^2}{2N}
  +O\!\left(\frac{t^4}{N^2}\right)\right).
\]
Multiplying and integrating against \(pt^{p-1}\,dt\) gives
\[
\begin{aligned}
  \Exp (M_N^{\rm full})^p
  &=p\int_0^\infty t^{p-1}e^{-t^2/2}\,dt  \\
  &\quad
  -\frac{ps_Nc_0}{\sqrt N}\int_0^\infty t^p e^{-t^2/2}\,dt \\
  &\quad
  +\frac{p}{N}\left(\frac{s_N}{2}+\frac{s_N(s_N-1)}{\pi}\right)
    \int_0^\infty t^{p+1}e^{-t^2/2}\,dt
  +O_p(N^{-3/2}).
\end{aligned}
\]
The three displayed integrals are
\[
  \frac{R_p}{p},\qquad
  2^{(p-1)/2}\Gamma\!\left(\frac{p+1}{2}\right),\qquad
  R_p,
\]
respectively.  Substituting them gives the full-grid moment expansion.  The
expectation statement is the case \(p=1\).
\end{proof}

\begin{corollary}[The DFT grid misses the continuum scale]\label{cor:gaussian-root-grid-misses}
Under the hypotheses and notation of Proposition~\ref{prop:gaussian-root-grid},
for every fixed \(T<\infty\) there is a constant \(C_T<\infty\) such that
\[
  \sup_{0\le \tau\le T}
  \left|
  N\Prob\Bigl\{\min_{z\in\mathcal U_N}\abs{G_m(z)}
       \le \frac{\tau}{\sqrt N}\Bigr\}
  -
  \left(s_N\sqrt{\frac2\pi}\,\tau+\frac{\tau^2}{2}\right)
  \right|
  \le \frac{C_T}{N}.
\]
In particular, along odd \(N\) the limiting coefficient is
\(\sqrt{2/\pi}\,\tau+\tau^2/2\), while along even \(N\) it is
\(2\sqrt{2/\pi}\,\tau+\tau^2/2\).  Thus
for \(0\le\tau\le T\),
\[
  \Prob\Bigl\{\min_{z\in\mathcal U_N}\abs{G_m(z)}
       \le \frac{\tau}{\sqrt N}\Bigr\}
  \le \frac{C_T}{N}.
\]
Thus the natural \(N\)-point DFT grid has vanishing detection probability at
the continuum boundary-minimum scale \(t\asymp N^{-1/2}\).
\end{corollary}

\begin{proof}
Put \(t=\tau/\sqrt N\) in the exact formula
\eqref{eq:gaussian-full-root-grid-exact}.  Set
\[
  a_N(\tau):=\alpha(\tau/N),
  \qquad
  b_N(\tau):=\frac{\lfloor(N-1)/2\rfloor}{N^2}\,\tau^2 .
\]
Uniformly for \(0\le\tau\le T\),
\[
  a_N(\tau)=\sqrt{\frac2\pi}\frac{\tau}{N}+O_T(N^{-3}),
  \qquad
  b_N(\tau)=\frac{\tau^2}{2N}+O_T(N^{-2}).
\]
Since \(s_N\le2\) and \(a_N+b_N=O_T(N^{-1})\), Taylor expansion gives
\begin{align*}
  \Prob\Bigl\{\min_{z\in\mathcal U_N}\abs{G_m(z)}\le
       \frac{\tau}{\sqrt N}\Bigr\}
  &=1-(1-a_N(\tau))^{s_N}e^{-b_N(\tau)}\\
  &=s_Na_N(\tau)+b_N(\tau)+O_T(N^{-2}).
\end{align*}
Multiplying by \(N\) proves the displayed estimate, and the final \(O_T(N^{-1})\)
bound follows by increasing \(C_T\).
\end{proof}

\begin{remark}[The $t^2$ point bound is a grid-resolution artifact]\label{rem:grid-artifact}
The excursion of $\{g\le t\}$ predicted above has angular length of order $\delta\asymp t\,m^{-3/2}$: at curve speed $\asymp m^{3/2}$, dropping the modulus below $t$ and recovering takes only $\Theta(t\,m^{-3/2})$ radians. A numerical scan of $\T$ on a grid of spacing $h$ therefore \emph{misses} such an excursion unless $h\ll t\,m^{-3/2}$; a coarse grid sees only the pointwise event $\{g(z_\ell)\le t\}$, of probability $\asymp t^2/m$ per non-real node. Proposition~\ref{prop:gaussian-root-grid} makes this exact on the \(N=m+1\) root grid: Corollary~\ref{cor:gaussian-root-grid-rayleigh} shows that the root-grid minimum has a non-degenerate Rayleigh limit at fixed \(t\), while the non-real small-\(t\) grid minimum has probability \(\asymp t^2\), and the real endpoint coordinates multiply the survival probability by the exact one-dimensional factors in \eqref{eq:gaussian-full-root-grid-exact}. Corollary~\ref{cor:gaussian-root-grid-misses} shows in particular that at the continuum scale \(t=\tau/\sqrt N\), fixed \(\tau\), this root grid sees the event with the exact first-order coefficient \((s_N\sqrt{2/\pi}\,\tau+\tau^2/2)/N\), up to \(O_\tau(N^{-2})\). The sharp continuum law is $\sqrt m\,t$ (Theorem~\ref{thm:uniform}): the $t^2$ pointwise rate is a sampling artifact, not the true small-ball exponent. We record this so the linear law is not re-broken by under-resolved numerics; a simulation must resolve the scale $t\,m^{-3/2}$ to see it.
\end{remark}

\subsection{From single arcs to the uniform bound}\label{sub:secondmoment}

\begin{remark}[From arc bounds to the circle]
The proof of Theorem~\ref{thm:uniform} is the bookkeeping suggested by the
single-arc scale.  Partition \(\T\) into \(K=\lfloor c_1m\rfloor\) arcs
\(I_1,\ldots,I_K\) of arclength \(\Theta(1/m)\).  After enlarging the two
endpoint arcs by a fixed number of neighbors, every remaining arc lies at
distance \(\ge\gamma_0/m\) from \(\{\pm1\}\), and
Theorem~\ref{thm:singlearc} gives
\[
  \Prob\{\min_{I_k}\abs{G_m}\le t\}\le C\frac{t}{\sqrt m}.
\]
A union bound over \(O(m)\) such arcs returns the whole-circle scale
\(C\sqrt m\,t\).  The endpoint arcs are no longer a conditional input:
Proposition~\ref{prop:endpoint-singlearc} proves their contribution is
\(O(t/\sqrt m)\), and Proposition~\ref{prop:endpoint-lower} shows this local
order is sharp.  Corollary~\ref{cor:endpoint-global-floor} records the resulting
global endpoint floor.  Thus the proved upper bound has the correct local
one-dimensional rates at both the non-real arcs and the two parity endpoints;
what remains for the sharp two-sided global law is the second-moment step that
turns the \(O(m)\) non-real local crossing opportunities into a lower-tail
probability and an eventual limiting law.
\end{remark}

\begin{remark}[Decorrelation and the matching lower bound]\label{rem:decorrelation}
The matching lower bound, and with it the exponential law, should follow from a second-moment argument in the spirit of the good-family estimate of \cite{PhaseRetrieval}: with $Y_t=\#\{k:\min_{I_k}\abs{G_m}\le t\}$ one needs
\begin{equation}\label{eq:pair-decorrelation}
  \Exp[\mathbf1\{\min_{I_k}\abs{G_m}\le t\}\,\mathbf1\{\min_{I_{k'}}\abs{G_m}\le t\}]
  \le (1+o(1))\,\Prob\{\min_{I_k}\abs{G_m}\le t\}\,\Prob\{\min_{I_{k'}}\abs{G_m}\le t\}
\end{equation}
for arcs at angular separation $\gg1/m$, the analogue of the two-slab estimate \cite[Lem.~4.3]{PhaseRetrieval}; the off-diagonal Dirichlet mass \eqref{eq:cov-again} between two separated arcs is $o(1)$ after normalization, so the joint field on \(I_k\cup I_{k'}\) should factorize up to \(o(1)\). Proposition~\ref{prop:bulk-secondmoment-transfer} records the corresponding count-level implication: once the factorial second moment of the bulk down-crossing count is controlled, the fixed-bulk lower tail follows immediately. We do not need this lower bound for the conditioning theorem, which uses only an upper bound: Hypothesis~\ref{hyp:boundary-law} for general laws and Corollary~\ref{prop:uniform-polylog} for Gaussian polynomial scales. The full two-sided law is the target of the minimum-modulus literature cited above.
\end{remark}

\section{From the boundary minimum to the finite section}\label{sec:main}

We now turn the boundary control into a bound on $\smin(T_n)$. The mechanism \eqref{eq:fs-reduction} is made precise as follows.

\begin{proposition}[Finite-section reduction]\label{prop:fs}
Let $P(z)=z^{-r}G_m(z)$ be a banded symbol with $G_m(\zeta)\ne0$ for all $\zeta\in\T$, and assume \(\wind(P)=0\); equivalently, in the balanced case, \(\nu=r\). Set
\[
  \mu:=\min_{z\in\T}\abs{G_m(z)}>0,
  \qquad
  \rho:=\min\bigl\{\bigl|\,\abs\zeta-1\,\bigr|:G_m(\zeta)=0\bigr\}>0,
  \qquad
  M:=\max_j\abs{\xi_j}.
\]
There is a constant $c_2=c_2(m)$, polynomial in $m$, and a threshold $n_0=n_0(m,\rho)$ such that for all $n\ge n_0$
\begin{equation}\label{eq:fs-precise}
  \smin\bigl(T_n(P)\bigr)\;\ge\;c_2(m)^{-1}\,\mu\,\rho^{\,r}.
\end{equation}
Consequently $\kdn(T_n)=\sigma_1(T_n)\smin(T_n)^{-1}\le c_2(m)\,(m+1)\,M\,\mu^{-1}\rho^{-r}$, since $\sigma_1(T_n)\le\norm{T_n}_\infty\le(m+1)M$.
\end{proposition}

\begin{proof}[Proof sketch]
This is the quantitative form of the Gohberg--Feldman inversion for banded symbols, e.g.\ \cite[Ch.~3 and Thm.~3.7]{BG}. Since $G_m$ has no zeros on $\T$ and \(\wind(P)=0\), $P$ admits a Wiener--Hopf factorization $P=P_-P_+$ with $P_\pm^{\pm1}\in H^\infty$ of the two half-planes; the finite-section inverses converge to $T(P^{-1})$ with an error that decays like $\rho^{\,n}$ once $n\ge n_0$. The norm of the limiting inverse is controlled by $\norm{P_+^{-1}}\norm{P_-^{-1}}$, and for a degree-$m$ polynomial symbol these factors are bounded by the reciprocal boundary minimum $\mu^{-1}$ times a factor accounting for the location of the zeros, which contributes $\rho^{-r}$ from the $r$ zeros nearest the circle from inside. Collecting the polynomial-in-$m$ prefactors into $c_2(m)$ gives \eqref{eq:fs-precise}. The norm bound $\sigma_1(T_n)\le(m+1)M$ is the row-sum bound.
\end{proof}

\begin{corollary}[Uniform threshold under a clearance floor]\label{prop:fs-uniform}
Fix \(m\) and \(0<\eta<1\).  There is a deterministic threshold
\(N_0(m,\eta)<\infty\) such that every degree-\(m\) balanced symbol \(P=z^{-r}G_m\)
with \(G_m\ne0\) on \(\T\), \(\wind(P)=0\), boundary minimum \(\mu>0\), and
circle clearance \(\rho\ge\eta\) satisfies
\[
  \smin\bigl(T_n(P)\bigr)\ge c_2(m)^{-1}\,\mu\,\rho^{\,r}
  \qquad(n\ge N_0(m,\eta)),
\]
with the same polynomial-in-\(m\) constant \(c_2(m)\) as in
Proposition~\ref{prop:fs}.  Consequently, after fixing a deterministic
clearance floor \(\eta_m\), the finite-section part of the argument is uniform
over all zero-winding symbols with \(\rho\ge\eta_m\).
\end{corollary}

\begin{proof}
The proof of Proposition~\ref{prop:fs} is deterministic.  If \(\rho\ge\eta\),
then the zeros of \(G_m\) avoid the annulus
\(\{1-\eta/2\le |z|\le1+\eta/2\}\).  The Wiener--Hopf factors and their inverses
are therefore rational functions whose poles stay a distance \(\gtrsim\eta\)
from \(\T\).  For fixed \(m\) this gives a uniform geometric tail bound for the
finite-section error, of the form
\[
  \norm{T_n(P)^{-1}-T(P)^{-1}_{[n]}}
  \le C(m,\eta)\,\mu^{-1}\rho^{-r}q(\eta)^n,
  \qquad 0<q(\eta)<1,
\]
where \(T(P)^{-1}_{[n]}\) denotes the \(n\times n\) corner of the limiting
inverse.  Choosing \(N_0(m,\eta)\) so that this tail is at most one half of the
limiting inverse bound gives the stated estimate.  The coefficient size is not
part of the threshold; it enters only later in the condition-number upper bound
through \(\sigma_1(T_n)\le(m+1)M\).
\end{proof}

\begin{proof}[Proof of Theorem~\ref{thm:main}]
Work in the balanced regime $r=s$, $m=2r$. We prove the result on the zero-winding event \(E0:=\{\nu=r\}\), intersected with three analytic events whose complements have probability \(o(1)\).
\begin{enumerate}[label=(E\arabic*)]
\item \textbf{Clearance.} Let \(E1:=\mathcal C_m\) be the event from Hypothesis~\ref{hyp:clearance}.  Its complement has probability \(o(1)\), and on \(E0\cap E1\),
\[
  \rho\ge\frac{1}{C_0\,m^2(\log m)^3},
\]
so \(G_m\) has no zeros on \(\T\).  Put
\[
  \eta_m:=\frac{1}{C_0m^2(\log m)^3}.
\]
By Corollary~\ref{prop:fs-uniform}, the finite-section estimate
\eqref{eq:fs-precise} is available uniformly on \(E0\cap E1\) for all
\(n\ge N_0(m,\eta_m)\).
\item \textbf{Boundary minimum.} Choose \(A_0>2\) and set
\(E2:=\{\mu\ge m^{-A_0}\}\).  Corollary~\ref{cor:variance-boundary-separation}
gives \(\Prob(E2)=1-o(1)\).  This is the only boundary-minimum input needed by
the finite-section reduction.
\item \textbf{Coefficient size.} Set \(E3:=\{M=\max_j\abs{\xi_j}\le m\}\). Since \(\Exp\abs{\xi_j}^2=1\), Chebyshev and a union bound give \(\Prob(E3)=1-o(1)\), and \(M\le m\) is more than sufficient.
\end{enumerate}
Set \(\Omega_m:=E1\cap E2\cap E3\). Then \(\Prob(\Omega_m)=1-o(1)\), and by the union bound
\[
  \Prob(E0\cap\Omega_m)=\Prob(E0)-o(1).
\]
On \(E0\cap\Omega_m\), Corollary~\ref{prop:fs-uniform} gives, for
$n\ge N_0(m,\eta_m)$,
\[
  \smin(T_n)\ge c_2(m)^{-1}\,\mu\,\rho^{\,r}
  \ge c_2(m)^{-1}\,m^{-A_0}\,\Bigl(\frac{1}{C_0m^2(\log m)^3}\Bigr)^{m/2}.
\]
On \(\Omega_m\), the lower bound is the deterministic quantity
\[
  \delta_m:=c_2(m)^{-1}m^{-A_0}\bigl(C_0m^2(\log m)^3\bigr)^{-m/2}>0,
\]
and the condition-number estimate in Proposition~\ref{prop:fs} gives the deterministic upper bound
\[
  \kdn(T_n)\le
  \Gamma_m:=c_2(m)(m+1)m\,m^{A_0}\bigl(C_0m^2(\log m)^3\bigr)^{m/2}.
\]
Given the prescribed exponent \(b>0\), take \(n_1(m,b)\) so that
\(n^{-b}\le\delta_m\) and \(n^b\ge\Gamma_m\) for all \(n\ge n_1(m,b)\).  Then
\[
  n_{0,b}(m):=\max\{N_0(m,\eta_m),n_1(m,b)\}
\]
is deterministic and gives \eqref{eq:main}.
\end{proof}

\begin{proof}[Proof of Corollary~\ref{cor:quant-main}]
Repeat the proof of Theorem~\ref{thm:main}, but keep the losses explicit.  Put
\[
  E1:=\mathcal C_m,\qquad
  E2:=\{\mu\ge m^{-A_0}\},\qquad
  E3:=\{M=\max_j|\xi_j|\le R_m\}.
\]
The assumed clearance estimate gives \(\Prob(E1^c)\le q_m\).  By
Corollary~\ref{cor:variance-boundary-separation},
\[
  \Prob(E2^c)\le C_{A_0,K}m^{-2(A_0-1)/3}.
\]
Finally, Chebyshev's inequality and the union bound give
\[
  \Prob(E3^c)
  \le\sum_{j=0}^m\Prob\{|\xi_j|>R_m\}
  \le\frac{m+1}{R_m^2}.
\]
Let \(\Omega_{m,A_0,R}:=E1\cap E2\cap E3\).  The displayed probability bound for
\(\Omega_{m,A_0,R}^c\) follows by the union bound, and intersecting with
\(\{\nu=r\}\) costs no more.

On \(\Omega_{m,A_0,R}\cap\{\nu=r\}\), set
\[
  \eta_m=(C_0m^2(\log m)^3)^{-1}.
\]
Corollary~\ref{prop:fs-uniform} applies for all \(n\ge N_0(m,\eta_m)\).  Hence
\[
  \smin(T_n)\ge
  c_2(m)^{-1}m^{-A_0}\bigl(C_0m^2(\log m)^3\bigr)^{-m/2}
  =:\delta_{m,A_0}>0,
\]
and the condition number is at most
\[
  \Gamma_{m,A_0,R}:=
  c_2(m)(m+1)R_m\,m^{A_0}\bigl(C_0m^2(\log m)^3\bigr)^{m/2}.
\]
Choose \(n_{0,b,A_0,R}(m)\ge N_0(m,\eta_m)\) so large that
\[
  n_{0,b,A_0,R}(m)^{-b}\le\delta_{m,A_0},\qquad
  n_{0,b,A_0,R}(m)^b\ge\Gamma_{m,A_0,R}.
\]
Then the same inequalities hold for every \(n\ge n_{0,b,A_0,R}(m)\).  This
deterministic threshold gives the claimed singular-value and condition-number
bounds.
\end{proof}

\begin{proof}[Proof of Theorem~\ref{thm:gaussian-main}]
For Gaussian coefficients, Hypothesis~\ref{hyp:clearance} is supplied by the
Konyagin--Schlag clearance theorem \cite{KMS}.  The conclusion is therefore
Theorem~\ref{thm:main}.  The sharper Gaussian boundary estimate
Corollary~\ref{prop:uniform-polylog} remains useful for the full boundary-law
discussion, but is no longer needed for finite-section conditioning.
\end{proof}

\begin{remark}[The honest content of the bound]\label{rem:honest-bound}
The bound \eqref{eq:main} is a statement about the $n$-asymptotics at fixed bandwidth $m$ on the stable subevent: for any prescribed polynomial exponent \(b>0\), the balanced finite sections do \emph{not} suffer the exponential-in-$n$ blow-up of the unbalanced regime once the winding number is zero and the symbol stays off $\T$, after increasing the deterministic threshold \(n_{0,b}(m)\). The theorem deliberately makes no high-probability claim for stability itself, which is constrained by the parity obstruction \eqref{eq:half}. The prefactors are exponentially small in $m$, reflecting the rare-but-possible near-degeneracies of the boundary minimum and the circle clearance; we make no claim that the $m$-dependence is sharp.
\end{remark}

\begin{remark}[Single arc versus the whole circle]\label{rem:onesided}
The single-arc bound (Theorem~\ref{thm:singlearc}) controls \(\abs{G_m}\) only on one arc \(I_0\), not on all of \(\T\), and so by itself does not give \eqref{eq:fs-precise}, which needs the global minimum \(\mu\). Its fixed-arc consequence would be a lower bound on the symbol restricted to a favorable arc, equivalently a localized resolvent estimate; promoting it to the full \(\smin(T_n)\) is exactly the role of the uniform bound \eqref{eq:uniform}.
\end{remark}

\section{The bridge to the parity obstruction}\label{sec:bridge}

We make precise the moral asserted in Section~\ref{sub:bridge}: the parity observable of \cite{Parity} and the boundary small-ball of this paper are two readings of the same boundary evaluations. For Rademacher coefficients in the balanced convention \(m=2r\), the exact zero story is subtler than in adjacent parities because \(m+1\) is odd.

\subsection{The parity observable is a boundary evaluation}
Recall from \cite[Lem.~2.1]{Parity} that, since $G_m$ has real coefficients, its non-real zeros come in conjugate pairs of equal modulus, so
\begin{equation}\label{eq:parity}
  \nu\equiv\#\{\text{real zeros of }G_m\text{ in }(-1,1)\}\equiv\mathbf1\{G_m(1)G_m(-1)<0\}\pmod2,
\end{equation}
the last congruence holding when $G_m(\pm1)\ne0$. The two numbers $G_m(1)=\sum_j\xi_j$ and $G_m(-1)=\sum_j(-1)^j\xi_j$ are precisely the evaluations of the boundary field $G_m|_\T$ at the two real points $z=\pm1$ --- the same two points at which, in Lemma~\ref{lem:onepoint}, the field collapses from two real dimensions to one and the small-ball exponent degrades from $t^2$ to $t$. Thus the parity obstruction lives at the two real boundary points, and the boundary small-ball estimate is its quantitative refinement: \eqref{eq:parity} reads only the \emph{sign} of the boundary field at $\pm1$, while $\min_\T\abs{G_m}$ reads its \emph{magnitude} along all of $\T$.

\subsection{The Bernoulli lattice obstruction}
\begin{proposition}[Endpoint lattice scale for Rademacher coefficients]\label{prop:rademacher-endpoint}
Let \(\varepsilon_0,\ldots,\varepsilon_m\) be independent Rademacher signs and
let \(G_m(z)=\sum_{j=0}^m\varepsilon_jz^j\).  Put \(N=m+1\).  There is an
absolute constant \(C<\infty\) such that, for each \(\zeta\in\{1,-1\}\), each
\(a\in\mathbb R\), and each \(t\ge0\),
\[
  \Prob\{\abs{G_m(\zeta)-a}\le t\}\le C\,\frac{1+t}{\sqrt N}.
\]
In particular,
\[
  \Prob\{\abs{G_m(\zeta)}\le t\}\le C\,\frac{1+t}{\sqrt N}.
\]
Moreover, if \(N\) is even then
\[
  \Prob\{G_m(\zeta)=0\}\asymp N^{-1/2},
\]
whereas if \(N\) is odd then
\[
  \Prob\{\abs{G_m(\zeta)}=1\}\asymp N^{-1/2}.
\]
Thus the endpoint small-ball scale for lattice coefficients is governed by the
lattice spacing; a density-type bound \(Ct/\sqrt N\), uniformly in the center,
cannot hold below that spacing.
\end{proposition}

\begin{proof}
For \(\zeta=1\), \(G_m(1)=\sum_{j=0}^m\varepsilon_j\).  For \(\zeta=-1\), the
signs \((-1)^j\varepsilon_j\) are again independent Rademacher signs, so
\(G_m(-1)\) has the same law.  It is therefore enough to consider
\(S_N=\sum_{j=1}^N\varepsilon_j\).  The variable \(S_N\) is supported on the
lattice \(\{k\in\mathbb Z:k\equiv N \pmod 2\}\), and
\[
  \Prob\{S_N=k\}=2^{-N}\binom{N}{(N+k)/2}
\]
whenever \(N+k\) is even.  By Stirling's formula, uniformly in \(k\),
\[
  \Prob\{S_N=k\}\le C N^{-1/2}.
\]
The interval \([a-t,a+t]\) has length \(2t\), hence contains at most
\(\lfloor t\rfloor+2\) points of this parity sublattice.  Therefore
\[
  \Prob\{\abs{S_N-a}\le t\}
  \le C(\lfloor t\rfloor+2)N^{-1/2}
  \le C'\frac{1+t}{\sqrt N}.
\]
When \(N\) is even,
\(\Prob\{S_N=0\}=2^{-N}\binom{N}{N/2}\asymp N^{-1/2}\).  When \(N\) is odd,
the two central values are \(\pm1\), and
\(\Prob\{|S_N|=1\}=2^{1-N}\binom{N}{(N+1)/2}\asymp N^{-1/2}\).
\end{proof}

\begin{proposition}[Root-of-unity lattice scale for Rademacher coefficients]\label{prop:rademacher-root-unity}
Let \(\varepsilon_0,\ldots,\varepsilon_m\) be independent Rademacher signs,
let \(G_m(z)=\sum_{j=0}^m\varepsilon_jz^j\), and put \(N=m+1\).  Fix a
root of unity \(\zeta\notin\{1,-1\}\), of order \(q\ge3\).  There is an
absolute constant \(C<\infty\) such that, for every \(a\in\C\) and every
\(t\ge0\),
\[
  \Prob\{\abs{G_m(\zeta)-a}\le t\}
  \le C\,\frac{q(1+t)^2}{N}.
\]
Thus, for each fixed order \(q\), a non-real root of unity has the
two-dimensional fixed-point lattice scale \(N^{-1}\), in contrast to the
one-dimensional endpoint scale \(N^{-1/2}\) in
Proposition~\ref{prop:rademacher-endpoint}.
\end{proposition}

\begin{proof}
It is enough to consider \(N\ge4q\), since otherwise the right-hand side is at
least a positive absolute constant and the claim follows by increasing \(C\).
Choose an integer \(k\in\{1,\ldots,q-1\}\) such that
\[
  \frac{\pi}{3}\le \frac{2\pi k}{q}\le \frac{2\pi}{3};
\]
such a \(k\) exists for every \(q\ge3\).  Then
\(\abs{\Im(\zeta^k)}\ge\sqrt3/2\).  For \(0\le\ell<q\), let
\[
  J_\ell=\{0\le j\le m:j\equiv \ell\pmod q\},\qquad
  S_\ell=\sum_{j\in J_\ell}\varepsilon_j .
\]
Then the \(S_\ell\)'s are independent Rademacher sums, the two classes
\(J_0,J_k\) satisfy \(\abs{J_0},\abs{J_k}\ge N/(2q)\), and
\[
  G_m(\zeta)=\sum_{\ell=0}^{q-1}\zeta^\ell S_\ell .
\]
Because \(\zeta^k\notin\R\) with imaginary part bounded below, the real-linear
map \((x,y)\mapsto x+\zeta^k y\) from \(\R^2\) to \(\C\) has inverse norm
bounded by an absolute constant.  Condition on all \(S_\ell\) with
\(\ell\notin\{0,k\}\).  For the resulting center \(b\in\C\), the set of real
pairs \((x,y)\) satisfying \(\abs{x+\zeta^k y-b}\le t\) has diameter \(O(1+t)\)
in both coordinate directions, hence contains at most \(C(1+t)^2\) points of
the parity lattice on which \((S_0,S_k)\) lives.

By Stirling's formula, a Rademacher sum of length \(M\) has point probabilities
bounded by \(CM^{-1/2}\).  Therefore each admissible pair \((S_0,S_k)=(x,y)\)
has probability at most
\[
  C\,\abs{J_0}^{-1/2}\abs{J_k}^{-1/2}\le C\,\frac qN.
\]
Multiplying by the \(O((1+t)^2)\) admissible pairs and then averaging over the
conditioned variables gives the stated bound.
\end{proof}

\begin{corollary}[Rademacher anti-concentration on a rational grid]\label{cor:rademacher-root-grid}
Let \(G_m(z)=\sum_{j=0}^m\varepsilon_jz^j\) with independent Rademacher signs,
put \(N=m+1\), and let \(\Omega_q=\{\zeta\in\C:\zeta^q=1\}\).  There is an
absolute constant \(C<\infty\) such that, for every \(q\ge1\) and \(t\ge0\),
\[
  \Prob\Bigl\{\min_{\zeta\in\Omega_q\setminus\{1,-1\}}\abs{G_m(\zeta)}
       \le t\Bigr\}
  \le C\,\frac{q^2(1+t)^2}{N},
\]
and, including the endpoint roots,
\[
  \Prob\Bigl\{\min_{\zeta\in\Omega_q}\abs{G_m(\zeta)}\le t\Bigr\}
  \le C\,\frac{1+t}{\sqrt N}
      +C\,\frac{q^2(1+t)^2}{N}.
\]
Thus rational grids of order \(q=o(\sqrt N)\) have no small Rademacher
boundary value with high probability at any fixed lattice scale, apart from the
endpoint contribution.
\end{corollary}

\begin{proof}
For the non-real roots, apply Proposition~\ref{prop:rademacher-root-unity} to
each \(\zeta\in\Omega_q\setminus\{1,-1\}\) and sum the resulting
\(O(q(1+t)^2/N)\) bounds over at most \(q\) roots.  The endpoint roots are
handled by Proposition~\ref{prop:rademacher-endpoint}; there are at most two of
them, and their contribution is \(O((1+t)/\sqrt N)\).
\end{proof}

\begin{corollary}[Bounded-denominator roots]\label{cor:rademacher-denom-grid}
Let \(G_m(z)=\sum_{j=0}^m\varepsilon_jz^j\) with independent Rademacher signs
and put \(N=m+1\).  For \(Q\ge1\), set
\[
  \mathcal R_Q:=\{\zeta\in\T:\zeta^q=1\ \text{for some}\ 1\le q\le Q\}.
\]
There is an absolute constant \(C<\infty\) such that, for every \(Q\ge1\) and
\(t\ge0\),
\[
  \Prob\Bigl\{\min_{\zeta\in\mathcal R_Q\setminus\{1,-1\}}\abs{G_m(\zeta)}
       \le t\Bigr\}
  \le C\,\frac{Q^3(1+t)^2}{N},
\]
and, including the endpoint roots,
\[
  \Prob\Bigl\{\min_{\zeta\in\mathcal R_Q}\abs{G_m(\zeta)}\le t\Bigr\}
  \le C\,\frac{1+t}{\sqrt N}
      +C\,\frac{Q^3(1+t)^2}{N}.
\]
Thus all rational angles of denominator \(Q=o(N^{1/3})\) are separated at any
fixed lattice scale with high probability, apart from the endpoint contribution.
\end{corollary}

\begin{proof}
Apply Corollary~\ref{cor:rademacher-root-grid} for each \(1\le q\le Q\) and
take a union bound.  The duplicate representation of the same root at several
denominators only makes this union bound larger.  Since
\(\sum_{q\le Q}q^2\le C Q^3\), the non-real bound follows.  The two endpoint
roots are handled once by Proposition~\ref{prop:rademacher-endpoint}.
\end{proof}

\begin{proposition}[Fixed-root exact atom local limit]
\label{prop:rademacher-small-denom-lower}
Let \(G_m(z)=\sum_{j=0}^m\varepsilon_jz^j\) with independent Rademacher signs
and put \(N=m+1\).  Let \(q\ge3\) be fixed, let \(\zeta\) be a primitive
\(q\)-th root of unity, and put \(d=\varphi(q)\).  For
\(0\le\ell<q\), let
\[
  M_\ell(N):=\#\{0\le j<N:j\equiv\ell\pmod q\},
  \qquad
  a_\ell(N)\equiv M_\ell(N)\pmod2,
\]
and let
\[
  L:\mathbb Z^q\to\mathbb Z[\zeta],\qquad
  L(s_0,\ldots,s_{q-1})=\sum_{\ell=0}^{q-1}s_\ell\zeta^\ell,
  \qquad K:=\ker L .
\]
Let \(\operatorname{covol}(K)\) denote the Euclidean covolume of \(K\) in
\(K\otimes_{\mathbb Z}\mathbb R\subset\mathbb R^q\), and set
\[
  \gamma_q:=\frac{(2q/\pi)^{\varphi(q)/2}}{\operatorname{covol}(K)}.
\]
Then, for all \(N\),
\[
  \Prob\{G_m(\zeta)=0\}\le C_q N^{-d/2}.
\]
Moreover fix a residue class \(R\pmod {2q}\), and write \(a^R=a(R)\) for the
constant value of the parity vector \(a(N)\) on \(N\equiv R\pmod {2q}\).  If
\[
  K\cap\{s\in\mathbb Z^q:s_\ell\equiv a_\ell(R)\pmod2
       \text{ for all }\ell\}\ne\varnothing ,
\]
then, along that residue class,
\[
  N^{\varphi(q)/2}\Prob\{G_m(\zeta)=0\}\longrightarrow \gamma_q .
\]
If the displayed intersection is empty, then
\(\Prob\{G_m(\zeta)=0\}=0\) for every \(N\equiv R\pmod {2q}\).  In particular
the old divisible case \(q\mid N\) is always admissible with this same local
limit constant, the primitive quadratic roots \(q=3,4,6\) have possible exact
atom scale \(N^{-1}\), and a primitive root of prime order \(p\) has possible
exact atom scale \(N^{-(p-1)/2}\).
\end{proposition}

\begin{proof}
Group the coefficients by residues modulo \(q\):
\[
  S_\ell=\sum_{\substack{0\le j\le m\\j\equiv \ell\pmod q}}\varepsilon_j,
  \qquad 0\le \ell<q.
\]
The variables \(S_0,\ldots,S_{q-1}\) are independent Rademacher sums of lengths
\(M_\ell(N)\), and \(M_\ell(N)=N/q+O_q(1)\).
The image of \(L\) has rank \(d:=\varphi(q)\), because
\(\mathbb Z[\zeta]\) is generated by \(1,\zeta,\ldots,\zeta^{d-1}\) and has
rank \(d\).  Hence the kernel lattice \(K\) has rank \(q-d\).  The event
\(\{G_m(\zeta)=0\}\) is exactly
\(\{(S_0,\ldots,S_{q-1})\in K\}\).

For the upper bound, choose a set \(J\subset\{0,\ldots,q-1\}\) of size \(q-d\)
such that the coordinate projection \(\pi_J:K\to\mathbb Z^J\) is injective;
this is possible by choosing a nonzero maximal minor in a basis matrix for the
fixed lattice \(K\).  If
\(p_\ell(s)=\Prob\{S_\ell=s\}\), Stirling's formula gives
\(\sup_s p_\ell(s)\le C_qN^{-1/2}\) for all large \(N\).  Therefore
\[
  \sum_{s\in K}\prod_{\ell=0}^{q-1}p_\ell(s_\ell)
  \le C_q N^{-d/2}\sum_{u\in\pi_J(K)}\prod_{j\in J}p_j(u_j)
  \le C_q N^{-d/2}.
\]
The last sum is at most \(1\), since it is a sub-sum of the product law of the
\(q-d\) free coordinates.

Now fix a residue class \(R\pmod {2q}\).  The parities
\(a_\ell(N)\) are constant on this class.  If no point of \(K\) has these
coordinate parities, then every point in the support of
\((S_0,\ldots,S_{q-1})\) lies outside \(K\), so the atom is zero.

Assume instead that \(a^R\in K\) has the required coordinate parities.  Then the
admissible support points in \(K\) are exactly the coset \(a^R+2K\): if
\(x\in K\cap 2\mathbb Z^q\), then \(x=2y\) with \(y\in\mathbb Z^q\), and
torsion-freeness of \(\mathbb Z[\zeta]\) gives \(y\in K\).

We now apply the local central limit theorem on this fixed lattice coset.  The
Hoeffding bound makes the contribution of
\(\max_\ell\abs{s_\ell}>N^{2/3}\) exponentially small.  Uniformly for
\(\max_\ell\abs{s_\ell}\le N^{2/3}\) on the support parities,
\[
  p_\ell(s_\ell)
  =\sqrt{\frac{2}{\pi M_\ell(N)}}
   \exp\!\left(-\frac{s_\ell^2}{2M_\ell(N)}\right)(1+o(1)).
\]
Since \(M_\ell(N)=N/q+O_q(1)\), uniformly in the same central range,
\[
  \prod_{\ell=0}^{q-1}p_\ell(s_\ell)
  =\left(\frac{2q}{\pi N}\right)^{q/2}
   \exp\!\left(-\frac{q\|s\|_2^2}{2N}\right)(1+o(1)).
\]
Let \(k=q-d=\operatorname{rank}K\).  The Riemann-sum limit for the coset
\(a^R+2K\) in \(K\otimes\mathbb R\) gives
\[
  \sum_{s\in a^R+2K}
   \exp\!\left(-\frac{q\|s\|_2^2}{2N}\right)
  =\frac{N^{k/2}}{\operatorname{covol}(2K)}
   \left(\frac{2\pi}{q}\right)^{k/2}+o(N^{k/2}).
\]
As \(\operatorname{covol}(2K)=2^k\operatorname{covol}(K)\), multiplying the
last two displays yields
\[
  \Prob\{G_m(\zeta)=0\}
  =\frac{(2q/\pi)^{d/2}}{\operatorname{covol}(K)}\,N^{-d/2}
   +o(N^{-d/2}).
\]

Finally, if \(q\mid N\), then all \(M_\ell(N)=M=N/q\) have the same parity.
When \(M\) is even the zero vector is an admissible point of \(K\); when \(M\)
is odd the vector \(\mathbf1=(1,\ldots,1)\) is admissible because
\(\sum_{\ell=0}^{q-1}\zeta^\ell=0\).  Thus every divisible residue class is
admissible and has the same local-limit constant \(\gamma_q\).
\end{proof}

\begin{corollary}[Finite cyclotomic-layer exact-zero asymptotic]
\label{cor:rademacher-finite-layer-zero}
Let \(\mathcal S\) be a fixed nonempty finite set of integers \(q\ge3\), and
for each \(q\in\mathcal S\) let \(E_q\) be the event \(G_m(\zeta_q)=0\), where
\(\zeta_q\) is any primitive \(q\)-th root of unity.  This event is independent
of the chosen primitive root.  Put
\[
  d_*=\min_{q\in\mathcal S}\varphi(q),
\]
and let \(\gamma_q\) be the covolume constant of
Proposition~\ref{prop:rademacher-small-denom-lower}.  Then
\[
  \Prob\Bigl(\bigcup_{q\in\mathcal S}E_q\Bigr)
  =N^{-d_*/2}
    \sum_{\substack{q\in\mathcal S\\ \varphi(q)=d_*}}
      \gamma_q\,\mathbf1_{\{N\bmod 2q\text{ is admissible for }q\}}
    +o_{\mathcal S}(N^{-d_*/2}).
\]
Thus every fixed finite rational grid has a complete leading exact-zero
asymptotic on its first nonempty cyclotomic degree layer.
\end{corollary}

\begin{proof}
Because \(G_m\) has integer coefficients, \(G_m(\zeta_q)=0\) for one primitive
\(q\)-th root if and only if the cyclotomic polynomial \(\Phi_q\) divides
\(G_m\), hence if and only if \(G_m\) vanishes at every primitive \(q\)-th root.
Thus \(E_q\) is well defined.

The single-order asymptotic is Proposition~\ref{prop:rademacher-small-denom-lower}.
The total contribution of orders with \(\varphi(q)>d_*\) is
\(o_{\mathcal S}(N^{-d_*/2})\).  It remains only to show that intersections
between distinct orders in the first layer are lower order.  Let \(q\ne r\) and
put \(L=\operatorname{lcm}(q,r)\).  Group the signs by residue classes modulo
\(L\):
\[
  S_a=\sum_{\substack{0\le j<N\\j\equiv a\pmod L}}\varepsilon_j,
  \qquad 0\le a<L.
\]
The two equations defining \(E_q\cap E_r\) are the kernel condition for the
\(\mathbb Z\)-linear map
\[
  T:\mathbb Z^L\to \mathbb Z[\zeta_q]\oplus\mathbb Z[\zeta_r],\qquad
  T(s)=\left(\sum_{a=0}^{L-1}s_a\zeta_q^a,\,
             \sum_{a=0}^{L-1}s_a\zeta_r^a\right).
\]
After tensoring with \(\mathbb Q\), this is the product quotient map from
\(\mathbb Q[x]/(x^L-1)\) to
\(\mathbb Q[x]/(\Phi_q)\oplus\mathbb Q[x]/(\Phi_r)\).  Since \(\Phi_q\) and
\(\Phi_r\) are distinct irreducible factors of \(x^L-1\), its rank is
\(\varphi(q)+\varphi(r)\).  The same projection-and-largest-atom argument used
for Proposition~\ref{prop:rademacher-small-denom-lower} therefore gives
\[
  \Prob(E_q\cap E_r)\le C_{q,r}N^{-(\varphi(q)+\varphi(r))/2}.
\]
For \(q,r\) in the first layer this is \(O_{\mathcal S}(N^{-d_*})\), hence
\(o_{\mathcal S}(N^{-d_*/2})\).  Inclusion--exclusion over the finite first
layer completes the proof.
\end{proof}

\begin{corollary}[Primitive quadratic exact atom constants]
\label{cor:rademacher-quadratic-atoms}
Let \(G_m(z)=\sum_{j=0}^m\varepsilon_jz^j\) with independent Rademacher signs,
and put \(N=m+1\).  If \(\zeta\) is a primitive third or sixth root of unity,
then
\[
  \Prob\{G_m(\zeta)=0\}=0\qquad\text{unless }3\mid N,
\]
while, as \(N\to\infty\) through multiples of \(3\),
\[
  N\Prob\{G_m(\zeta)=0\}\longrightarrow \frac{2\sqrt3}{\pi}.
\]
If \(\zeta\) is a primitive fourth root of unity, then
\[
  \Prob\{G_m(\zeta)=0\}=0\qquad\text{unless }4\mid N,
\]
while, as \(N\to\infty\) through multiples of \(4\),
\[
  N\Prob\{G_m(\zeta)=0\}\longrightarrow \frac4\pi.
\]
Thus the cyclotomic degree-two exact-zero floor is completely visible at the
level of divisibility classes and constants: cubic and primitive sixth roots
require \(3\mid N\) and have constant \(2\sqrt3/\pi\), while quartic roots
require \(4\mid N\) and have constant \(4/\pi\).
\end{corollary}

\begin{proof}
For \(q=3\), write \(S_\ell=\sum_{j\equiv\ell\pmod3}\varepsilon_j\).  Since
\(1+\zeta+\zeta^2=0\), the identity
\[
  S_0+S_1\zeta+S_2\zeta^2=0
\]
is equivalent to \(S_0=S_1=S_2\).  Such equality is compatible with the support
parities exactly when \(M_0(N),M_1(N),M_2(N)\) have the same parity.  Writing
\(N=3a+r\), the counts are
\[
  (M_0,M_1,M_2)=(a+\mathbf 1_{r\ge1},\,a+\mathbf 1_{r\ge2},\,a),
  \qquad r=0,1,2,
\]
so the parities agree exactly when \(r=0\), i.e. \(3\mid N\).  When \(N=3L\),
the three residue sums are independent copies of a length-\(L\) Rademacher sum.
If \(p_L(s)\) denotes the mass function of such a sum, then
\[
  \Prob\{G_m(\zeta)=0\}=\sum_s p_L(s)^3.
\]
Stirling's formula gives, uniformly for \(\abs{s}\le L^{2/3}\) on the support
parity,
\[
  p_L(s)=\sqrt{\frac{2}{\pi L}}\exp\!\left(-\frac{s^2}{2L}\right)(1+o(1)),
\]
and the Hoeffding tail makes the complementary range contribute \(o(L^{-1})\)
to the cubic sum.  Since the support lattice has spacing \(2\),
\[
\begin{aligned}
  \sum_s p_L(s)^3
  &=\frac12\int_{\mathbb R}
      \left(\frac{2}{\pi L}\right)^{3/2}
      \exp\!\left(-\frac{3x^2}{2L}\right)\,dx+o(L^{-1})  \\
  &=\frac{2}{\pi\sqrt3}\,\frac1L+o(L^{-1}).
\end{aligned}
\]
Multiplying by \(N=3L\) gives the limit \(2\sqrt3/\pi\).

For \(q=4\), a primitive root is \(\zeta=\pm i\).  Grouping the residue sums
modulo \(4\),
\[
  G_m(\zeta)=S_0+\zeta S_1-S_2-\zeta S_3,
\]
so \(G_m(\zeta)=0\) is equivalent to \(S_0=S_2\) and \(S_1=S_3\).  The necessary
and sufficient parity conditions are therefore
\[
  M_0\equiv M_2\pmod2,\qquad M_1\equiv M_3\pmod2.
\]
For \(N=4a+r\),
\[
  (M_0,M_1,M_2,M_3)=
  (a+\mathbf 1_{r\ge1},\,a+\mathbf 1_{r\ge2},\,
   a+\mathbf 1_{r\ge3},\,a),
\]
and the two congruences hold exactly when \(r=0\).  When \(N=4L\), the four
residue sums are independent copies of a length-\(L\) Rademacher sum, so
\[
  \Prob\{G_m(\zeta)=0\}
  =\left(\sum_s p_L(s)^2\right)^2.
\]
Vandermonde's identity and Stirling's formula give
\[
  \sum_s p_L(s)^2
  =4^{-L}\binom{2L}{L}
  =\frac1{\sqrt{\pi L}}(1+o(1)).
\]
Thus \(\Prob\{G_m(\zeta)=0\}=(\pi L)^{-1}+o(L^{-1})\), and multiplying by
\(N=4L\) gives the limit \(4/\pi\).

For \(q=6\), let \(\zeta\) be primitive, so \(\zeta^3=-1\) and
\(\zeta^2-\zeta+1=0\).  Put
\[
  D_0=S_0-S_3,\qquad D_1=S_1-S_4,\qquad D_2=S_2-S_5 .
\]
Then
\[
  G_m(\zeta)=D_0+D_1\zeta+D_2\zeta^2
  =(D_0-D_2)+(D_1+D_2)\zeta,
\]
so \(G_m(\zeta)=0\) is equivalent to \(D_0=D_2=-D_1\).  The parity of \(D_j\) is
\(M_j+M_{j+3}\pmod2\).  For \(N=6a+r\),
\[
\begin{aligned}
  M_0+M_3&\equiv \mathbf1_{r\ge1}+\mathbf1_{r\ge4},\\
  M_1+M_4&\equiv \mathbf1_{r\ge2}+\mathbf1_{r\ge5},\\
  M_2+M_5&\equiv \mathbf1_{r\ge3}
  \qquad\pmod2.
\end{aligned}
\]
These three parities are equal exactly for \(r=0\) and \(r=3\), equivalently
\(3\mid N\).  When they are unequal, no zero is possible.  When \(N=3L\), the
random variables \(D_0,D_1,D_2\) are independent and each has the law of a
length-\(L\) Rademacher sum: if \(N=6a\) then \(L=2a\) and \(D_j\) is the
difference of two length-\(a\) sums, while if \(N=6a+3\) then \(L=2a+1\) and
\(D_j\) is the difference of a length-\(a+1\) sum and a length-\(a\) sum.  By
symmetry,
\[
  \Prob\{G_m(\zeta)=0\}
  =\sum_s p_L(s)p_L(-s)p_L(s)
  =\sum_s p_L(s)^3.
\]
The \(q=3\) computation therefore gives the same limit \(2\sqrt3/\pi\).
\end{proof}

\begin{corollary}[Pairwise primitive-quadratic overlap constants]
\label{cor:rademacher-quadratic-overlaps}
Let \(E_q\) denote the event that \(G_m\) vanishes at a primitive \(q\)-th
root, for \(q\in\{3,4,6\}\).  Then
\[
  \Prob(E_3\cap E_6)=0\qquad\text{unless }6\mid N,
\]
and, as \(N\to\infty\) through multiples of \(6\),
\[
  N^2\Prob(E_3\cap E_6)\longrightarrow \frac{48}{\pi^2}.
\]
Moreover
\[
  \Prob(E_3\cap E_4)=\Prob(E_4\cap E_6)=0
  \qquad\text{unless }12\mid N,
\]
and, as \(N\to\infty\) through multiples of \(12\),
\[
  N^2\Prob(E_3\cap E_4)\longrightarrow \frac{8\sqrt3}{\pi^2},
  \qquad
  N^2\Prob(E_4\cap E_6)\longrightarrow \frac{8\sqrt3}{\pi^2}.
\]
\end{corollary}

\begin{proof}
First group the coefficients modulo \(6\), and write the corresponding sums as
\(T_0,\ldots,T_5\).  The event \(E_3\) is
\[
  T_0+T_3=T_1+T_4=T_2+T_5,
\]
while the calculation in the preceding proof shows that \(E_6\) is
\[
  T_0-T_3=T_2-T_5=-(T_1-T_4).
\]
These two systems are equivalent to
\[
  T_0=T_2=T_4,\qquad T_1=T_3=T_5.
\]
The parity requirements for the three even residue classes and the three odd
residue classes are simultaneously compatible exactly when \(6\mid N\).  If
\(N=6L\), the six sums are independent length-\(L\) Rademacher sums.  Hence
\[
  \Prob(E_3\cap E_6)=\left(\sum_t p_L(t)^3\right)^2
  =\left(\frac{2}{\pi\sqrt3}\,L^{-1}+o(L^{-1})\right)^2,
\]
which gives the constant \(48/\pi^2\) after multiplying by \(N^2=36L^2\).

For \(E_3\cap E_4\), the individual admissibility criteria in
Corollary~\ref{cor:rademacher-quadratic-atoms} already force \(12\mid N\).
Assume \(N=12L\).  Identify the residue classes modulo \(12\) with the
\(3\times4\) grid \((a,b)\in\mathbb Z/3\mathbb Z\times\mathbb Z/4\mathbb Z\),
and let \(X_{a,b}\) be the corresponding independent length-\(L\) sums.  Let
\[
  R_a=\sum_{b=0}^3X_{a,b},\qquad
  C_b=\sum_{a=0}^2X_{a,b}.
\]
Then \(E_3\cap E_4\) is the fixed-rank lattice event
\[
  R_0=R_1=R_2,\qquad C_0=C_2,\qquad C_1=C_3 .
\]
Equivalently \(AX=0\), where
\[
  AX=(R_0-R_1,\ R_0-R_2,\ C_0-C_2,\ C_1-C_3).
\]
The four rows of \(A\) have Gram matrix
\[
  AA^T=
  \begin{pmatrix}
    8&4&0&0\\
    4&8&0&0\\
    0&0&6&0\\
    0&0&0&6
  \end{pmatrix},
\]
whose determinant is \(1728\).  The \(4\times4\) minor using columns
\[
  (0,0),\qquad (0,1),\qquad (1,0),\qquad (2,2)
\]
has determinant \(1\).  Thus the standard covolume formula for the kernel of a
full-rank integer map gives
\[
  \operatorname{covol}(\ker A)=\sqrt{\det(AA^T)}=24\sqrt3 .
\]
The same fixed-lattice local central limit argument used in
Proposition~\ref{prop:rademacher-small-denom-lower} gives
\[
  \Prob(E_3\cap E_4)
  =\frac{(2/\pi)^2}{24\sqrt3}\,L^{-2}+o(L^{-2}).
\]
Since \(N=12L\), this is \(8\sqrt3/(\pi^2N^2)+o(N^{-2})\).

Finally, for \(E_4\cap E_6\), the individual criteria force \(12\mid N\) again.
With residues grouped modulo \(12\), the four defining equations have Gram
matrix
\[
  \begin{pmatrix}
    6&0&0&0\\
    0&6&0&0\\
    0&0&8&-4\\
    0&0&-4&8
  \end{pmatrix},
\]
and the minor using residues \(0,1,2,3\) has determinant \(1\).  The kernel
covolume is again \(24\sqrt3\), and the same local limit gives the same
constant \(8\sqrt3/\pi^2\).
\end{proof}

\begin{corollary}[Primitive-quadratic atom expansions to second order]
\label{cor:rademacher-quadratic-second-order}
Let \(E_q\) denote the event that \(G_m\) vanishes at a primitive \(q\)-th root,
for \(q\in\{3,4,6\}\).  Then
\[
  \Prob(E_3)=\mathbf1_{\{3\mid N\}}
  \left(\frac{2\sqrt3}{\pi N}-\frac{2\sqrt3}{\pi N^2}\right)
  +o(N^{-2}),
\]
\[
  \Prob(E_6)=\mathbf1_{\{3\mid N\}}
  \left(\frac{2\sqrt3}{\pi N}-\frac{2\sqrt3}{\pi N^2}\right)
  +o(N^{-2}),
\]
and
\[
  \Prob(E_4)=\mathbf1_{\{4\mid N\}}
  \left(\frac4{\pi N}-\frac4{\pi N^2}\right)+o(N^{-2}).
\]
\end{corollary}

\begin{proof}
The inadmissible classes are exactly those identified in
Corollary~\ref{cor:rademacher-quadratic-atoms}.  It remains to sharpen the
admissible binomial sums.  Let \(p_L\) be the mass function of a length-\(L\)
Rademacher sum.  Stirling's formula, uniformly on the support for
\(\abs{s}\le L^{1/2+\eta}\) with any fixed \(0<\eta<1/6\), gives
\[
\begin{aligned}
  p_L(s)
  &=\sqrt{\frac{2}{\pi L}}\exp\!\left(-\frac{x^2}{2}\right)
    \left(1-\frac{x^4-6x^2+3}{12L}
      +O_\eta(L^{-2}(1+\abs{x}^{12}))\right),\\
  x&=\frac{s}{\sqrt L}.
\end{aligned}
\]
The complement contributes \(O(e^{-cL^{2\eta}})\) to the cubic sum.  Using the
spacing \(2/\sqrt L\) of the support lattice and the Gaussian Riemann-sum limit,
\[
\begin{aligned}
  \sum_s p_L(s)^3
  &=\frac12\left(\frac2{\pi L}\right)^{3/2}L^{1/2}
    \int_\mathbb R e^{-3x^2/2}
    \left(1-\frac{x^4-6x^2+3}{4L}\right)\,dx
    +o(L^{-2})  \\
  &=\frac{2}{\pi\sqrt3}\,L^{-1}
    -\frac{2}{3\pi\sqrt3}\,L^{-2}+o(L^{-2}).
\end{aligned}
\]
The identities in the proof of Corollary~\ref{cor:rademacher-quadratic-atoms}
show that both \(E_3\) and \(E_6\), on their admissible classes, are given by
this cubic sum with \(L=N/3\).  This yields the first two expansions.

For \(E_4\), on \(N=4L\),
\[
  \Prob(E_4)=\left(\sum_s p_L(s)^2\right)^2
  =\left(4^{-L}\binom{2L}{L}\right)^2.
\]
The central-binomial expansion
\[
  4^{-L}\binom{2L}{L}
  =\frac1{\sqrt{\pi L}}\left(1-\frac1{8L}+O(L^{-2})\right)
\]
therefore gives
\[
  \Prob(E_4)=\frac1{\pi L}-\frac1{4\pi L^2}+o(L^{-2})
  =\frac4{\pi N}-\frac4{\pi N^2}+o(N^{-2}).
\]
\end{proof}

\begin{corollary}[Primitive degree-four exact atom constants]
\label{cor:rademacher-degree-four-atoms}
Let \(G_m(z)=\sum_{j=0}^m\varepsilon_jz^j\) with independent Rademacher signs,
and put \(N=m+1\).  If \(\zeta\) is a primitive fifth or tenth root of unity,
then
\[
  \Prob\{G_m(\zeta)=0\}=0\qquad\text{unless }5\mid N,
\]
while, as \(N\to\infty\) through multiples of \(5\),
\[
  N^2\Prob\{G_m(\zeta)=0\}\longrightarrow \frac{20\sqrt5}{\pi^2}.
\]
If \(\zeta\) is a primitive eighth root of unity, then
\[
  \Prob\{G_m(\zeta)=0\}=0\qquad\text{unless }8\mid N,
\]
while, as \(N\to\infty\) through multiples of \(8\),
\[
  N^2\Prob\{G_m(\zeta)=0\}\longrightarrow \frac{64}{\pi^2}.
\]
If \(\zeta\) is a primitive twelfth root of unity, then
\[
  \Prob\{G_m(\zeta)=0\}=0\qquad\text{unless }6\mid N,
\]
while, as \(N\to\infty\) through multiples of \(6\),
\[
  N^2\Prob\{G_m(\zeta)=0\}\longrightarrow \frac{48}{\pi^2}.
\]
\end{corollary}

\begin{proof}
Let \(p_L\) denote the mass function of a length-\(L\) Rademacher sum.  We first
record the local-limit sum used for the fifth and tenth roots:
\[
  \sum_t p_L(t)^5
  =\frac12\int_{\mathbb R}
       \left(\frac{2}{\pi L}\right)^{5/2}
       \exp\!\left(-\frac{5x^2}{2L}\right)\,dx+o(L^{-2})
  =\frac{4}{\sqrt5\,\pi^2}\,L^{-2}+o(L^{-2}).
\]
The Hoeffding tail and Stirling approximation justify the same central-range
Riemann-sum replacement as in the proof of
Corollary~\ref{cor:rademacher-quadratic-atoms}.

For \(q=5\), the identity \(1+\zeta+\cdots+\zeta^4=0\) makes
\[
  S_0+S_1\zeta+\cdots+S_4\zeta^4=0
\]
equivalent to \(S_0=\cdots=S_4\).  The support parities of the five residue
sums agree exactly when the five residue-class counts modulo \(5\) have the
same parity, which for \(N=5a+r\) happens exactly at \(r=0\).  When \(N=5L\),
therefore,
\[
  \Prob\{G_m(\zeta)=0\}=\sum_t p_L(t)^5,
\]
and multiplying the displayed asymptotic by \(N^2=25L^2\) gives
\(20\sqrt5/\pi^2\).

For \(q=10\), write
\[
  D_a=S_a-S_{a+5},\qquad 0\le a<5.
\]
Since \(\zeta^5=-1\) and
\(1-\zeta+\zeta^2-\zeta^3+\zeta^4=0\), the equation \(G_m(\zeta)=0\) is
equivalent to
\[
  D_a=t(-1)^a,\qquad 0\le a<5,
\]
for some integer \(t\).  The parity of \(D_a\) is
\(M_a+M_{a+5}\pmod2\), the count of the residue class \(a\) modulo \(5\).
These five parities agree exactly when \(5\mid N\).  If \(N=5L\), each \(D_a\)
has the law of a length-\(L\) Rademacher sum, and symmetry gives again
\[
  \Prob\{G_m(\zeta)=0\}=\sum_t p_L(t)^5.
\]
Thus the tenth-root constant is the same \(20\sqrt5/\pi^2\).

For \(q=8\), put \(D_a=S_a-S_{a+4}\) for \(0\le a<4\).  A primitive eighth root
has degree \(4\), so \(1,\zeta,\zeta^2,\zeta^3\) are linearly independent over
\(\mathbb Q\), and \(G_m(\zeta)=0\) is equivalent to \(D_a=0\) for every
\(a<4\).  This requires \(M_a+M_{a+4}\) to be even for all \(a\), i.e. every
residue-class count modulo \(4\) is even; for \(N=4a+r\) this is equivalent to
\(r=0\) and \(a\) even, or \(8\mid N\).  When \(N=8L\), the four equalities are
independent and
\[
  \Prob\{G_m(\zeta)=0\}
  =\left(\sum_t p_L(t)^2\right)^4.
\]
By Vandermonde's identity,
\[
  \sum_t p_L(t)^2=4^{-L}\binom{2L}{L}
  =\frac1{\sqrt{\pi L}}(1+o(1)),
\]
and multiplying by \(N^2=64L^2\) gives \(64/\pi^2\).

For \(q=12\), put \(D_a=S_a-S_{a+6}\) for \(0\le a<6\).  The two relations
\[
  1-\zeta^2+\zeta^4=0,\qquad
  \zeta-\zeta^3+\zeta^5=0
\]
show that \(G_m(\zeta)=0\) is equivalent to
\[
  D_0=D_4=-D_2,\qquad D_1=D_5=-D_3 .
\]
The needed parity equalities say that the three even residue-class counts
modulo \(6\) have the same parity and that the three odd residue-class counts
modulo \(6\) have the same parity.  For \(N=6a+r\), a direct check of
\(0\le r<6\) shows that this happens exactly when \(r=0\), i.e. \(6\mid N\).
If \(N=6L\), each \(D_a\) has the law of a length-\(L\) Rademacher sum, so
\[
  \Prob\{G_m(\zeta)=0\}
  =\left(\sum_t p_L(t)^3\right)^2.
\]
The cubic local-limit sum computed in
Corollary~\ref{cor:rademacher-quadratic-atoms} is
\[
  \sum_t p_L(t)^3=\frac{2}{\pi\sqrt3}\,L^{-1}+o(L^{-1}),
\]
and multiplying its square by \(N^2=36L^2\) gives \(48/\pi^2\).
\end{proof}

\begin{corollary}[Bounded-denominator exact zeros]
\label{cor:rademacher-denom-zero}
Let \(G_m(z)=\sum_{j=0}^m\varepsilon_jz^j\) with independent Rademacher signs,
put \(N=m+1\), and let
\[
  \mathcal R_Q^\circ:=
  \{\zeta\in\T:\zeta^q=1\text{ for some }1\le q\le Q\}
  \setminus\{1,-1\}.
\]
For every fixed \(Q\ge1\) there is \(C_Q<\infty\) such that, for all \(N\),
\[
  \Prob\{\exists\,\zeta\in\mathcal R_Q^\circ:\ G_m(\zeta)=0\}
  \le \frac{C_Q}{N}.
\]
Moreover,
\[
  \Prob\{\exists\,\zeta\in\mathcal R_Q^\circ:
    \varphi(\operatorname{ord}\zeta)>2,\ G_m(\zeta)=0\}
  \le \frac{C_Q}{N^2}.
\]
The first bound is sharp at fixed \(Q\) as soon as \(Q\ge3\): if \(3\mid N\),
then
\[
  \Prob\{\exists\,\zeta\in\mathcal R_Q^\circ:\ G_m(\zeta)=0\}
  \ge \frac{c}{N}
\]
with an absolute \(c>0\).
\end{corollary}

\begin{proof}
Every point of \(\mathcal R_Q^\circ\) is a primitive root of some order
\(3\le q\le Q\).  Proposition~\ref{prop:rademacher-small-denom-lower} gives
\[
  \Prob\{G_m(\zeta)=0\}\le C_q N^{-\varphi(q)/2}
\]
for a primitive root of order \(q\).  Summing over the finitely many roots with
order at most \(Q\) gives the \(C_Q/N\) bound, since
\(\varphi(q)\ge2\) for every non-real root of unity.  If
\(\varphi(q)>2\), then \(\varphi(q)\ge4\), giving the \(C_Q/N^2\) bound after
the same finite union.

For sharpness, take a primitive cubic root \(\zeta\).  When \(3\mid N\), the
divisible case in Proposition~\ref{prop:rademacher-small-denom-lower} is
admissible and \(\varphi(3)=2\), so
\[
  \Prob\{G_m(\zeta)=0\}\ge cN^{-1}
\]
after decreasing \(c\) to absorb the finitely many small \(N\).  This event is
contained in the bounded-denominator union whenever \(Q\ge3\).
\end{proof}

\begin{corollary}[Peeled bounded-denominator cyclotomic layers]
\label{cor:rademacher-denom-zero-layer}
With the notation of Corollary~\ref{cor:rademacher-denom-zero}, fix \(D_0\ge0\)
and assume that there is an order \(3\le q\le Q\) with
\(\varphi(q)>D_0\).  Put
\[
  d_{Q,D_0}:=\min\{\varphi(q):3\le q\le Q,\ \varphi(q)>D_0\}.
\]
Then
\[
\begin{aligned}
&\Prob\{\exists\,\zeta\in\mathcal R_Q^\circ:
     \varphi(\operatorname{ord}\zeta)>D_0,\ G_m(\zeta)=0\} \\
&\qquad =
N^{-d_{Q,D_0}/2}
  \sum_{\substack{3\le q\le Q\\ \varphi(q)=d_{Q,D_0}}}
    \gamma_q\,\mathbf1_{\{N\bmod 2q\text{ is admissible for }q\}}
  +o_Q(N^{-d_{Q,D_0}/2}).
\end{aligned}
\]
Thus every fixed bounded-denominator exact-zero grid has a leading coefficient
on each peeled cyclotomic degree layer; only continuum or growing-denominator
questions remain outside this finite lattice expansion.
\end{corollary}

\begin{proof}
The event in the first line is the union of the primitive-order events \(E_q\)
from Corollary~\ref{cor:rademacher-finite-layer-zero} over the finite set
\[
  \mathcal S_{Q,D_0}:=\{q:3\le q\le Q,\ \varphi(q)>D_0\}.
\]
Indeed, integer coefficients make vanishing at one primitive \(q\)-th root
equivalent to vanishing at every primitive \(q\)-th root.  Applying
Corollary~\ref{cor:rademacher-finite-layer-zero} to \(\mathcal S_{Q,D_0}\)
gives the displayed formula.
\end{proof}

\begin{corollary}[Bounded-denominator exact-zero leading coefficient]
\label{cor:rademacher-denom-zero-leading}
With the notation of Corollary~\ref{cor:rademacher-denom-zero},
\[
  \Prob\{\exists\,\zeta\in\mathcal R_Q^\circ:\ G_m(\zeta)=0\}
  =\frac{c_Q(N)}{N}+o_Q(N^{-1}),
\]
where the periodic coefficient is
\[
  c_Q(N)=
  \frac{2\sqrt3}{\pi}\mathbf1_{\{Q\ge3\}}\mathbf1_{\{3\mid N\}}
  +\frac4\pi\mathbf1_{\{Q\ge4\}}\mathbf1_{\{4\mid N\}}
  +\frac{2\sqrt3}{\pi}\mathbf1_{\{Q\ge6\}}\mathbf1_{\{3\mid N\}} .
\]
Equivalently, the only order \(N^{-1}\) exact-zero contributions in a fixed
bounded-denominator non-real grid are the primitive cubic pair, the primitive
quartic pair, and the primitive sixth pair; all higher-degree roots and all
overlaps among these three pairs are lower order.
\end{corollary}

\begin{proof}
Let \(E_q\) be the event that \(G_m(\zeta)=0\) for a primitive \(q\)-th root
\(\zeta\).  Since the coefficients are real, the two conjugate primitive roots
of a fixed order give the same event.  Corollary~\ref{cor:rademacher-quadratic-atoms}
gives
\[
\begin{aligned}
  \Prob(E_3)&=\frac{2\sqrt3}{\pi N}\mathbf1_{\{3\mid N\}}+o(N^{-1}),\\
  \Prob(E_4)&=\frac{4}{\pi N}\mathbf1_{\{4\mid N\}}+o(N^{-1}),\\
  \Prob(E_6)&=\frac{2\sqrt3}{\pi N}\mathbf1_{\{3\mid N\}}+o(N^{-1}).
\end{aligned}
\]
The remaining primitive orders have \(\varphi(q)>2\), and their total
contribution is \(O_Q(N^{-2})\) by Corollary~\ref{cor:rademacher-denom-zero}.

Corollary~\ref{cor:rademacher-quadratic-overlaps} gives the sharper explicit
\(N^{-2}\) constants for all pairwise intersections among \(E_3,E_4,E_6\), in
particular these intersections are \(O(N^{-2})\).  Inclusion--exclusion over
the subset of \(\{3,4,6\}\) with order at most \(Q\), together with the
\(O_Q(N^{-2})\) higher-degree contribution, gives the displayed asymptotic.
\end{proof}

\begin{corollary}[First nonquadratic bounded-denominator coefficient]
\label{cor:rademacher-denom-zero-degree-four-leading}
With the notation of Corollary~\ref{cor:rademacher-denom-zero},
\[
  \Prob\{\exists\,\zeta\in\mathcal R_Q^\circ:
     \varphi(\operatorname{ord}\zeta)>2,\ G_m(\zeta)=0\}
  =\frac{d_Q(N)}{N^2}+o_Q(N^{-2}),
\]
where
\[
\begin{aligned}
  d_Q(N)=&
  \frac{20\sqrt5}{\pi^2}\mathbf1_{\{Q\ge5\}}\mathbf1_{\{5\mid N\}}
  +\frac{64}{\pi^2}\mathbf1_{\{Q\ge8\}}\mathbf1_{\{8\mid N\}} \\
  &+\frac{20\sqrt5}{\pi^2}\mathbf1_{\{Q\ge10\}}\mathbf1_{\{5\mid N\}}
  +\frac{48}{\pi^2}\mathbf1_{\{Q\ge12\}}\mathbf1_{\{6\mid N\}} .
\end{aligned}
\]
Equivalently, after the primitive quadratic roots are removed, the only
order-\(N^{-2}\) exact-zero contributions in a fixed bounded-denominator
non-real grid come from primitive orders \(5,8,10,12\); all higher cyclotomic
orders and all overlaps are lower order.
\end{corollary}

\begin{proof}
This is Corollary~\ref{cor:rademacher-denom-zero-layer} with \(D_0=2\), together
with the explicit degree-four constants and admissible classes from
Corollary~\ref{cor:rademacher-degree-four-atoms}.  The only orders with
\(\varphi(q)=4\) are \(q=5,8,10,12\).
\end{proof}

\begin{corollary}[Bounded-denominator exact-zero expansion through \(N^{-2}\)]
\label{cor:rademacher-denom-zero-second-order}
With the notation of Corollary~\ref{cor:rademacher-denom-zero},
\[
  \Prob\{\exists\,\zeta\in\mathcal R_Q^\circ:\ G_m(\zeta)=0\}
  =\frac{c_Q(N)}{N}+\frac{e_Q(N)}{N^2}+o_Q(N^{-2}),
\]
where \(c_Q(N)\) is the coefficient in
Corollary~\ref{cor:rademacher-denom-zero-leading}.  Let \(d_Q(N)\) be the
degree-four coefficient displayed in
Corollary~\ref{cor:rademacher-denom-zero-degree-four-leading}.  Then
\[
\begin{aligned}
  e_Q(N)=&\,d_Q(N)
  -\frac{2\sqrt3}{\pi}\mathbf1_{\{Q\ge3\}}\mathbf1_{\{3\mid N\}}
  -\frac4\pi\mathbf1_{\{Q\ge4\}}\mathbf1_{\{4\mid N\}}  \\
  &-\frac{2\sqrt3}{\pi}\mathbf1_{\{Q\ge6\}}\mathbf1_{\{3\mid N\}}
  -\frac{8\sqrt3}{\pi^2}\mathbf1_{\{Q\ge4\}}\mathbf1_{\{12\mid N\}} \\
  &-\frac{48}{\pi^2}\mathbf1_{\{Q\ge6\}}\mathbf1_{\{6\mid N\}}
  -\frac{8\sqrt3}{\pi^2}\mathbf1_{\{Q\ge6\}}\mathbf1_{\{12\mid N\}} .
\end{aligned}
\]
Thus the full unpeeled second-order term is not just the first nonquadratic
cyclotomic layer: it also contains the second-order corrections to the
degree-two single atoms and the pairwise overlaps among the degree-two order
events.
\end{corollary}

\begin{proof}
Let \(E_q\) be the primitive-order event as above.  The single-event
contributions from \(q=3,4,6\) are given to order \(N^{-2}\) by
Corollary~\ref{cor:rademacher-quadratic-second-order}.  The single-event
contributions from \(\varphi(q)=4\) sum to \(d_Q(N)N^{-2}+o_Q(N^{-2})\) by
Corollary~\ref{cor:rademacher-denom-zero-degree-four-leading}.  All primitive
orders with \(\varphi(q)\ge6\) contribute \(O_Q(N^{-3})\) by
Proposition~\ref{prop:rademacher-small-denom-lower}.

For intersections, Corollary~\ref{cor:rademacher-quadratic-overlaps} gives the
three degree-two pairwise terms.  They enter with a minus sign in
inclusion--exclusion and give exactly the last three displayed terms in
\(e_Q(N)\).  Every other pairwise intersection contains either one degree-two
and one degree-four-or-higher event, or two degree-four-or-higher events.  The
product-quotient rank argument used in
Corollary~\ref{cor:rademacher-finite-layer-zero} gives rank at least \(6\) in
the first case and at least \(8\) in the second, hence these intersections are
\(O_Q(N^{-3})\).  The same argument gives \(O_Q(N^{-3})\) for every triple
intersection, including \(E_3\cap E_4\cap E_6\), because the three distinct
cyclotomic factors \(\Phi_3,\Phi_4,\Phi_6\) have total rank \(6\).  Higher
inclusion--exclusion terms are smaller.  Combining these terms gives the
displayed expansion.
\end{proof}

\begin{corollary}[Bounded-denominator exact zeros including endpoints]
\label{cor:rademacher-denom-zero-with-endpoints}
Let \(G_m(z)=\sum_{j=0}^m\varepsilon_jz^j\) with independent Rademacher signs,
put \(N=m+1\), and let
\[
  \mathcal R_Q:=\{\zeta\in\T:\zeta^q=1\text{ for some }1\le q\le Q\}.
\]
Define
\[
  b_N:=\mathbf1_{\{2\mid N\}}2^{-N}\binom{N}{N/2},
  \qquad
  h_N:=\mathbf1_{\{4\mid N\}}
       \left(2^{-N/2}\binom{N/2}{N/4}\right)^2 .
\]
For every fixed \(Q\ge1\),
\[
  \Prob\{\exists\,\zeta\in\mathcal R_Q:\ G_m(\zeta)=0\}
  =
  (1+\mathbf1_{\{Q\ge2\}})b_N
  -\mathbf1_{\{Q\ge2\}}h_N
  +\frac{c_Q(N)}{N}+o_Q(N^{-1}),
\]
where \(c_Q(N)\) is the non-real coefficient in
Corollary~\ref{cor:rademacher-denom-zero-leading}.  Thus the full
bounded-denominator root-of-unity grid splits, through order \(N^{-1}\), into
the exact endpoint-union term and the leading non-real cyclotomic contribution.
\end{corollary}

\begin{proof}
Let \(E_+=\{G_m(1)=0\}\), \(E_-=\{G_m(-1)=0\}\), and
\[
  F_Q=\{\exists\,\zeta\in\mathcal R_Q^\circ:\ G_m(\zeta)=0\}.
\]
The endpoint variables \(G_m(1)=\sum_{j=0}^{N-1}\varepsilon_j\) and
\(G_m(-1)=\sum_{j=0}^{N-1}(-1)^j\varepsilon_j\) both have the law of a
length-\(N\) Rademacher sum, hence
\[
  \Prob(E_+)=\Prob(E_-)=b_N.
\]
For the joint endpoint event, write
\[
  A=\sum_{\substack{0\le j<N\\ j\ {\rm even}}}\varepsilon_j,\qquad
  B=\sum_{\substack{0\le j<N\\ j\ {\rm odd}}}\varepsilon_j .
\]
Then \(G_m(1)=A+B\) and \(G_m(-1)=A-B\), so \(E_+\cap E_-\) is the event
\(\{A=0,\ B=0\}\).  This is impossible unless both parity classes have even
cardinality.  Since the two parity classes have equal size \(N/2\) when \(N\)
is even, this happens exactly when \(4\mid N\), and then independence gives
\[
  \Prob(E_+\cap E_-)
  =\left(2^{-N/2}\binom{N/2}{N/4}\right)^2.
\]
Therefore the endpoint contribution is
\[
  \Prob(E_+\cup E_-)
  =2b_N-h_N
\]
when \(Q\ge2\), while for \(Q=1\) it is just \(b_N\).

It remains to check that intersections between endpoint zeros and non-real
bounded-denominator zeros are \(o_Q(N^{-1})\).  Fix a primitive non-real
\(q\)-th root with \(3\le q\le Q\).  The event \(G_m(1)=0\) is the divisibility
condition \(\Phi_1\mid G_m\), and \(G_m(-1)=0\) is \(\Phi_2\mid G_m\).  Grouping
coefficients modulo the relevant least common multiple, the simultaneous
endpoint and primitive-\(q\) conditions are the kernel condition for the product
quotient map onto the distinct cyclotomic factors \(\Phi_1,\Phi_q\), or
\(\Phi_2,\Phi_q\).  Its rational rank is \(1+\varphi(q)\).  The
projection-and-largest-atom argument used in
Corollary~\ref{cor:rademacher-finite-layer-zero} gives
\[
  \Prob(E_\pm\cap E_q)\le C_q N^{-(1+\varphi(q))/2}=O_q(N^{-3/2}).
\]
With both endpoints and \(E_q\) the rank is \(2+\varphi(q)\), so the bound is
even smaller.  Summing over the finite set \(3\le q\le Q\) gives
\(\Prob((E_+\cup E_-)\cap F_Q)=o_Q(N^{-1})\).  Finally,
Corollary~\ref{cor:rademacher-denom-zero-leading} gives
\[
  \Prob(F_Q)=\frac{c_Q(N)}{N}+o_Q(N^{-1}),
\]
and inclusion--exclusion with the endpoint union proves the formula.
\end{proof}

\begin{proposition}[Fixed-radius lower bound at rational angles]
\label{prop:rademacher-root-fixed-radius-lower}
Let \(G_m(z)=\sum_{j=0}^m\varepsilon_jz^j\) with independent Rademacher signs
and put \(N=m+1\).  Let \(q\ge3\) be fixed, and let \(\zeta\) be a primitive
\(q\)-th root of unity.  There are constants \(c_q,t_q>0\) such that, for all
\(N\ge1\),
\[
  \Prob\{\abs{G_m(\zeta)}\le t_q\}\ge \frac{c_q}{N}.
\]
Thus the fixed-point upper bound \(O_q(N^{-1})\) from
Proposition~\ref{prop:rademacher-root-unity} is sharp at every fixed non-real
rational angle at a constant lattice scale, even though the exact atom may be
smaller when \(\varphi(q)>2\).
\end{proposition}

\begin{proof}
It is enough to prove the estimate for \(N\ge4q\); finitely many smaller \(N\)
are absorbed by increasing \(t_q\) and decreasing \(c_q\).  Choose
\(k\in\{1,\ldots,q-1\}\) with
\(\abs{\Im(\zeta^k)}\ge\sqrt3/2\), as in the proof of
Proposition~\ref{prop:rademacher-root-unity}, and group the residue-class sums
\[
  S_\ell=\sum_{\substack{0\le j\le m\\j\equiv \ell\pmod q}}\varepsilon_j,
  \qquad 0\le \ell<q.
\]
Let \(M_\ell=\#\{0\le j<N:j\equiv\ell\pmod q\}\).  Then
\[
  \frac{N}{2q}\le M_\ell\le \frac{2N}{q}
  \qquad(0\le\ell<q)
\]
for \(N\ge4q\).
Condition on all \(S_\ell\) with \(\ell\notin\{0,k\}\), and write
\[
  B=\sum_{\ell\notin\{0,k\}}\zeta^\ell S_\ell .
\]
The event
\[
  E=\{\abs{S_\ell}\le C_q\sqrt N\text{ for every }\ell\notin\{0,k\}\}
\]
has probability at least \(c_q>0\), by the central binomial lower bound on a
window of length \(\asymp_q\sqrt{M_\ell}\), applied independently to the
\(q-2\) conditioned residue classes.  On \(E\), \(\abs B\le C_q\sqrt N\).

The real-linear map \(T(x,y)=x+\zeta^k y\) has inverse norm bounded in terms of
q.  The image of the parity lattice on which \((S_0,S_k)\) lives is one of the
four translates of the fixed lattice \(2\mathbb Z+2\zeta^k\mathbb Z\) in \(\C\).
Let \(t_q\) be a common covering radius for these four translates.  For every
realized \(B\), there is an admissible parity pair \((x_B,y_B)\) such that
\[
  \abs{x_B+\zeta^k y_B+B}\le t_q .
\]
If \(E\) holds, the inverse-norm bound gives
\(\abs{x_B}+\abs{y_B}\le C_q\sqrt N\), after increasing \(C_q\).
Since \(M_0,M_k\asymp_q N\), the local central binomial estimate at bounded
standard-deviation scale gives
\[
  \Prob\{S_0=x_B,\ S_k=y_B\}\ge c_q N^{-1}.
\]
Consequently, conditionally on \(E\),
\[
  \Prob\{\abs{G_m(\zeta)}\le t_q\mid (S_\ell)_{\ell\notin\{0,k\}}\}
  \ge c_q N^{-1}.
\]
Multiplying by \(\Prob(E)\ge c_q\) proves the claim.
\end{proof}

\begin{remark}[Fixed rational angles versus the continuum]\label{rem:root-unity-lattice}
Proposition~\ref{prop:rademacher-root-unity} settles the simplest discrete
analogue beyond the two parity endpoints: a root of unity of order \(q\) has
small-ball probability \(O(q/N)\), and fixed non-real rational angles therefore
have the expected two-dimensional \(N^{-1}\) fixed-point anti-concentration;
Proposition~\ref{prop:rademacher-root-fixed-radius-lower} shows that this
scale is sharp at every fixed order at a constant lattice radius.
Corollary~\ref{cor:rademacher-root-grid} packages this into an all-points bound
on a fixed rational grid, and Corollary~\ref{cor:rademacher-denom-grid} extends
the same elementary method to all denominators \(q\le Q\) at the cost \(Q^3/N\).
Proposition~\ref{prop:rademacher-small-denom-lower} identifies the exact zero
atom at every fixed primitive root: on the admissible parity residue classes it
has a local limit \(\gamma_qN^{-\varphi(q)/2}\), with \(\gamma_q\) given by a
kernel covolume, while on the inadmissible classes it is absent.
Corollary~\ref{cor:rademacher-finite-layer-zero} sums these local limits over
any fixed finite set of cyclotomic orders and shows that only the first
nonempty totient layer contributes to the leading exact-zero asymptotic.
Corollary~\ref{cor:rademacher-denom-zero} promotes this to a bounded-denominator
exact-zero grid statement: for fixed \(Q\), non-real roots of denominator at most
\(Q\) have total exact-zero probability \(O_Q(N^{-1})\), this is sharp when the
cubic roots are present, and after removing the cyclotomic degree-two roots the
bound improves to \(O_Q(N^{-2})\).  Corollary~\ref{cor:rademacher-denom-zero-layer}
gives the general peeled-layer formula after removing all primitive orders below
any fixed cyclotomic degree threshold.  Corollary~\ref{cor:rademacher-denom-zero-leading}
then identifies the exact leading \(N^{-1}\) coefficient: it is the sum of the
primitive quadratic pair constants that are present on the relevant divisibility
classes.  Corollaries~\ref{cor:rademacher-quadratic-overlaps} and
\ref{cor:rademacher-quadratic-second-order} compute the degree-two pairwise
overlaps and the second-order corrections to the degree-two single atoms, and
Corollary~\ref{cor:rademacher-denom-zero-second-order} combines these with the
degree-four layer to give the full unpeeled fixed-\(Q\) expansion through
\(N^{-2}\).  Corollaries~\ref{cor:rademacher-degree-four-atoms} and
\ref{cor:rademacher-denom-zero-degree-four-leading} separately identify the
first nonquadratic peeled layer, where primitive orders \(5,8,10,12\) give the
complete order-\(N^{-2}\) coefficient after the degree-two roots are removed.
Corollary~\ref{cor:rademacher-denom-zero-with-endpoints} then adds back
\(z=\pm1\): through order \(N^{-1}\), the full bounded-denominator
root-of-unity grid is the exact endpoint-union central-binomial term plus the
same non-real coefficient \(c_Q(N)/N\).
Thus the two-dimensional \(N^{-1}\) scale is a possible genuine atom floor
exactly at the primitive quadratic roots \(q=3,4,6\), every fixed
bounded-denominator cyclotomic layer has a leading coefficient, and the
unpeeled non-real bounded-denominator exact-zero union is explicit through
order \(N^{-2}\), with the endpoint-inclusive grid explicit through
order \(N^{-1}\).
Fixed-radius small balls therefore cannot, in general, be read off from the zero
atom alone.  This is still far from a Rademacher analogue of
Theorem~\ref{thm:uniform}.  The continuum problem requires uniform control over
moving arcs, large denominators, and irrational angles, where the relevant input
is a density-level local central limit theorem for the boundary field, not just
the binomial atom bound used above.
\end{remark}

\begin{remark}[Boundary lattice obstruction]\label{rem:bernoulli-lattice}
For Rademacher coefficients and the balanced convention \(m=2r\), both \(G_m(1)\) and \(G_m(-1)\) are odd sums of signs, hence neither can vanish. Thus the two parity points do \emph{not} give an atom at \(t=0\) in the balanced model considered in Theorem~\ref{thm:main}. What remains true, and is the point relevant here, is that the boundary values are lattice-valued rather than density-valued: Proposition~\ref{prop:rademacher-endpoint} gives the correct endpoint scale and shows that small-ball estimates for \(\abs{G_m(\pm1)}\) cannot behave like bounded-density estimates below the lattice spacing, while Proposition~\ref{prop:rademacher-root-unity} gives the corresponding two-dimensional fixed-point scale at each fixed-order non-real root of unity, with an explicit \(q\)-loss when the order is \(q\). In adjacent parity conventions, where the number of signs is even, exact boundary zeros at \(z=\pm1\) occur with probability of order \(m^{-1/2}\). This is why any discrete analogue of Theorem~\ref{thm:uniform} must be formulated only above the appropriate lattice scale.
\end{remark}

\subsection{The shared moral with phase retrieval}
The same continuity warning appears in \cite{PhaseRetrieval}. There, for $\pm1$ measurement matrices $A\in\{\pm1\}^{(2m-1)\times m}$, a pigeonhole argument in the first two coordinates shows that at least $m$ of the $2m-1$ rows lie in one of the two hyperplanes $\{x_1=x_2\}$, $\{x_1=-x_2\}$, so some $m\times m$ submatrix is singular and the least-singular-value functional $\omega(A)$ vanishes \emph{deterministically}, for every sign matrix \cite[Rem.~6.2]{PhaseRetrieval}. The conclusion drawn there is that continuity of the entry law is necessary for any lower bound on $\omega$. The present paper exhibits the Toeplitz mirror of this warning at the boundary-evaluation level: lattice-valued boundary fields cannot be treated by bounded-density small-ball estimates all the way down to zero. We regard this as the unifying moral of the two companions: the parity coin of \cite{Parity}, the Bernoulli pigeonhole of \cite{PhaseRetrieval}, and the boundary lattice obstruction of Remark~\ref{rem:bernoulli-lattice} are three faces of a single boundary-evaluation obstruction.

\section{Open problems}\label{sec:open}

We close with the problems left open, in decreasing order of how much of the present program they would complete.

\begin{problem}[Sharp constants and the lower tail]\label{prob:sharp-boundary}
Determine the matching lower tail and limiting law behind
Theorem~\ref{thm:uniform}.  The proof gives the correct upper scale
\(\Prob\{\min_{\T}\abs{G_m}\le t\}\le C\sqrt m\,t\) and, at the continuum
threshold \(t=\tau/\sqrt N\), the sharp whole-circle probability upper constant
\(\sqrt{\pi/3}\,\tau\), but it does not prove the matching lower bound or the
exponential minimum-modulus law expected from
\cite{YakirZeitouni,CookNguyen}.  Proposition~\ref{prop:endpoint-lower} proves
the endpoint contribution at order \(t/\sqrt m\), and
Corollaries~\ref{cor:endpoint-two-point},
\ref{cor:endpoint-continuum-negligible}, and
\ref{cor:endpoint-neighborhood-continuum} give the exact bivariate endpoint
formula, sharp second-order endpoint expansion, and continuum-scale
\(O_\tau(N^{-1})\) endpoint-neighborhood contribution before
Corollary~\ref{cor:endpoint-global-floor} makes the resulting global endpoint
floor explicit.  Proposition~\ref{prop:away-firstmoment}
proves the correct away-from-endpoint first moment for level crossings, and
Proposition~\ref{prop:bulk-kr-constant} computes the fixed-bulk down-crossing
density as \(t\sqrt N/\sqrt{12\pi}\) per unit arclength, and
Corollary~\ref{cor:bulk-continuum-intensity} converts it at
\(t=\tau/\sqrt N\) into the explicit fixed-bulk limit
\((2\pi-4\delta)\tau/\sqrt{12\pi}\) for down-crossings, while
Corollary~\ref{cor:bulk-prob-upper} gives the corresponding fixed-bulk
first-moment probability upper bound with the same down-crossing constant.
Proposition~\ref{prop:bulk-secondmoment-transfer} proves that a bound on the
factorial second moment of the fixed-bulk down-crossing count gives a
Paley--Zygmund lower tail with leading small-\(\tau\) constant
\((2\pi-4\delta)/\sqrt{12\pi}\), and
Corollary~\ref{cor:bulk-firstorder-sandwich} combines this with the sharp
fixed-bulk upper bound to give the complete first-order fixed-bulk law as
\(\tau\downarrow0\).  Corollary~\ref{cor:global-prob-upper} then combines the
fixed-bulk upper bound with endpoint-neighborhood separation to give the sharp
whole-circle continuum upper constant \(\sqrt{\pi/3}\), and
Corollary~\ref{cor:global-secondmoment-transfer} shows that uniformity as
\(\delta\downarrow0\) gives the whole-circle lower-tail constant
\(\sqrt{\pi/3}\); together, Corollary~\ref{cor:global-firstorder-sandwich}
gives the whole-circle first-order law under that uniform factorial
second-moment input.  Thus
the remaining global lower-tail problem is not the one-dimensional endpoint
small ball, the two-endpoint expansion, the continuum-scale endpoint
neighborhood separation, the endpoint floor, the bulk first-moment constant, the
sharp global upper constant, the resulting subcritical and logarithmic
inverse-moment control,
or the fixed-bulk/global first-order transfer, but the anti-clustering and
decorrelation analysis that proves the required second-factorial-moment bound
for the \(O(m)\) local down-crossing opportunities uniformly through the
endpoint-cutoff removal and ultimately upgrades it to an exponential
minimum-modulus law, including the transition beyond the exact two-endpoint
expansion at the parity points \(z=\pm1\).
\end{problem}

\begin{problem}[Non-Gaussian universality for boundary separation]\label{prob:bernoulli}
Prove Hypothesis~\ref{hyp:boundary-law}, or the sharper analogue of \eqref{eq:uniform}, for bounded-density laws at full generality. Proposition~\ref{prop:nongaussian-grid} proves the fixed-point and deterministic-grid anti-concentration, and Corollary~\ref{cor:bd-uniform-speed} shows that a deterministic \(Lm^{3/2}\) global speed cutoff gives the correct uniform continuum rate with \(L\) as the speed parameter. Corollary~\ref{cor:bd-speedtail-to-boundary} reduces the stated polynomial-scale hypothesis to proving a polynomially small speed tail at the natural \(\sqrt{\log m}\) level; Corollary~\ref{cor:bd-polymoment-boundary} verifies that tail under uniformly bounded moments of all orders, and Corollary~\ref{cor:bd-subexp-boundary} verifies it under a subexponential bounded-density assumption. Corollary~\ref{cor:variance-boundary-separation} and Theorem~\ref{thm:main} show that this natural-scale tail is not needed for the finite-section conditioning conclusion under circle clearance: the tail-free polynomial separation floor \(\min_{\T}\abs{G_m}>m^{-A}\), for every \(A>1\), already supplies the one boundary-minimum threshold needed there. Thus the missing non-Gaussian step is now the full cutoff-free continuum speed/crossing law for bounded-density coefficients, or a local-limit replacement for it, not point anti-concentration and not the clearance-conditional conditioning theorem. For discrete laws such as Rademacher, the rational-angle lattice cases are now separated: Proposition~\ref{prop:rademacher-endpoint} gives the one-dimensional endpoint scale, Proposition~\ref{prop:rademacher-root-unity} gives an \(O(q/N)\) bound at a non-real root of unity of order \(q\), Proposition~\ref{prop:rademacher-root-fixed-radius-lower} gives the matching \(N^{-1}\) fixed-radius lower bound at every fixed non-real rational angle, Corollary~\ref{cor:rademacher-root-grid} gives the resulting \(q^2/N\) fixed-grid union bound, and Corollary~\ref{cor:rademacher-denom-grid} controls all denominators \(q\le Q\) with loss \(Q^3/N\). Proposition~\ref{prop:rademacher-small-denom-lower} identifies the exact atom local limit at every fixed primitive \(q\)-th root as \(\gamma_qN^{-\varphi(q)/2}\) on the admissible parity residue classes and zero on the inadmissible classes, and Corollary~\ref{cor:rademacher-denom-zero} shows that a fixed bounded-denominator non-real grid has total exact-zero probability \(O_Q(N^{-1})\), sharp once cubic roots are included, improving to \(O_Q(N^{-2})\) away from cyclotomic degree two. Corollary~\ref{cor:rademacher-denom-zero-with-endpoints} adds the endpoint zeros back through order \(N^{-1}\), as an exact central-binomial endpoint-union term plus the same non-real coefficient. Corollary~\ref{cor:rademacher-quadratic-atoms} makes those degree-two admissible classes explicit: \(q=3,6\) require \(3\mid N\), while \(q=4\) requires \(4\mid N\). The remaining discrete analogue is therefore a continuum local-CLT theorem above the relevant lattice scale, uniformly over moving arcs, large denominators, and irrational angles, not a fixed-rational-angle or exact-zero calculation. The minimum-modulus universality of Cook--Nguyen \cite{CookNguyen} establishes an analogous law for random trigonometric polynomials; transporting it to the banded boundary field $G_m|_\T$ requires the density-level delocalization input flagged in \cite[Rem.~6.3]{PhaseRetrieval}, since a Berry--Esseen estimate is insufficient at the microscopic scale and a local CLT for the boundary field is needed. The note of caution recorded in Remark~\ref{rem:grid-artifact} applies a fortiori to any numerical study of the discrete case: the Gaussian law is $\sqrt m\,t$, with excursions of angular length $\asymp t\,m^{-3/2}$, and the apparent \(m\,t^2\) behaviour seen on coarse grids is a sampling artifact, not the small-ball exponent.
Corollary~\ref{cor:rademacher-quadratic-atoms} also records the corresponding
leading constants on those classes: \(2\sqrt3/\pi\) for \(q=3,6\) and
\(4/\pi\) for \(q=4\).  Corollaries~\ref{cor:rademacher-quadratic-overlaps} and
\ref{cor:rademacher-quadratic-second-order} compute the pairwise \(N^{-2}\)
overlap constants and the second-order single-event corrections among these
three degree-two events.  More generally,
Corollary~\ref{cor:rademacher-finite-layer-zero} gives the leading exact-zero
asymptotic on the first nonempty cyclotomic degree layer for any fixed finite
set of orders.  Corollary~\ref{cor:rademacher-denom-zero-leading} specializes
this to the full leading coefficient for every fixed bounded-denominator
exact-zero grid, and Corollary~\ref{cor:rademacher-denom-zero-second-order}
upgrades it to the full unpeeled expansion through \(N^{-2}\).
Corollary~\ref{cor:rademacher-denom-zero-with-endpoints} adds the two real
endpoints and records the full bounded-denominator root-of-unity expansion
through order \(N^{-1}\).
Corollary~\ref{cor:rademacher-denom-zero-degree-four-leading} separately gives
the first nonquadratic layer after removing primitive orders \(3,4,6\).  More generally,
Corollary~\ref{cor:rademacher-denom-zero-layer} gives the next coefficient after
peeling off all orders below any fixed cyclotomic degree threshold.  These
finite-layer calculations remove the fixed-rational exact-zero problem from the
remaining discrete continuum problem.
\end{problem}

\begin{problem}[Sharp $m$-dependence of the conditioning]\label{prob:sharp-m}
Determine the correct $m$-dependence of $\smin(T_n)$ in the balanced regime on the stable event \(\nu=r\). Proposition~\ref{prop:fs} loses a factor $\rho^{r}$ with $\rho\asymp m^{-2}$, hence $\smin\gtrsim m^{-O(m)}$ at fixed large $n$; is the true typical size only polynomially small in $m$ after conditioning on stability? This connects to the conjectural disk-count fluctuation and local-limit problem isolated in \cite[Conj.~1.3]{Parity}.
\end{problem}

\begin{problem}[Joint law of boundary minimum and parity]\label{prob:joint}
Quantify the dependence between the boundary minimum $\mu=\min_\T\abs{G_m}$ and the parity $\mathbf1\{G_m(1)G_m(-1)<0\}$. Both are driven by the boundary field; \eqref{eq:parity} and Lemma~\ref{lem:onepoint} suggest that small $\mu$ is correlated with the field being small near a real point, where the parity is decided. A joint local limit theorem for $(G_m(1),G_m(-1),\mu)$ would unify the anti-concentration of \cite{Parity} with the conditioning bound here.
\end{problem}

\subsection*{Acknowledgments}
The author thanks the Randomstrasse 101 project for the circle of questions around stable phase retrieval, and acknowledges that this paper is a companion to \cite{Parity} and \cite{PhaseRetrieval}, from which it borrows, respectively, the operator-theoretic framework and the negative-second-moment small-ball machinery. The research was supported by the European Research Council (ERC) under the European Union's Horizon Europe research and innovation programme (grant agreement No.\ 101041711), by the Israel Science Foundation (grant number 2258/19) and by the Israel Science Foundation (ISF Grant 4101/25).

\subsection*{Declaration on generative AI}
During the preparation of this work the author used generative AI tools to accelerate drafting and revision. All mathematical claims, proofs, and bibliographic information were subsequently reviewed and edited by the author, who takes full responsibility for the content of the manuscript.

\begin{thebibliography}{9}

\bibitem{Parity}
Y.~Shmalo,
\emph{A parity obstruction to the stability of random banded Toeplitz matrices},
companion manuscript, 2026.

\bibitem{PhaseRetrieval}
Y.~Shmalo,
\emph{Extreme least singular values of Gaussian row submatrices and a phase retrieval stability problem},
companion manuscript, 2026.

\bibitem{BG}
A.~B\"ottcher and S.~M.~Grudsky,
\emph{Spectral Properties of Banded Toeplitz Matrices},
SIAM, Philadelphia, 2005.

\bibitem{KMS}
S.~V.~Konyagin and W.~Schlag,
\emph{Lower bounds for the absolute value of random polynomials on a neighborhood of the unit circle},
Trans.\ Amer.\ Math.\ Soc.\ \textbf{351} (1999), no.~12, 4963--4980.

\bibitem{YakirZeitouni}
O.~Yakir and O.~Zeitouni,
\emph{The minimum modulus of Gaussian trigonometric polynomials},
Israel J.\ Math.\ \textbf{245} (2021), no.~1, 543--566.
\href{https://arxiv.org/abs/2006.08943}{arXiv:2006.08943}.

\bibitem{CookNguyen}
N.~A.~Cook and H.~H.~Nguyen,
\emph{Universality of the minimum modulus for random trigonometric polynomials},
Discrete Anal.\ \textbf{2021:20}, 46 pp.
\href{https://arxiv.org/abs/2101.07203}{arXiv:2101.07203}.

\bibitem{RudelsonVershynin2015}
M.~Rudelson and R.~Vershynin,
\emph{Small ball probabilities for linear images of high-dimensional distributions},
Int.\ Math.\ Res.\ Not.\ IMRN \textbf{2015}, no.~19, 9594--9617.
\href{https://arxiv.org/abs/1402.4492}{arXiv:1402.4492}.

\bibitem{FukNagaev}
D.~H.~Fuk and S.~V.~Nagaev,
\emph{Probability inequalities for sums of independent random variables},
Theory Probab.\ Appl.\ \textbf{16} (1971), no.~4, 643--660.

\bibitem{VershyninHDP}
R.~Vershynin,
\emph{High-Dimensional Probability: An Introduction with Applications in Data Science},
Cambridge Series in Statistical and Probabilistic Mathematics, vol.~47,
Cambridge University Press, Cambridge, 2018.

\bibitem{Videnskii}
V.~S.~Videnskii,
\emph{Extremal estimates for the derivative of a trigonometric polynomial on an interval shorter than its period},
Soviet Math.\ Dokl.\ \textbf{1} (1960), 5--8.

\bibitem{NagyTotik}
B.~Nagy and V.~Totik,
\emph{Bernstein's inequality for algebraic polynomials on circular arcs},
Constr.\ Approx.\ \textbf{37} (2013), no.~2, 223--232.

\bibitem{BalanWang2015}
R.~Balan and Y.~Wang,
\emph{Invertibility and robustness of phaseless reconstruction},
Appl.\ Comput.\ Harmon.\ Anal.\ \textbf{38} (2015), no.~3, 469--488.

\bibitem{RudelsonVershynin2009}
M.~Rudelson and R.~Vershynin,
\emph{The smallest singular value of a random rectangular matrix},
Comm.\ Pure Appl.\ Math.\ \textbf{62} (2009), no.~12, 1707--1739.

\bibitem{TaoVu2010}
T.~Tao and V.~Vu,
\emph{Random matrices: the distribution of the smallest singular values},
Geom.\ Funct.\ Anal.\ \textbf{20} (2010), no.~1, 260--297.

\bibitem{Edelman1988}
A.~Edelman,
\emph{Eigenvalues and condition numbers of random matrices},
SIAM J.\ Matrix Anal.\ Appl.\ \textbf{9} (1988), no.~4, 543--560.

\bibitem{HKPV}
J.~B.~Hough, M.~Krishnapur, Y.~Peres and B.~Vir\'ag,
\emph{Zeros of Gaussian Analytic Functions and Determinantal Point Processes},
University Lecture Series, vol.~51, Amer.\ Math.\ Soc., Providence, RI, 2009.

\end{thebibliography}

\end{document}
